\chapter{非线性优化}
\label{ch:nlopt}

% ════════════════════════════════════════════════════════════════
%  P1 收尾章·全吸收范本：吸收
%  十四讲6(1027): MAP→最小二乘全推导/SD-牛顿-GN-LM 逐步/曲线拟合手写GN·Ceres·g2o 三完整例
%  Barfoot4(03ext): Newton/GN 两推导/H≈JᵀJ 来由与缺陷/line search+LM/正规方程(误差形式)
%                    /白化Cholesky/Laplace=逆Hessian协方差/批量块结构/ML+Box(1971)偏差/SWF边缘化=Schur
%  Barfoot5(04ext): M-估计(L2/Cauchy/GM)/IRLS/鲁棒=逆Wishart先验MAP/RANSAC/NEES-NIS/协方差估计
%  Handbook1(hb_01): 因子图→MAP→NLS/白化/Cholesky-vs-QR/√SAM/稀疏=fill-in/COLAMD/变量消元=分解/Bayes树·iSAM
%  右扰动为主；教学层遵循 v5.0 理论教学规范（动机→反面→历史→理论→陷阱→练习 + 符号表/速查/故障排查）。
% ════════════════════════════════════════════════════════════════

\begin{practice}[前置自测]
开始之前，先试答下列问题。若已能流畅作答，可只速读本章表格与速查卡；若卡住，正是本章要为你解决的。
\begin{enumerate}
  \item 上一章把 SLAM 化成最小二乘 $\min_{\mathbf{x}}\tfrac12\sum_i\|\mathbf{e}_i\|^2$，但 $\mathbf{e}_i$ 非线性时为何不能像线性情形那样一步解出？（答错 $\to$ \cref{sec:nlopt-why}）
  \item 高斯牛顿法用什么近似牛顿法的海森矩阵？这个近似什么时候站得住、什么时候崩？（答错 $\to$ \cref{sec:nlopt-gn}）
  \item 列文伯格–马夸尔特的阻尼项 $\lambda$ 调大调小，分别让算法更像哪种方法？什么是“增益比 $\rho$”？（答错 $\to$ \cref{sec:nlopt-lm}）
  \item 解 $\boldsymbol{\Lambda}\Delta\mathbf{x}=\mathbf{b}$ 时，Cholesky 与 QR 各有什么优劣？为什么绝不直接求 $\boldsymbol{\Lambda}^{-1}$？（答错 $\to$ \cref{sec:nlopt-solve}）
  \item BA 的信息矩阵高达上万维，为什么还能实时求解？“Schur 消元 / 边缘化 / fill-in” 各指什么？（答错 $\to$ \cref{sec:nlopt-sparse}）
  \item 一个误匹配（外点）会怎样毁掉最小二乘？“鲁棒核函数”如何补救，它和概率（逆 Wishart 先验）有什么关系？（答错 $\to$ \cref{sec:nlopt-robust}）
\end{enumerate}
\end{practice}

\section*{本章目标}
学完本章，你应当能够：
\begin{enumerate}
  \item \textbf{推导}从最大后验（MAP）到加权最小二乘的完整链条，说清高斯假设在哪一步、独立性假设在哪一步被用到；
  \item \textbf{推导}最速下降、牛顿、高斯牛顿三种增量方程，解释 GN 用 $\mathbf{J}^\top\mathbf{W}^{-1}\mathbf{J}$ 近似海森的依据（丢掉哪一项）与失效情形；
  \item \textbf{掌握}列文伯格–马夸尔特/信赖域，理解 $\lambda$ 如何在高斯牛顿与最速下降间插值、增益比 $\rho$ 如何接受/拒绝步并调 $\lambda$；
  \item \textbf{选择}解正规方程的数值方法（白化、Cholesky 与 QR、平方根信息 $\sqrt{}$SAM），并说清它们的复杂度与数值稳定性差别；
  \item \textbf{利用}稀疏性：把 BA 的信息矩阵按因子图结构做 Schur 消元/变量消元，把代价从 $O(n^3)$ 降到 $O(\text{相机}^3)$，并解释 fill-in 与 COLAMD；
  \item \textbf{在流形上优化}：用右扰动把位姿增量定义为 $\mathbf{T}\,\mathrm{Exp}(\Delta\boldsymbol{\xi})$，对接 \cref{ch:lie} 与求解器；
  \item \textbf{鲁棒化}：在\emph{前端}用 RANSAC/成对一致性最大化（PCM，最大团）剔除外点、在\emph{后端}用 M-估计/鲁棒核 $+$ IRLS 降权外点，\emph{证明}“鲁棒核 $=$ 某协方差逆 Wishart 先验下的 MAP”并点明其为 Black--Rangarajan 对偶之实例，再用渐进非凸（GNC）对付坏初值；
  \item \textbf{提取不确定度}：收敛后由拉普拉斯近似得逆海森即后验协方差（与克拉默–拉奥下界的关系），并判断估计是否一致（NEES/NIS）；
  \item \textbf{（选读）升级视角}：把后端嵌进端到端学习（双层优化/隐式微分/海森-向量积），并用半定松弛/SE-Sync 不靠初值地求全局最优、附最优性证书；
  \item \textbf{运行}手写高斯牛顿、Ceres、g2o 三套完整曲线拟合例并对比。
\end{enumerate}

\subsection*{本章知识导航}
本章主线：\emph{怎么解}上一章立起来的最小二乘——从 MAP 立目标，到迭代下降、高斯牛顿/LM，到数值分解与利用稀疏结构、流形右扰动，再到鲁棒化、协方差与一致性，最后落到三套软件实践。

\begin{center}
\small
\begin{tikzpicture}[node distance=6mm and 7mm, >=Stealth,
  blk/.style={draw, rounded corners, align=center, font=\small, inner sep=3pt, fill=main!6, draw=main!60!black},
  grp/.style={draw, rounded corners, align=center, font=\small, inner sep=3pt, fill=second!6, draw=second!60!black},
  ln/.style={->, thick, gray!60!black}]
\node[blk] (map) {MAP\\{\scriptsize $\to$最小二乘}};
\node[blk, right=of map] (iter) {迭代下降\\{\scriptsize 梯度/牛顿}};
\node[blk, right=of iter] (gn) {高斯牛顿\\{\scriptsize $\mathbf{J}^\top\mathbf{W}^{-1}\mathbf{J}$}};
\node[blk, right=of gn] (lm) {LM·信赖域\\{\scriptsize $\rho,\lambda$}};
\node[grp, below=of lm] (solve) {白化·Cholesky/QR\\{\scriptsize $\sqrt{}$SAM}};
\node[grp, left=of solve] (sparse) {稀疏·Schur·消元\\{\scriptsize COLAMD}};
\node[grp, left=of sparse] (ba) {BA·位姿图\\{\scriptsize 重投影·箭头}};
\node[grp, left=of ba] (man) {流形优化\\{\scriptsize 右扰动}};
\node[grp, left=of man] (robust) {鲁棒核·IRLS\\{\scriptsize 逆Wishart}};
\node[blk, below=of robust, fill=third!8, draw=third!60!black] (cov) {协方差\\{\scriptsize Laplace·NEES}};
\node[blk, right=of cov, fill=third!8, draw=third!60!black] (code) {实践\\{\scriptsize Ceres/g2o}};
\draw[ln] (map)--(iter); \draw[ln] (iter)--(gn); \draw[ln] (gn)--(lm);
\draw[ln] (lm)--(solve);
\draw[ln] (solve)--(sparse); \draw[ln] (sparse)--(ba); \draw[ln] (ba)--(man); \draw[ln] (man)--(robust);
\draw[ln] (robust)--(cov); \draw[ln] (cov)--(code);
\end{tikzpicture}
\end{center}

\noindent\textbf{推荐路径（骨干 / 次优 / 前沿）}：
\emph{骨干}（任何读者必跟）= MAP$\to$最小二乘$\to$迭代$\to$GN$\to$LM、数值分解 $+$ 稀疏 Schur/消元、以及 \cref{sec:nlopt-ba} 的 BA（重投影误差$+$箭头稀疏$+$Ceres/g2o BA 实践），后端、VIO 反复用；
\emph{次优}（够用即可）= \cref{sec:nlopt-posegraph} 位姿图、流形右扰动（细节已在 \cref{ch:lie}）、协方差与一致性、软件实践；
\emph{前沿/选读}（点到为止）= \cref{sec:nlopt-robust} 的前端 PCM/最大团、Black--Rangarajan 对偶与 GNC、“鲁棒 $=$ 逆 Wishart 先验 MAP”的完整证明，\cref{sec:nlopt-cov} 的 ML 偏差闭式，\cref{sec:nlopt-gvi} 的高斯变分推断（GVI/ESGVI——把点估计升级为分布估计、接回因子图，含立体相机三方对照），\cref{sec:nlopt-diffopt} 的可微优化（双层优化/隐式微分/海森-向量积——让学习前端从几何后端反传受益），\cref{sec:nlopt-certifiable} 的可证最优（Shor 半定松弛/SE-Sync/Riemannian Staircase/最优性证书——不靠初值的全局求解），以及附录的连续时间与 SWF 边缘化。带（选读）标记的小节为进阶/选读。

\subsection*{前置知识桥接}
\cref{ch:slam_est} 把 SLAM 钉成 $\min_{\mathbf{x}}J(\mathbf{x})=\tfrac12\sum_i\|\mathbf{e}_i\|_{\boldsymbol{\Sigma}_i}^2$ 并给了\emph{线性}情形的正规方程；本章解其\emph{非线性}版。\cref{ch:eskf} 已埋两处伏笔：“EKF $\approx$ 优化的一次迭代”、“IEKF $=$ 单时刻的高斯牛顿”——本章把“优化”讲透并\emph{证明}这两条。位姿增量用 \cref{ch:lie} 的右扰动 $\mathbf{T}\,\mathrm{Exp}(\Delta\boldsymbol{\xi})$ 与右扰动雅可比。概率基元（高斯密度、贝叶斯法则、马氏距离、边缘化/条件化）见 \cref{ch:slam_est}，本章在用到处会就地重述，无需翻回。

\subsection*{如果跳过本章会怎样}
\begin{itemize}
  \item 你能写出 BA/位姿图的目标函数，却\emph{解不动}它——不知道 GN/LM、不知道如何用稀疏 Schur 把上万维问题压到可实时；
  \item Ceres/g2o/GTSAM 里的 \texttt{LinearSolver}、\texttt{RobustKernel}、\texttt{LocalParameterization}、\texttt{SPARSE\_SCHUR} 会像黑话；
  \item 真实数据里一条误匹配就能把整条轨迹拉飞，而你不知道前端 RANSAC/PCM 如何\emph{剔}外点、后端鲁棒核如何\emph{降权}外点、GNC 为何不是“黑魔法”；
  \item 你不知道 GN/LM 只保证\emph{局部}最优、初值差会陷局部极小，更不知道半定松弛/SE-Sync 如何\emph{不靠初值}求全局最优并附\emph{最优性证书}；也不知道如何让学习前端经“反传几何后端”受益（可微优化）；
  \item 你将无法理解 \cref{ch:vo}、\cref{ch:vio} 后端“为什么这样建图、这样求解、这样估协方差”。
\end{itemize}

\begin{table}[H]\centering\small
\caption{本章阅读方式建议}
\begin{tabular}{lll}
\toprule
阅读方式 & 时间 & 适合谁 \\
\midrule
精读（含全部推导、例与练习） & 6--8 小时 & 要写/调 BA、位姿图、VIO 后端 \\
速读（跳过 QR/消元/ML 偏差细节） & 2--3 小时 & 已懂 GN/LM、想抓稀疏 Schur 与右扰动 \\
速查（只看小结三张表 $+$ 算法卡） & 15 分钟 & 写代码时回来查公式 \\
\bottomrule
\end{tabular}
\end{table}

\begin{note}
\textbf{本章记号与约定.}\;
目标 $J(\mathbf{x})=\tfrac12\sum_i\|\mathbf{e}_i\|_{\boldsymbol{\Sigma}_i}^2=\tfrac12\mathbf{e}(\mathbf{x})^\top\mathbf{W}^{-1}\mathbf{e}(\mathbf{x})$，权 $\mathbf{W}=\operatorname{diag}(\boldsymbol{\Sigma}_i)$（取协方差时即马氏距离/负对数似然）。\emph{残差} $\mathbf{e}(\mathbf{x})$、其\emph{雅可比} $\mathbf{J}=\partial\mathbf{e}/\partial\mathbf{x}$；目标的梯度记 $\nabla J$、海森记 $\mathbf{H}_J$（以区别于残差雅可比 $\mathbf{J}$）。增量 $\Delta\mathbf{x}$（或李代数语境记 $\delta\mathbf{x}$）；信息矩阵 $\boldsymbol{\Lambda}=\mathbf{J}^\top\mathbf{W}^{-1}\mathbf{J}$。鲁棒核 $\rho$。先验 $\check{\mathbf{x}},\check{\mathbf{P}}$；后验 $\hat{\mathbf{x}},\hat{\mathbf{P}}$；线性化工作点 $\mathbf{x}_{\mathrm{op}}$。
\textbf{流形约定}：当 $\mathbf{x}$ 含位姿时，更新用\emph{右扰动} $\mathbf{x}\leftarrow\mathbf{x}\boxplus\Delta\mathbf{x}$，位姿部分即 $\mathbf{T}\leftarrow\mathbf{T}\,\mathrm{Exp}(\Delta\boldsymbol{\xi})$（\cref{ch:lie}）。

\textbf{跨书记号差异（务必先看，三书在此打架）}：
（1）十四讲\cite{gaoxiang2019slam14} §6.2 把残差雅可比 $\mathbf{J}$ 定义为\emph{列向量}（标量残差的梯度），故写海森近似 $\mathbf{H}\approx\mathbf{J}\mathbf{J}^\top$、正规方程 $\mathbf{J}\mathbf{J}^\top\Delta\mathbf{x}=-\mathbf{J}f$；本书与多数文献（Ceres/g2o）把 $\mathbf{J}=\partial\mathbf{e}/\partial\mathbf{x}$ 当\emph{矩阵}（行$=$残差维、列$=$状态维），写 $\mathbf{J}^\top\mathbf{W}^{-1}\mathbf{J}$。两者“矩阵相同、符号互为转置”，本书统一用矩阵约定。
（2）十四讲把运动噪声协方差记 $R_k$、观测噪声协方差记 $Q_{k,j}$（与卡尔曼惯例\emph{相反}！）；Barfoot\cite{barfoot2024state} 用 $\mathbf{Q}$（过程）/$\mathbf{R}$（观测）。本书统一 $\boldsymbol{\Sigma}_w$（过程）/$\boldsymbol{\Sigma}_v$（观测），把字母 $\mathbf{R}$ 留给旋转。
（3）Barfoot 记 $\mathbf{H}=-\partial\mathbf{e}/\partial\mathbf{x}$（含负号），与本书 $\mathbf{J}=\partial\mathbf{e}/\partial\mathbf{x}$ 差一负号（不影响正规方程，因左右各出现偶数次）；Handbook\cite{carlone2026handbook} 把白化雅可比记 $\mathbf{A}$、Cholesky 上三角因子也记 $\mathbf{R}$（与旋转撞名，本书该处改记 $\mathbf{U}$ 并显式说明）。
\end{note}

% ════════════════════════════════════════════════════════════════
\section{从最大后验到最小二乘}
\label{sec:nlopt-map}

\paragraph{动机：噪声让方程不再精确成立。}
经典 SLAM 模型由运动方程与观测方程组成\cite{gaoxiang2019slam14}：
\begin{equation}
  \begin{cases}
    \mathbf{x}_k=f(\mathbf{x}_{k-1},\mathbf{u}_k)+\mathbf{w}_k, & \mathbf{w}_k\sim\mathcal{N}(\mathbf{0},\boldsymbol{\Sigma}_{w,k}),\\
    \mathbf{z}_{k,j}=h(\mathbf{y}_j,\mathbf{x}_k)+\mathbf{v}_{k,j}, & \mathbf{v}_{k,j}\sim\mathcal{N}(\mathbf{0},\boldsymbol{\Sigma}_{v,k,j}),
  \end{cases}
  \label{eq:nlopt-slam-model}
\end{equation}
其中 $\mathbf{x}_k$ 是相机位姿（可用 $\mathbf{T}_k\in\SE(3)$ 表达），$\mathbf{y}_j$ 是路标，$\mathbf{z}_{k,j}$ 是像素观测，$\mathbf{u}_k$ 是输入。观测方程对针孔相机即 $s\,\mathbf{z}_{k,j}=\mathbf{K}(\mathbf{R}_k\mathbf{y}_j+\mathbf{t}_k)$（$\mathbf{K}$ 为内参，$s$ 为深度，即 $\mathbf{R}_k\mathbf{y}_j+\mathbf{t}_k$ 的第三分量，\cref{ch:camera}）。因噪声 $\mathbf{w},\mathbf{v}$ 存在，等式\emph{必不精确成立}——即便高精度相机也只是近似。于是与其“假设数据必须符合方程”，不如问：\emph{如何在带噪数据里做准确的状态估计}\cite{gaoxiang2019slam14}。

\paragraph{反面：不带概率地“硬解”会怎样。}
若直接把 \cref{eq:nlopt-slam-model} 当作等式联立求解，方程数（观测数）通常远多于未知数，系统超定且互相矛盾（每个噪声把对应方程顶歪一点），\emph{无解}。强行取某个“折中”又缺乏依据——凭什么相信这条观测多于那条？答案是\emph{按不确定度加权}：观测越准（协方差越小、信息越大），它说话越算数。这正是下面 MAP$\to$最小二乘要给出的、有概率依据的“加权折中”。

\paragraph{历史：从 Gauss 的最小二乘到 SLAM 的 MAP。}
最小二乘可追到 Gauss（1795，行星轨道拟合）与 Legendre（1805）。把它放进贝叶斯框架（先验$\times$似然$=$后验、最大化后验）是现代状态估计的主线；Barfoot\cite{barfoot2024state} 与 Handbook\cite{carlone2026handbook,dellaert2017factor} 都以 MAP 为起点。SLAM 后端今天普遍采用\emph{批量}最大后验/最大似然，而非历史上的滤波\cite{gaoxiang2019slam14}。

\paragraph{两种处理方式（三书共识）。}\cite{gaoxiang2019slam14,barfoot2024state}
\begin{itemize}
  \item \textbf{增量/滤波（incremental）}：数据随时间到来，持当前估计、用新数据更新（EKF 及衍生）。只关心当前时刻 $\mathbf{x}_k$。
  \item \textbf{批量（batch）}：把 $0\sim k$ 全部输入与观测“攒”起来一并处理，估计整条轨迹与地图。在更大范围内最优化，被认为优于传统滤波\cite{gaoxiang2019slam14}，是当前主流；极端是离线 SfM。
  \item \textbf{折中}：固定历史、只对当前邻近的若干轨迹优化——\emph{滑动窗口估计}（\cref{sec:nlopt-appendix} 的 SWF）。
\end{itemize}
本章以批量非线性优化为主。

\paragraph{核心推导：MAP$\to$加权最小二乘（逐步，无跳步）。}
设全部时刻状态 $\mathbf{x}=\{\mathbf{x}_1,\dots,\mathbf{x}_N\}$、路标 $\mathbf{y}=\{\mathbf{y}_1,\dots,\mathbf{y}_M\}$。状态估计 $=$ 求后验 $P(\mathbf{x},\mathbf{y}\mid\mathbf{z},\mathbf{u})$\cite{gaoxiang2019slam14}。直接求整个分布难，但求其\emph{峰值}（最大后验，MAP）可行。用贝叶斯法则（\cref{thm:nlopt-map}）。

\begin{theorem}[MAP $\Rightarrow$ 加权最小二乘]\label{thm:nlopt-map}
高斯先验 $+$ 高斯（加性）观测噪声 $+$ 各因子条件独立时，最大后验估计等价于求解
\begin{equation}
  \min_{\mathbf{x},\mathbf{y}}\ J=\sum_k\mathbf{e}_{u,k}^\top\boldsymbol{\Sigma}_{w,k}^{-1}\mathbf{e}_{u,k}
  +\sum_{k,j}\mathbf{e}_{z,k,j}^\top\boldsymbol{\Sigma}_{v,k,j}^{-1}\mathbf{e}_{z,k,j},
  \label{eq:nlopt-batch-ls}
\end{equation}
即\emph{各误差项的马氏距离（信息矩阵加权二次型）之和}最小，其中 $\mathbf{e}_{u,k}=\mathbf{x}_k-f(\mathbf{x}_{k-1},\mathbf{u}_k)$，$\mathbf{e}_{z,k,j}=\mathbf{z}_{k,j}-h(\mathbf{x}_k,\mathbf{y}_j)$。
\end{theorem}
\begin{proof}
\emph{Step 1（贝叶斯法则，丢分母）}。
\begin{equation}
  P(\mathbf{x},\mathbf{y}\mid\mathbf{z},\mathbf{u})=\frac{P(\mathbf{z},\mathbf{u}\mid\mathbf{x},\mathbf{y})\,P(\mathbf{x},\mathbf{y})}{P(\mathbf{z},\mathbf{u})}\propto\underbrace{P(\mathbf{z},\mathbf{u}\mid\mathbf{x},\mathbf{y})}_{\text{似然}}\underbrace{P(\mathbf{x},\mathbf{y})}_{\text{先验}}.
  \label{eq:nlopt-bayes}
\end{equation}
分母 $P(\mathbf{z},\mathbf{u})$ 与待估状态无关，求 $\arg\max$ 时可丢。故 $\arg\max$ 后验 $=\arg\max(\text{似然}\times\text{先验})$。无先验信息（取均匀先验）时退化为最大似然 $\arg\max P(\mathbf{z},\mathbf{u}\mid\mathbf{x},\mathbf{y})$。Handbook\cite{carlone2026handbook} 与本式同形（其 $\mathbf{x}^{\text{MAP}}=\arg\max p(\mathbf{z}\mid\mathbf{x})p(\mathbf{x})$）。

\emph{Step 2（取负对数，把“最大化乘积”变“最小化求和”）}。$N$ 维高斯密度
\begin{equation}
  P(\mathbf{a})=\frac{1}{\sqrt{(2\pi)^N\det\boldsymbol{\Sigma}}}\exp\!\Big(-\tfrac12(\mathbf{a}-\boldsymbol{\mu})^\top\boldsymbol{\Sigma}^{-1}(\mathbf{a}-\boldsymbol{\mu})\Big),
  \label{eq:nlopt-gauss}
\end{equation}
取负对数
\begin{equation}
  -\ln P(\mathbf{a})=\tfrac12\ln\big((2\pi)^N\det\boldsymbol{\Sigma}\big)+\tfrac12(\mathbf{a}-\boldsymbol{\mu})^\top\boldsymbol{\Sigma}^{-1}(\mathbf{a}-\boldsymbol{\mu}).
  \label{eq:nlopt-neglog}
\end{equation}
对数单调递增，故\emph{最大化 $P$} $=$ \emph{最小化 $-\ln P$}。第一项 $\tfrac12\ln((2\pi)^N\det\boldsymbol{\Sigma})$ 与待估量无关、可略，\textbf{只剩右侧二次型}。

\emph{Step 3（单次观测的二次型）}。观测 $\mathbf{z}_{k,j}=h(\mathbf{y}_j,\mathbf{x}_k)+\mathbf{v}_{k,j}$、$\mathbf{v}\sim\mathcal{N}(\mathbf{0},\boldsymbol{\Sigma}_{v,k,j})$ 给出 $P(\mathbf{z}_{k,j}\mid\mathbf{x}_k,\mathbf{y}_j)=\mathcal{N}(h(\mathbf{y}_j,\mathbf{x}_k),\boldsymbol{\Sigma}_{v,k,j})$，于是单次最大似然
\begin{equation}
  (\mathbf{x}_k,\mathbf{y}_j)^*=\arg\min\ \big(\mathbf{z}_{k,j}-h(\mathbf{x}_k,\mathbf{y}_j)\big)^\top\boldsymbol{\Sigma}_{v,k,j}^{-1}\big(\mathbf{z}_{k,j}-h(\mathbf{x}_k,\mathbf{y}_j)\big),
  \label{eq:nlopt-mahal}
\end{equation}
即\emph{马哈拉诺比斯距离}（马氏距离），也可看作由\emph{信息矩阵} $\boldsymbol{\Sigma}^{-1}$ 加权的欧氏距离\cite{gaoxiang2019slam14}。

\emph{Step 4（独立性因式分解）}。设各时刻输入、观测相互独立（且输入与观测独立）：
\begin{equation}
  P(\mathbf{z},\mathbf{u}\mid\mathbf{x},\mathbf{y})=\prod_k P(\mathbf{u}_k\mid\mathbf{x}_{k-1},\mathbf{x}_k)\prod_{k,j}P(\mathbf{z}_{k,j}\mid\mathbf{x}_k,\mathbf{y}_j).
  \label{eq:nlopt-factorize}
\end{equation}

\emph{Step 5（合并）}。对 \cref{eq:nlopt-factorize} 取负对数：乘积变求和，每项按 Step 2--3 化为二次型，与状态无关的常数项丢弃，定义误差 $\mathbf{e}_{u,k},\mathbf{e}_{z,k,j}$，即得 \cref{eq:nlopt-batch-ls}。其解等价于状态的最大似然/最大后验估计。\hfill$\square$
\end{proof}

\begin{insight}
\cref{thm:nlopt-map} 的链条是全书估计的“总纲”：\textbf{高斯} $\Rightarrow$ 负对数是二次型（最小二乘）；\textbf{独立} $\Rightarrow$ 二次型可拆成一项项之和（稀疏，\cref{sec:nlopt-sparse}）；\textbf{信息矩阵加权} $\Rightarrow$ 越准的观测权重越大。换言之，最小二乘不是凑出来的损失，而是“在高斯独立假设下，找最可能产生这批观测的状态”。
\end{insight}

\paragraph{SLAM 最小二乘的三个结构特点（十四讲，逐条）。}\cite{gaoxiang2019slam14}
（1）目标由许多误差项的加权二次型组成；总维数很高，但\emph{每项只与一两个变量有关}（运动误差只连 $\mathbf{x}_{k-1},\mathbf{x}_k$；观测误差只连 $\mathbf{x}_k,\mathbf{y}_j$）——这带来\emph{稀疏}结构（\cref{sec:nlopt-sparse}）。
（2）用李代数表示增量则问题\emph{无约束}；用 $\mathbf{R},\mathbf{T}$ 直接当变量则带约束 $\mathbf{R}^\top\mathbf{R}=\mathbf{I},\det\mathbf{R}=1$，优化困难——这正是 \cref{ch:lie} 右扰动的价值（\cref{sec:nlopt-manifold}）。
（3）二次型度量误差时，误差分布决定权重；但这种二次度量\emph{对外点极敏感}（\cref{sec:nlopt-robust}）——这为后文的鲁棒核埋下伏笔（\cref{sec:nlopt-robust}）。

\paragraph{例：线性系统的批量 MAP 有闭式（十四讲 §6.1.3，全录）。}\cite{gaoxiang2019slam14}
沿 $x$ 轴运动的小车，运动 $\mathbf{x}_k=\mathbf{x}_{k-1}+\mathbf{u}_k+\mathbf{w}_k$（$\mathbf{w}_k\sim\mathcal{N}(0,\boldsymbol{\Sigma}_{w,k})$）、观测 $\mathbf{z}_k=\mathbf{x}_k+\mathbf{n}_k$（$\mathbf{n}_k\sim\mathcal{N}(0,\boldsymbol{\Sigma}_{v,k})$），$k=1,2,3$，初值 $\mathbf{x}_0$ 已知。批量变量 $\mathbf{x}=[\mathbf{x}_0,\mathbf{x}_1,\mathbf{x}_2,\mathbf{x}_3]^\top$。误差 $\mathbf{e}_{u,k}=\mathbf{x}_k-\mathbf{x}_{k-1}-\mathbf{u}_k$、$\mathbf{e}_{z,k}=\mathbf{z}_k-\mathbf{x}_k$，目标 $\min\sum_k\mathbf{e}_{u,k}^\top\boldsymbol{\Sigma}_{w,k}^{-1}\mathbf{e}_{u,k}+\sum_k\mathbf{e}_{z,k}^\top\boldsymbol{\Sigma}_{v,k}^{-1}\mathbf{e}_{z,k}$。因系统\emph{线性}，可写成 $\mathbf{y}-\mathbf{H}\mathbf{x}=\mathbf{e}\sim\mathcal{N}(\mathbf{0},\boldsymbol{\Sigma})$，其中（上块为运动、下块为观测）
\begin{equation}
  \mathbf{H}=\begin{bmatrix}
  1&-1&0&0\\ 0&1&-1&0\\ 0&0&1&-1\\
  0&1&0&0\\ 0&0&1&0\\ 0&0&0&1
  \end{bmatrix},\qquad
  \boldsymbol{\Sigma}=\operatorname{diag}(\boldsymbol{\Sigma}_{w,1},\boldsymbol{\Sigma}_{w,2},\boldsymbol{\Sigma}_{w,3},\boldsymbol{\Sigma}_{v,1},\boldsymbol{\Sigma}_{v,2},\boldsymbol{\Sigma}_{v,3}).
  \label{eq:nlopt-linear-H}
\end{equation}
于是 $\mathbf{x}^*_{\text{map}}=\arg\min\mathbf{e}^\top\boldsymbol{\Sigma}^{-1}\mathbf{e}$，唯一解即\emph{加权（广义）最小二乘}正规方程的解：
\begin{equation}
  \mathbf{x}^*_{\text{map}}=(\mathbf{H}^\top\boldsymbol{\Sigma}^{-1}\mathbf{H})^{-1}\mathbf{H}^\top\boldsymbol{\Sigma}^{-1}\mathbf{y}.
  \label{eq:nlopt-gls}
\end{equation}
（推导：$\min\|\boldsymbol{\Sigma}^{-1/2}(\mathbf{y}-\mathbf{H}\mathbf{x})\|^2$ 对 $\mathbf{x}$ 求导置零，$\mathbf{H}^\top\boldsymbol{\Sigma}^{-1}(\mathbf{H}\mathbf{x}-\mathbf{y})=\mathbf{0}$，即得。无权版即十四讲习题 1 的 $\mathbf{x}=(\mathbf{A}^\top\mathbf{A})^{-1}\mathbf{A}^\top\mathbf{b}$。）\textbf{线性高斯时一步出解，无需迭代}——这正是下一节要打破的“非线性”假设。

\begin{pitfall}[概念误区：把 MAP 的“峰值”当成后验的“均值”]
MAP 给的是后验\emph{众数（mode/峰值）}，不是\emph{均值（mean）}。线性高斯时后验恰为高斯、众数$=$均值，两者重合；\emph{一旦观测模型非线性，后验不再对称，众数 $\neq$ 均值}\cite{barfoot2024state}。\cref{ch:eskf} 的立体相机贯穿例（深度 $y=fb/x+n$，\cref{sec:ekf}）正是这一点的最小例证：即便先验与噪声都高斯，后验也被视差–深度的非线性\emph{掰歪}成不对称，MAP 落在被拉偏的众数上、系统性地偏离均值约 $-33$\,cm。本章一切 GN/LM 求的都是\emph{众数}（MAP）；要均值得用采样/sigma 点（\cref{ch:eskf}），或本章末 \cref{sec:nlopt-gvi} 的高斯变分推断（GVI），后者把同一立体相机例的偏差压到约 $0.28$\,cm。这也是“MAP/最大似然\emph{有偏}”的根（\cref{sec:nlopt-cov} 的 ML 偏差）。
\end{pitfall}

\begin{pitfall}[概念误区：以为“信息矩阵加权”可有可无]
有人把 \cref{eq:nlopt-mahal} 里的 $\boldsymbol{\Sigma}^{-1}$ 当装饰、直接用裸欧氏 $\|\mathbf{z}-h\|^2$。后果：量纲不同的观测（像素 vs 米 vs 弧度）被错误地等权相加，单位大的项主导一切；各向异性噪声（如纵向比横向准）的方向信息全丢。正确做法是\emph{白化}（\cref{sec:nlopt-solve}）：左乘 $\boldsymbol{\Sigma}^{-1/2}$ 把马氏范数化成标准欧氏范数，量纲与各向异性都被吸收。
\end{pitfall}

\begin{practice}
\begin{enumerate}
  \item （推导，草稿纸）完整重做 \cref{thm:nlopt-map} 的 Step 2--5，特别说明“对数单调”和“常数项可略”分别用在哪一步。答案要点：单调用于把 $\arg\max P$ 换成 $\arg\min(-\ln P)$；常数项 $\tfrac12\ln((2\pi)^N\det\boldsymbol{\Sigma})$ 不含 $\mathbf{x}$ 故不影响 $\arg\min$。
  \item 证明超定线性方程 $\mathbf{A}\mathbf{x}=\mathbf{b}$ 的最小二乘解为 $\mathbf{x}=(\mathbf{A}^\top\mathbf{A})^{-1}\mathbf{A}^\top\mathbf{b}$（即 \cref{eq:nlopt-gls} 无权版）。要点：$\min\|\mathbf{A}\mathbf{x}-\mathbf{b}\|^2$ 求导得 $2\mathbf{A}^\top(\mathbf{A}\mathbf{x}-\mathbf{b})=\mathbf{0}$。
  \item （概念）若把上例 \cref{eq:nlopt-linear-H} 的小车运动改成 $\mathbf{x}_k=\cos(\mathbf{x}_{k-1})+\mathbf{u}_k+\mathbf{w}_k$，\cref{eq:nlopt-gls} 还成立吗？为什么？（答错 $\to$ \cref{sec:nlopt-why}）
\end{enumerate}
\end{practice}

% ════════════════════════════════════════════════════════════════
\section{为什么要迭代}
\label{sec:nlopt-why}

\paragraph{问题：非线性最小二乘。}
把 \cref{thm:nlopt-map} 的目标简记为
\begin{equation}
  \min_{\mathbf{x}}\ J(\mathbf{x})=\tfrac12\big\|\mathbf{e}(\mathbf{x})\big\|^2=\tfrac12\sum_i\|\mathbf{e}_i(\mathbf{x})\|_{\boldsymbol{\Sigma}_i}^2,
  \label{eq:nls}
\end{equation}
其中 $\mathbf{e}(\mathbf{x})$ 堆叠所有（经信息矩阵加权的）残差。系数 $\tfrac12$ 只为求导方便、不影响最优解。

\paragraph{为什么不能一步解出。}
解析路线是令 $\mathrm{d}J/\mathrm{d}\mathbf{x}=\mathbf{0}$ 求驻点\cite{gaoxiang2019slam14}。$\mathbf{e}$ \emph{线性}时这是线性方程（正规方程），可直接解（\cref{eq:nlopt-gls}）；$\mathbf{e}$ \emph{非线性}时（相机投影、李群位姿、$\cos/\mathrm{atan2}$ 等），该方程一般\emph{无闭式解}，且求它需要知道目标函数的\emph{全局}性质，而这通常不可能。解此方程“可能得到极大、极小或鞍点，需逐一比较函数值”——对高维问题不现实。

\paragraph{迭代下降框架（三书一致）。}\cite{gaoxiang2019slam14,carlone2026handbook,barfoot2024state}
于是改求\emph{局部}下降：从初值出发，反复找一个让目标变小的增量。
\begin{enumerate}
  \item 给定初值 $\mathbf{x}_0$；
  \item 第 $k$ 次迭代，寻找增量 $\Delta\mathbf{x}_k$ 使 $\|\mathbf{e}(\mathbf{x}_k\boxplus\Delta\mathbf{x}_k)\|^2$（即 $J(\mathbf{x}_k\boxplus\Delta\mathbf{x}_k)$）下降；
  \item 若 $\|\Delta\mathbf{x}_k\|<\varepsilon$（或 $\|\nabla J\|<\varepsilon$），停止；
  \item 否则更新 $\mathbf{x}_{k+1}=\mathbf{x}_k\boxplus\Delta\mathbf{x}_k$，返回 2。
\end{enumerate}
关键在于：每步只需 $\mathbf{e}$ 在当前点的\emph{局部}性质（可线性化），增量的计算就简单很多——“求导函数为零”被换成“不断寻找下降增量”。各算法的差别只在\emph{如何局部近似目标}与\emph{如何据近似求增量}\cite{carlone2026handbook}。更新第 4 步当 $\mathbf{x}$ 含位姿时用右扰动 $\boxplus$（\cref{sec:nlopt-manifold}），普通向量则就是加法 $\mathbf{x}+\Delta\mathbf{x}$。

\begin{pitfall}[初值决定成败]
非线性优化\emph{几乎所有}迭代法都依赖一个好初值\cite{gaoxiang2019slam14,carlone2026handbook}：目标非凸，初值差极易陷局部极小，甚至发散。Handbook\cite{carlone2026handbook} 提醒：\cref{eq:nls} 一般非凸，局部方法不保证全局最优——这催生了\emph{可证最优（certifiably optimal）}求解器（如位姿图的 SE-Sync，留作前沿）。SLAM 中初值常来自 ICP、PnP、对极几何（\cref{ch:vo}）等几何方法。“写出目标函数”只是一半，“给好初值”是另一半。
\end{pitfall}

\begin{practice}
\begin{enumerate}
  \item 把 \cref{eq:nls} 的“解析法 vs 迭代法”用一句话各自概括，并说明为什么相机重投影残差属于后者。
  \item （概念）举一个一维非凸 $J(x)$，使从两个不同初值出发的梯度下降收敛到不同极小。
\end{enumerate}
\end{practice}

% ════════════════════════════════════════════════════════════════
\section{下降方向：梯度法与牛顿法}
\label{sec:nlopt-gradient}

把\emph{目标} $J$ 在 $\mathbf{x}_k$ 处二阶泰勒展开（注意是 $J$ 不是残差 $\mathbf{e}$）\cite{gaoxiang2019slam14,barfoot2024state}：
\begin{equation}
  J(\mathbf{x}_k+\Delta\mathbf{x})\approx J(\mathbf{x}_k)+\nabla J^\top\Delta\mathbf{x}+\tfrac12\Delta\mathbf{x}^\top\mathbf{H}_J\Delta\mathbf{x},
  \label{eq:taylor-J}
\end{equation}
$\nabla J$ 为梯度、$\mathbf{H}_J$ 为海森（对称，二次近似有良定义极小要求其正定）。

\paragraph{最速下降（一阶梯度法）。}
只留一阶项 $\nabla J^\top\Delta\mathbf{x}$，要它最负，取反梯度方向
\begin{equation}
  \Delta\mathbf{x}^*=-\nabla J\quad(\text{配步长 }\alpha,\ \Delta\mathbf{x}=-\alpha\nabla J).
  \label{eq:nlopt-sd}
\end{equation}
步长 $\alpha$ 可按线搜索或经验定\cite{gaoxiang2019slam14}。一阶近似下目标必降，但\emph{过于贪心、走锯齿路线}，近极小处收敛慢。对 \cref{eq:nls} 的最小二乘目标，把它局部近似为 $J\approx\|\mathbf{A}(\mathbf{x}-\mathbf{x}_k)-\mathbf{b}\|^2$（$\mathbf{A},\mathbf{b}$ 为白化雅可比与残差，\cref{sec:nlopt-solve}），则在 $\mathbf{x}_k$ 处\emph{精确梯度} $\nabla J=-2\mathbf{A}^\top\mathbf{b}$\cite{carlone2026handbook}。

\paragraph{牛顿法（二阶梯度法）。}
留二阶项、对 $\Delta\mathbf{x}$ 求导置零。对 \cref{eq:taylor-J} 求导：零次项导数为 $\mathbf{0}$，一次项 $\nabla J^\top\Delta\mathbf{x}$ 导数为 $\nabla J$，二次项 $\tfrac12\Delta\mathbf{x}^\top\mathbf{H}_J\Delta\mathbf{x}$ 导数为 $\mathbf{H}_J\Delta\mathbf{x}$（$\mathbf{H}_J$ 对称），令和为零得\emph{牛顿增量方程}
\begin{equation}
  \mathbf{H}_J\,\Delta\mathbf{x}^*=-\nabla J.
  \label{eq:nlopt-newton}
\end{equation}
Barfoot\cite{barfoot2024state} 把这写成 $\big(\partial^2 J/\partial\mathbf{x}\partial\mathbf{x}^\top\big|_{\mathbf{x}_{\mathrm{op}}}\big)\delta\mathbf{x}^*=-\big(\partial J/\partial\mathbf{x}\big|_{\mathbf{x}_{\mathrm{op}}}\big)^\top$，并给三点评注：(i) \emph{局部}收敛（初值够近才保证，复杂非线性目标这已是最好期望）；(ii) 收敛率\emph{二次}（远快于梯度下降）；(iii) \emph{难实现}——须算海森 $\mathbf{H}_J$，“问题规模大时非常困难，通常避免”\cite{gaoxiang2019slam14,barfoot2024state}。

\paragraph{两法的通病与出路。}\cite{gaoxiang2019slam14}
最速下降\emph{贪心、锯齿、慢}；牛顿\emph{要海森、贵}。一般问题可用拟牛顿（用历史梯度近似海森）；而\emph{最小二乘问题}有更省的专用法——高斯牛顿与 LM，它们\emph{不需显式海森}（\cref{sec:nlopt-gn,sec:nlopt-lm}）。

\begin{pitfall}[思维陷阱：把“梯度法/牛顿法/GN”当成线性递进的“越新越好”]
三者不是“改进链”，而是\emph{不同的局部近似策略}：最速下降用\emph{一阶}近似（只信梯度方向）、牛顿用\emph{完整二阶}（信精确曲率，但要海森）、高斯牛顿用\emph{最小二乘专属的近似二阶}（用 $\mathbf{J}^\top\mathbf{W}^{-1}\mathbf{J}$ 替海森）。选哪个要问：能不能算海森？残差是否小？问题是否病态？（强非线性$+$大残差时 GN 的近似海森反而差，见 \cref{sec:nlopt-gn}。）
\end{pitfall}

\begin{practice}
\begin{enumerate}
  \item 由 \cref{eq:taylor-J} 完整推出 \cref{eq:nlopt-newton}，写清每一项求导。
  \item （概念）为什么牛顿法在二次函数上\emph{一步}到极小，而最速下降不能？
  \item 验证最小二乘局部近似下 $\nabla J=-2\mathbf{A}^\top\mathbf{b}$，并说明它与 \cref{eq:gn-normal} 右端的关系。
\end{enumerate}
\end{practice}

% ════════════════════════════════════════════════════════════════
\section{高斯牛顿法}
\label{sec:nlopt-gn}

\paragraph{核心思想：线性化\emph{残差}，而非目标。}
高斯牛顿（GN）把\emph{残差} $\mathbf{e}(\mathbf{x})$ 一阶展开（这是与牛顿法的关键区别——牛顿展开的是目标 $J$）\cite{gaoxiang2019slam14,barfoot2024state}：
\begin{equation}
  \mathbf{e}(\mathbf{x}_k\boxplus\Delta\mathbf{x})\approx\mathbf{e}(\mathbf{x}_k)+\mathbf{J}\Delta\mathbf{x},\qquad \mathbf{J}=\frac{\partial\mathbf{e}}{\partial\mathbf{x}}\bigg|_{\mathbf{x}_k}.
  \label{eq:gn-linearize}
\end{equation}
代回 \cref{eq:nls} 得关于 $\Delta\mathbf{x}$ 的\emph{线性}最小二乘 $\min_{\Delta\mathbf{x}}\tfrac12(\mathbf{e}+\mathbf{J}\Delta\mathbf{x})^\top\mathbf{W}^{-1}(\mathbf{e}+\mathbf{J}\Delta\mathbf{x})$。

\begin{theorem}[高斯牛顿增量方程]\label{thm:gn}
\cref{eq:gn-linearize} 的最优增量满足
\begin{equation}
  \boxed{\ \underbrace{(\mathbf{J}^\top\mathbf{W}^{-1}\mathbf{J})}_{\textstyle\boldsymbol{\Lambda}\ \approx\,\text{海森}}\,\Delta\mathbf{x}^*
  =\underbrace{-\mathbf{J}^\top\mathbf{W}^{-1}\mathbf{e}(\mathbf{x}_k)}_{\textstyle\text{梯度的相反数}}\ }
  \label{eq:gn-normal}
\end{equation}
即“增量方程 / 高斯牛顿方程 / 正规方程”。GN 用 $\boldsymbol{\Lambda}=\mathbf{J}^\top\mathbf{W}^{-1}\mathbf{J}$ \emph{近似}牛顿法的海森 $\mathbf{H}_J$，从而\textbf{无需任何二阶导数}。
\end{theorem}
\begin{proof}
\emph{途径一（展开二次型）}\cite{gaoxiang2019slam14}。把 \cref{eq:gn-linearize} 代入 $\tfrac12\|\mathbf{e}+\mathbf{J}\Delta\mathbf{x}\|_{\mathbf{W}}^2$ 并展开：
\[
  \tfrac12\big(\mathbf{e}+\mathbf{J}\Delta\mathbf{x}\big)^\top\mathbf{W}^{-1}\big(\mathbf{e}+\mathbf{J}\Delta\mathbf{x}\big)
  =\tfrac12\Big(\mathbf{e}^\top\mathbf{W}^{-1}\mathbf{e}+2\,\mathbf{e}^\top\mathbf{W}^{-1}\mathbf{J}\Delta\mathbf{x}+\Delta\mathbf{x}^\top\mathbf{J}^\top\mathbf{W}^{-1}\mathbf{J}\Delta\mathbf{x}\Big).
\]
对 $\Delta\mathbf{x}$ 求导置零：$\mathbf{J}^\top\mathbf{W}^{-1}\mathbf{e}+\mathbf{J}^\top\mathbf{W}^{-1}\mathbf{J}\Delta\mathbf{x}=\mathbf{0}$，即 \cref{eq:gn-normal}。

\emph{途径二（GN 是牛顿法的近似）}\cite{barfoot2024state}。目标 $J=\tfrac12\mathbf{e}^\top\mathbf{W}^{-1}\mathbf{e}$，先白化 $\mathbf{u}=\mathbf{L}\mathbf{e}$（$\mathbf{L}^\top\mathbf{L}=\mathbf{W}^{-1}$，\cref{sec:nlopt-solve}），$J=\tfrac12\mathbf{u}^\top\mathbf{u}$。其精确雅可比与海森（记 $\mathbf{U}=\partial\mathbf{u}/\partial\mathbf{x}$，$u_i$ 为 $\mathbf{u}$ 的分量）：
\[
  \frac{\partial J}{\partial\mathbf{x}}=\mathbf{u}^\top\mathbf{U},\qquad
  \mathbf{H}_J=\frac{\partial^2 J}{\partial\mathbf{x}\partial\mathbf{x}^\top}=\mathbf{U}^\top\mathbf{U}+\sum_{i}u_i\frac{\partial^2 u_i}{\partial\mathbf{x}\partial\mathbf{x}^\top}.
\]
\textbf{近似的依据}：在极小附近残差 $u_i$ 小（理想为零），含\emph{残差$\times$二阶导}的第二项相对 $\mathbf{U}^\top\mathbf{U}$ 可忽略，故 $\mathbf{H}_J\approx\mathbf{U}^\top\mathbf{U}$，\emph{不含任何二阶导}。代入牛顿步 \cref{eq:nlopt-newton} 得 $\mathbf{U}^\top\mathbf{U}\,\delta\mathbf{x}^*=-\mathbf{U}^\top\mathbf{u}$；用 $\mathbf{u}=\mathbf{L}\mathbf{e}$、$\mathbf{U}=\mathbf{L}\mathbf{J}$、$\mathbf{L}^\top\mathbf{L}=\mathbf{W}^{-1}$ 回代即 \cref{eq:gn-normal}。
\hfill$\square$
\end{proof}

\begin{insight}
\cref{eq:gn-normal} 与 \cref{ch:slam_est} 的正规方程\emph{同形}——GN 就是“在当前线性化点解一次 \cref{ch:slam_est} 的正规方程，更新，再线性化”，反复迭代。这也兑现了 \cref{ch:eskf} 的两处伏笔：\textbf{EKF $\approx$ GN 的一次迭代}（线性化只做一次、不迭代到收敛）；\textbf{IEKF $=$ 单时刻 MAP 的 GN}（每步把工作点设为上次后验均值、迭代重线性化到收敛，收敛到全后验的\emph{众数}即 MAP，证明见 \cref{thm:nlopt-iekf-map}）。滤波是优化的特例。
\end{insight}

\paragraph{“先线性化残差”是流形求导的捷径（Barfoot 的 GN shortcut）。}
\cref{eq:gn-linearize} 这条“先对残差线性化、再组装最小二乘”的路子，Barfoot\cite{barfoot2024state} 称之为 Gauss-Newton \emph{shortcut}：不必先写出目标 $J$ 的二阶导，只要把残差对增量线性化即可。它正是 \cref{sec:nlopt-manifold} 在李群上求导的母本：把 $\Delta\mathbf{x}$ 理解为右扰动 $\delta\boldsymbol{\xi}$，$\mathbf{J}$ 用右扰动雅可比（\cref{ch:lie}），其余照旧。Barfoot 也由“批量误差线性化 $\to$ 堆叠 $\to$ 解 $(\mathbf{H}^\top\mathbf{W}^{-1}\mathbf{H})\delta\mathbf{x}=\mathbf{H}^\top\mathbf{W}^{-1}\mathbf{e}$”得到与 \cref{eq:gn-normal} 同形的批量更新（他记 $\mathbf{H}=-\partial\mathbf{e}/\partial\mathbf{x}$，差一负号不影响）。

\paragraph{GN 算法步骤。}\cite{gaoxiang2019slam14}
（1）给初值 $\mathbf{x}_0$；（2）第 $k$ 次迭代算雅可比 $\mathbf{J}(\mathbf{x}_k)$ 与残差 $\mathbf{e}(\mathbf{x}_k)$；（3）解 \cref{eq:gn-normal} 得 $\Delta\mathbf{x}_k$；（4）$\|\Delta\mathbf{x}_k\|$ 足够小则停，否则 $\mathbf{x}_{k+1}=\mathbf{x}_k\boxplus\Delta\mathbf{x}_k$ 返回 2。

\begin{pitfall}[GN 的半正定与病态——$\mathbf{H}\approx\mathbf{J}^\top\mathbf{W}^{-1}\mathbf{J}$ 近似的缺陷]
$\mathbf{J}^\top\mathbf{W}^{-1}\mathbf{J}$ 只\emph{半}正定，实际数据里可能\emph{奇异或病态（ill-conditioned）}\cite{gaoxiang2019slam14}：此时增量不稳定、算法不收敛，几何上意味着“局部不像二次函数”（对应 \cref{ch:slam_est} 的不可观测/gauge 自由度）。更严重：即使 $\boldsymbol{\Lambda}$ 非奇异、非病态，若 $\Delta\mathbf{x}$ \emph{太大}，\cref{eq:gn-linearize} 的线性近似失效，可能让目标\emph{变大}\cite{gaoxiang2019slam14,barfoot2024state}。根因是 GN 丢掉了 $\sum_i u_i\,\partial^2u_i/\partial\mathbf{x}\partial\mathbf{x}^\top$ 项：当残差\emph{大}（拟合差/外点）或问题\emph{强非线性}（$\partial^2u_i$ 大）时，该项不可略，近似海森失真，GN 可能失败。两个补丁：\emph{线搜索}（取 $\mathbf{x}\leftarrow\mathbf{x}\boxplus\alpha\Delta\mathbf{x}$，$\alpha\in(0,1]$，因 $\Delta\mathbf{x}$ 仍是下降方向）；以及更彻底的——LM（\cref{sec:nlopt-lm}）。
\end{pitfall}

\paragraph{完整数值例 1：手写高斯牛顿拟合曲线（十四讲 §6.3.1，全录）。}\cite{gaoxiang2019slam14}
拟合 $y=\exp(ax^2+bx+c)+w$（$w\sim\mathcal{N}(0,\sigma^2)$），待估 $a,b,c$。误差 $e_i=y_i-\exp(ax_i^2+bx_i+c)$，其对参数的导数
\begin{equation}
  \frac{\partial e_i}{\partial a}=-x_i^2\,\mathrm{e}^{(\cdot)},\quad
  \frac{\partial e_i}{\partial b}=-x_i\,\mathrm{e}^{(\cdot)},\quad
  \frac{\partial e_i}{\partial c}=-\mathrm{e}^{(\cdot)},\qquad \mathrm{e}^{(\cdot)}:=\exp(ax_i^2+bx_i+c),
  \label{eq:nlopt-curve-jac}
\end{equation}
即 $\mathbf{J}_i=[\partial e_i/\partial a,\partial e_i/\partial b,\partial e_i/\partial c]^\top$（$3\times1$）。增量方程（标量观测、信息权 $\sigma^{-2}$）
\begin{equation}
  \Big(\sum_{i=1}^{N}\mathbf{J}_i\,\sigma^{-2}\,\mathbf{J}_i^\top\Big)\Delta\mathbf{x}_k=\sum_{i=1}^{N}-\mathbf{J}_i\,\sigma^{-2}\,e_i.
  \label{eq:nlopt-curve-gn}
\end{equation}
（此处 $\mathbf{J}_i\mathbf{J}_i^\top$ 是 $3\times3$ 外积，与本书矩阵约定 $\sum_i\mathbf{J}_i^\top\sigma^{-2}\mathbf{J}_i$ 数值相同。）

\begin{codebox}[C++]{手写高斯牛顿曲线拟合（slambook2/ch6/gaussNewton.cpp，关键体）}
double ar=1.0, br=2.0, cr=1.0;       // 真实参数
double ae=2.0, be=-1.0, ce=5.0;      // 估计初值
int N=100; double w_sigma=1.0, inv_sigma=1.0/w_sigma;
// ... 生成带噪数据 x_data, y_data（y=exp(ar*x*x+br*x+cr)+gaussian noise）...
double cost=0, lastCost=0;
for (int iter=0; iter<100; iter++) {
  Matrix3d H = Matrix3d::Zero();     // Hessian = J^T * Sigma^-1 * J（GN 近似海森）
  Vector3d b = Vector3d::Zero();     // 右端 g = -J^T Sigma^-1 e
  cost = 0;
  for (int i=0; i<N; i++) {
    double xi=x_data[i], yi=y_data[i];
    double error = yi - exp(ae*xi*xi + be*xi + ce);   // 残差 e_i
    Vector3d J;                                        // 雅可比（列向量）
    J[0] = -xi*xi*exp(ae*xi*xi + be*xi + ce);          // de/da
    J[1] = -xi*exp(ae*xi*xi + be*xi + ce);             // de/db
    J[2] = -exp(ae*xi*xi + be*xi + ce);                // de/dc
    H += inv_sigma*inv_sigma * J * J.transpose();      // 累加 J Sigma^-1 J^T
    b += -inv_sigma*inv_sigma * error * J;             // 累加 -J Sigma^-1 e
    cost += error*error;
  }
  Vector3d dx = H.ldlt().solve(b);   // 用 LDLT(Cholesky 变体)解 H dx = b，绝不求逆
  if (isnan(dx[0])) break;           // 数值崩溃保护
  if (iter>0 && cost>=lastCost) break;  // 代价不再下降则停（朴素 GN 无信赖域回退）
  ae += dx[0]; be += dx[1]; ce += dx[2];
  lastCost = cost;
}
// 结果: estimated abc = 0.890912, 2.1719, 0.943629（真值 1,2,1，受噪声影响）
\end{codebox}

\noindent\textbf{运行输出与分析}\cite{gaoxiang2019slam14}：目标函数 9 次迭代后趋于收敛（更新量趋零），代价 $\sum e^2$ 从 $3.2\times10^6$ 降到 $101.94$；终值 $(a,b,c)\approx(0.8909,2.1719,0.9436)$ 与真值 $(1,2,1)$ 接近（差异来自噪声）。在 i7-8700 上约 \emph{0.2 ms}。注意手写 GN 没有信赖域回退，靠 \texttt{cost>=lastCost} 这一朴素准则止步——这正是 LM（\cref{sec:nlopt-lm}）要补的。

\begin{practice}
\begin{enumerate}
  \item 完整推导 \cref{eq:gn-normal}（两种途径），并写出多残差堆叠时 $\mathbf{J},\mathbf{W}$ 的分块。
  \item 说明 GN 丢掉的“残差 $\times$ 二阶导”项在何种问题上不可忽略（提示：强非线性 $+$ 大残差），并据此解释为何曲线拟合例 GN 收敛得很好。
  \item （概念）解释“为什么 $\mathbf{J}^\top\mathbf{W}^{-1}\mathbf{J}$ 不正定时解不稳定”，联系 \cref{ch:slam_est} 的可观测性。
  \item （编程）把上面 \texttt{H.ldlt().solve(b)} 换成显式 \texttt{H.inverse()*b}，在病态数据上观察数值差异。答案要点：求逆放大条件数，分解更稳（\cref{sec:nlopt-solve}）。
\end{enumerate}
\end{practice}

% ════════════════════════════════════════════════════════════════
\section{列文伯格–马夸尔特与信赖域}
\label{sec:nlopt-lm}

\paragraph{动机：GN 的近似只在小范围可信。}
GN 的线性化（\cref{eq:gn-linearize}）只在展开点附近有效，故给增量加一个\emph{信赖区域}（trust region）：区域内信任近似，区域外不信\cite{gaoxiang2019slam14,carlone2026handbook}。LM 因此被称\emph{阻尼牛顿法}（damped Newton），一般比 GN 更健壮（但可能更慢）。

\paragraph{增益比 $\rho$：度量近似好坏。}\cite{gaoxiang2019slam14,carlone2026handbook}
\begin{equation}
  \rho=\frac{J(\mathbf{x})-J(\mathbf{x}\boxplus\Delta\mathbf{x})}{L(\mathbf{0})-L(\Delta\mathbf{x})}
  =\frac{\text{实际下降}}{\text{模型预测下降}},
  \label{eq:lm-rho}
\end{equation}
$L(\Delta\mathbf{x})$ 为 $J$ 在当前点的二次（线性化）近似模型。$\rho\approx1$ 近似好（可放大信赖域）；$\rho$ 太小说明近似差（缩小信赖域、\emph{拒绝}该步）。
（两种分母对照。本书取 Handbook\cite{carlone2026handbook} 的标准版：分母是\emph{目标} $J$ 的\emph{二次模型}下降 $L(\mathbf{0})-L(\Delta\mathbf{x})$，其中二次模型 $L(\Delta\mathbf{x})=\tfrac12\|\mathbf{e}+\mathbf{J}\Delta\mathbf{x}\|_{\mathbf{W}^{-1}}^2=J(\mathbf{x})+(\nabla J)^\top\Delta\mathbf{x}+\tfrac12\Delta\mathbf{x}^\top\boldsymbol{\Lambda}\Delta\mathbf{x}$（梯度 $\nabla J=\mathbf{J}^\top\mathbf{W}^{-1}\mathbf{e}$、$\boldsymbol{\Lambda}=\mathbf{J}^\top\mathbf{W}^{-1}\mathbf{J}$）。另一种文献\cite{gaoxiang2019slam14} 把分母简化为目标的\emph{一阶}（梯度）预测下降：
\begin{equation}
  \rho=\frac{J(\mathbf{x})-J(\mathbf{x}\boxplus\Delta\mathbf{x})}{-\,(\nabla J)^\top\Delta\mathbf{x}}
  \qquad(\text{分子、分母均为\emph{标量目标}的下降量；下降步使 }(\nabla J)^\top\Delta\mathbf{x}<0),
  \label{eq:lm-rho-res}
\end{equation}
两式动机一致——都是“实际下降 / 模型预测下降”，只是标准版用二次模型 $L$（更准）、简化版用一阶预测（更省）。）

\begin{proposition}[LM 阻尼增量方程]\label{prop:lm}
带信赖域约束 $\|\mathbf{D}\Delta\mathbf{x}\|^2\le\mu$ 的子问题，经拉格朗日乘子 $\lambda$ 化为
\begin{equation}
  \big(\mathbf{J}^\top\mathbf{W}^{-1}\mathbf{J}+\lambda\mathbf{D}^\top\mathbf{D}\big)\Delta\mathbf{x}^*=-\mathbf{J}^\top\mathbf{W}^{-1}\mathbf{e}.
  \label{eq:lm}
\end{equation}
$\mathbf{D}=\mathbf{I}$（约束在球内）为 \emph{Levenberg}；$\mathbf{D}^\top\mathbf{D}=\operatorname{diag}(\mathbf{J}^\top\mathbf{W}^{-1}\mathbf{J})$（按各维梯度大小缩放）为 \emph{Marquardt}\cite{gaoxiang2019slam14,carlone2026handbook}。
\end{proposition}
\begin{proof}
信赖域子问题 $\min_{\Delta\mathbf{x}}\tfrac12\|\mathbf{e}+\mathbf{J}\Delta\mathbf{x}\|_{\mathbf{W}}^2$ s.t. $\|\mathbf{D}\Delta\mathbf{x}\|^2\le\mu$。构造拉格朗日函数
\[
  \mathcal{L}(\Delta\mathbf{x},\lambda)=\tfrac12\|\mathbf{e}+\mathbf{J}\Delta\mathbf{x}\|_{\mathbf{W}}^2+\tfrac{\lambda}{2}\big(\|\mathbf{D}\Delta\mathbf{x}\|^2-\mu\big),
\]
对 $\Delta\mathbf{x}$ 求导置零：$\mathbf{J}^\top\mathbf{W}^{-1}(\mathbf{e}+\mathbf{J}\Delta\mathbf{x})+\lambda\mathbf{D}^\top\mathbf{D}\Delta\mathbf{x}=\mathbf{0}$，即 \cref{eq:lm}。Barfoot\cite{barfoot2024state} 由“GN 不保证收敛”直接给同式 $(\mathbf{U}^\top\mathbf{U}+\lambda\mathbf{D})\delta\mathbf{x}^*=-\mathbf{U}^\top\mathbf{u}$；Handbook\cite{carlone2026handbook} 给 $(\mathbf{A}^\top\mathbf{A}+\lambda\mathbf{I})\Delta=\mathbf{A}^\top\mathbf{b}$（Levenberg）与 $\lambda\operatorname{diag}(\mathbf{A}^\top\mathbf{A})$（Marquardt），均与本式一致。\hfill$\square$
\end{proof}

\paragraph{$\lambda$ 在 GN 与最速下降间插值（三书一致的退化分析）。}\cite{gaoxiang2019slam14,barfoot2024state,carlone2026handbook}
\cref{eq:lm} 的两极行为：
\begin{itemize}
  \item $\lambda\to0$：$\mathbf{J}^\top\mathbf{W}^{-1}\mathbf{J}$ 主导，LM $\to$ \textbf{高斯牛顿}（二次近似好、收敛快）；
  \item $\lambda\to\infty$（取 $\mathbf{D}=\mathbf{I}$）：$\lambda\mathbf{I}$ 主导，$\Delta\mathbf{x}^*\approx-\tfrac1\lambda\mathbf{J}^\top\mathbf{W}^{-1}\mathbf{e}$，正是\textbf{最速下降}方向上的极小步（小步、稳）。
\end{itemize}
故 LM 在 GN 与 SD 间\emph{自然插值}。两种阻尼都可在贝叶斯意义上解释为\emph{给所有变量加一个零均值先验}\cite{carlone2026handbook}。Marquardt 的对角缩放使：在目标\emph{平坦}方向（梯度小、对角元小）走更大步、\emph{陡峭}方向走更小步。

\paragraph{接受/拒绝与 $\lambda$ 调度（Handbook 启发式）。}\cite{carlone2026handbook,gaoxiang2019slam14}
迭代里按 $\rho$ 调 $\lambda$：
\begin{itemize}
  \item 步被\emph{接受}（$\rho$ 大、残差下降）：\emph{减小} $\lambda$（如 $\div10$，趋近 GN、快），用新线性化点继续；
  \item 步被\emph{拒绝}（$\rho$ 小或目标变大）：\emph{增大} $\lambda$（如 $\times10$，趋近 SD、稳），\emph{重解}修改后的 \cref{eq:lm}。
\end{itemize}
十四讲 §6.2.3 给具体阈值：$\rho>\tfrac34$ 则 $\mu=2\mu$（放大信赖域），$\rho<\tfrac14$ 则 $\mu=0.5\mu$（缩小），$\rho$ 大于某阈值才接受步——倍数与阈值均为经验值。\textbf{LM 拒绝会让残差变大的步}，故在病态/强非线性下比 GN 稳健。工程准则：问题性质好用 GN，接近病态用 LM。

\paragraph{Powell 狗腿（Dog-Leg）：省一次分解。}\cite{carlone2026handbook}
LM 每次改 $\lambda$ 都要\emph{重新分解}修改后的矩阵（最贵环节）。Powell 狗腿改为：\emph{分别}算 GN 步与最速下降步，在信赖域内把两者折成一条“狗腿”折线（先沿 SD、到信赖域边界处\emph{急转}朝 GN，停在边界）。若组合步被拒，GN/SD 方向仍有效、只换组合比例，\textbf{每次状态更新只需一次分解}，故常比 LM 高效。Dog-Leg 维护\emph{显式}信赖域，用增益比 \cref{eq:lm-rho}（$\rho<0.25$ 缩小、$\rho>0.75$ 放大并接受）调半径。

\begin{algo}[列文伯格–马夸尔特（信赖域）]\label{alg:nlopt-lm}
给初值 $\mathbf{x}_0$、初始半径 $\mu$（或初始 $\lambda$）。重复：
(1) 在 $\mathbf{x}_k$ 算 $\mathbf{J},\mathbf{e}$；
(2) 解 \cref{eq:lm} 得 $\Delta\mathbf{x}_k$；
(3) 按 \cref{eq:lm-rho} 算 $\rho$；
(4) $\rho$ 大（如 $>\tfrac34$）则接受 $\mathbf{x}_{k+1}=\mathbf{x}_k\boxplus\Delta\mathbf{x}_k$ 并减小 $\lambda$/放大 $\mu$；$\rho$ 小（如 $<\tfrac14$）则拒绝、增大 $\lambda$/缩小 $\mu$、重解；
(5) 收敛（$\|\Delta\mathbf{x}\|$ 或 $\|\nabla J\|$ 小）则止。
\end{algo}

\begin{pitfall}[思维陷阱：“LM 的 $\lambda$ 越大越好”]
$\lambda$ 大趋最速下降（每步小、稳，但收敛\emph{慢}、可能停在远处）；$\lambda$ 小趋 GN（快，但二次近似差时可能发散）。没有“越大越好”——应\emph{按 $\rho$ 自适应}：近似好就减 $\lambda$ 提速，近似差就加 $\lambda$ 求稳。固定大 $\lambda$ 会让本可二次收敛的问题退化成龟速一阶下降。
\end{pitfall}

\begin{pitfall}[思维陷阱：把增益比 $\rho$ 与鲁棒核 $\rho$ 混为一谈]
本章两处都用了字母 $\rho$：\cref{eq:lm-rho} 的 $\rho$ 是\emph{标量增益比}（实际下降/预测下降，调信赖域用）；\cref{sec:nlopt-robust} 的 $\rho(u)$ 是\emph{鲁棒核函数}（把残差范数映成代价）。二者毫无关系，只是历史上撞了符号。读代码/论文时按上下文区分：前者是一个数、后者是一个函数。
\end{pitfall}

\begin{practice}
\begin{enumerate}
  \item 由信赖域子问题的拉格朗日函数完整推出 \cref{eq:lm}。
  \item 取 $\mathbf{D}=\mathbf{I}$，验证 $\lambda\to\infty$ 时增量趋于 $-\tfrac1\lambda\mathbf{J}^\top\mathbf{W}^{-1}\mathbf{e}$（最速下降）。
  \item （概念）说明 Dog-Leg 相比 LM “省一次矩阵分解”省在哪、为何 GN/SD 方向被拒后仍可复用。
  \item 写出 $\rho<0$ 的几何含义（实际\emph{不降反升}），并说明此时该如何调 $\lambda$。
\end{enumerate}
\end{practice}

% ════════════════════════════════════════════════════════════════
\section{解正规方程：白化、Cholesky 与 QR}
\label{sec:nlopt-solve}

每次迭代都要解一个线性系统 $\boldsymbol{\Lambda}\Delta\mathbf{x}=\mathbf{b}$（$\boldsymbol{\Lambda}=\mathbf{J}^\top\mathbf{W}^{-1}\mathbf{J}$，$\mathbf{b}=-\mathbf{J}^\top\mathbf{W}^{-1}\mathbf{e}$）。怎么解得又快又稳，是 BA 能否实时的关键\cite{carlone2026handbook}。

\paragraph{白化：把加权范数化成标准范数。}\cite{carlone2026handbook,barfoot2024state}
定义 $\boldsymbol{\Sigma}^{-1/2}$ 为 $\boldsymbol{\Sigma}$ 的矩阵平方根之逆，则马氏范数化为标准欧氏范数：
\begin{equation}
  \|\mathbf{e}\|_{\boldsymbol{\Sigma}}^2=\mathbf{e}^\top\boldsymbol{\Sigma}^{-1}\mathbf{e}=\big(\boldsymbol{\Sigma}^{-1/2}\mathbf{e}\big)^\top\big(\boldsymbol{\Sigma}^{-1/2}\mathbf{e}\big)=\big\|\boldsymbol{\Sigma}^{-1/2}\mathbf{e}\big\|_2^2.
  \label{eq:nlopt-whiten}
\end{equation}
对每个因子的雅可比与残差左乘 $\boldsymbol{\Sigma}_i^{-1/2}$，定义\emph{白化}量
\begin{equation}
  \mathbf{A}_i=\boldsymbol{\Sigma}_i^{-1/2}\mathbf{J}_i,\qquad \mathbf{b}_i=\boldsymbol{\Sigma}_i^{-1/2}\big(-\mathbf{e}_i\big),
  \label{eq:nlopt-Ai}
\end{equation}
堆叠成标准最小二乘 $\min_{\Delta\mathbf{x}}\|\mathbf{A}\Delta\mathbf{x}-\mathbf{b}\|_2^2$。Barfoot\cite{barfoot2024state} 用 Cholesky 实现：$\mathbf{u}=\mathbf{L}\mathbf{e}$，$\mathbf{L}^\top\mathbf{L}=\mathbf{W}^{-1}$，把 $\tfrac12\mathbf{e}^\top\mathbf{W}^{-1}\mathbf{e}$ 变成 $\tfrac12\mathbf{u}^\top\mathbf{u}$。白化\emph{消去量纲}（像素/米/弧度），不同观测可放进同一目标；标量时即“每项除以标准差 $\sigma_i$”\cite{carlone2026handbook}。

\paragraph{法方程与 Cholesky（平方根信息）。}\cite{carlone2026handbook,dellaert2017factor}
满秩时唯一解由\emph{法方程} $(\mathbf{A}^\top\mathbf{A})\Delta\mathbf{x}=\mathbf{A}^\top\mathbf{b}$ 给出；信息矩阵 $\boldsymbol{\Lambda}=\mathbf{A}^\top\mathbf{A}$ 对称正定，作 Cholesky 分解
\begin{equation}
  \boldsymbol{\Lambda}=\mathbf{A}^\top\mathbf{A}=\mathbf{R}^\top\mathbf{R}\quad(\mathbf{R}\ \text{上三角}),
  \qquad \mathbf{R}^\top\mathbf{y}=\mathbf{A}^\top\mathbf{b},\ \ \mathbf{R}\Delta\mathbf{x}=\mathbf{y}.
  \label{eq:cholesky}
\end{equation}
先\emph{前代}解 $\mathbf{y}$、再\emph{回代}解 $\Delta\mathbf{x}$。$\mathbf{R}$ 即\emph{平方根信息矩阵}——这族方法因此称 $\sqrt{}$SAM（square-root smoothing and mapping）\cite{dellaert2017factor}。
（Handbook\cite{carlone2026handbook} 把 Cholesky 因子记作上三角 $\mathbf{R}$；Golub--Van Loan 习惯下三角 $\mathbf{L}=\mathbf{R}^\top$。本书此处 $\mathbf{R}$ 指 Cholesky 因子，\emph{非}旋转矩阵。）

\paragraph{复杂度（Cholesky）。}\cite{carlone2026handbook}
稠密时 Cholesky 分解需 $n^3/3$ flops；含算对称 $\mathbf{A}^\top\mathbf{A}$ 的一半，整算法约 $(m+n/3)n^2$ flops（$\mathbf{A}$ 为 $m\times n$）。也可用 LDU/LDL$^\top$ 变体\emph{避免开平方根}（即代码里的 \texttt{ldlt()}）。

\paragraph{QR：更稳，不显式构造 $\boldsymbol{\Lambda}$。}\cite{carlone2026handbook}
对 $\mathbf{A}$ 直接做 QR（Householder 反射逐列消零）
\begin{equation}
  \mathbf{A}=\mathbf{Q}\begin{bmatrix}\mathbf{R}\\\mathbf{0}\end{bmatrix},\qquad \begin{bmatrix}\mathbf{d}\\\mathbf{e}'\end{bmatrix}=\mathbf{Q}^\top\mathbf{b},
  \label{eq:nlopt-qr}
\end{equation}
$\mathbf{Q}$ 为 $m\times m$ 正交阵（通常\emph{不显式构造}，把 $\mathbf{b}$ 作为附加列随 $\mathbf{A}$ 一起变换即得 $\mathbf{Q}^\top\mathbf{b}$）。由 $\mathbf{Q}$ 正交：
\[
  \|\mathbf{A}\Delta\mathbf{x}-\mathbf{b}\|_2^2=\|\mathbf{Q}^\top\mathbf{A}\Delta\mathbf{x}-\mathbf{Q}^\top\mathbf{b}\|_2^2=\|\mathbf{R}\Delta\mathbf{x}-\mathbf{d}\|_2^2+\|\mathbf{e}'\|_2^2,
\]
回代 $\mathbf{R}\Delta\mathbf{x}=\mathbf{d}$ 即解，$\|\mathbf{e}'\|_2^2$ 为残差平方和。QR 与 Cholesky 得到\emph{相同}的 $\mathbf{R}$（因 $\mathbf{A}^\top\mathbf{A}=[\mathbf{R}^\top\,\mathbf{0}]\mathbf{Q}^\top\mathbf{Q}[\mathbf{R};\mathbf{0}]=\mathbf{R}^\top\mathbf{R}$，至多对角符号差），但 QR \emph{不构造} $\mathbf{A}^\top\mathbf{A}$（避免条件数平方放大），\textbf{数值更稳}；代价 $2(m-n/3)n^2$，约为 Cholesky 的两倍\cite{carlone2026handbook}。

\begin{insight}[Cholesky vs QR 的取舍]
两者解\emph{同一}问题、给\emph{同一} $\mathbf{R}$。差别在数值：QR 直接作用于 $\mathbf{A}$，条件数 $\kappa(\mathbf{A})$；Cholesky 先形成 $\mathbf{A}^\top\mathbf{A}$，条件数变 $\kappa(\mathbf{A})^2$（病态问题精度损失加倍）。但 QR 慢约一倍。实践中：\emph{稀疏} Cholesky/LDL 因更快、且稀疏问题上优势\emph{不止常数倍}，是多数 SLAM 的默认；只在数值病态时才换 QR\cite{carlone2026handbook,dellaert2017factor}。
\end{insight}

\begin{note}
\textbf{为什么绝不直接求逆}：解 $\boldsymbol{\Lambda}\Delta\mathbf{x}=\mathbf{b}$ \emph{绝不}先算 $\boldsymbol{\Lambda}^{-1}$（稠密、$O(n^3)$、数值差、破坏稀疏），而是\emph{分解 $+$ 前后代}。手写实现里 Eigen 的 \texttt{ldlt().solve()} 即此意。需要某变量\emph{边缘协方差}时（\cref{sec:nlopt-cov}），也只回代 $\boldsymbol{\Lambda}^{-1}$ 的\emph{对应块}，不整体求逆。
\end{note}

\paragraph{历史：平方根滤波。}\cite{dellaert2017factor,carlone2026handbook}
分解信息矩阵的思想在序贯估计里叫\emph{平方根信息滤波}（SRIF），1969 年为 JPL 水手 10 号金星任务开发。Maybeck 名言：“许多实践者有充分理由主张平方根滤波应\emph{总是}优先于标准卡尔曼递推”——平方根带来更精确、更稳定的算法。$\sqrt{}$SAM\cite{dellaert2017factor} 把这一思想用到批量平滑。

\begin{practice}
\begin{enumerate}
  \item 推导 \cref{eq:nlopt-whiten}，并说明白化后为什么不同量纲的观测可以相加。
  \item 验证 QR 与 Cholesky 给出相同的 $\mathbf{R}$（用 $\mathbf{A}^\top\mathbf{A}=\mathbf{R}^\top\mathbf{R}$）。
  \item （概念）一个条件数 $\kappa(\mathbf{A})=10^4$ 的问题，用 Cholesky 解会损失约几位有效数字？为什么 QR 更好？
\end{enumerate}
\end{practice}

% ════════════════════════════════════════════════════════════════
\section{稀疏性：SLAM 为何可解}
\label{sec:nlopt-sparse}

BA/位姿图的信息矩阵动辄上万、上百万维，稠密分解 $O(n^3)$ 不可行。SLAM 可解，全靠\emph{稀疏}\cite{carlone2026handbook,gaoxiang2019slam14,dellaert2017factor}。

\paragraph{稀疏来自因子图结构。}\cite{dellaert2017factor,carlone2026handbook}
每个残差只关联一两个变量（\cref{thm:nlopt-map} 结构特点 1），故白化雅可比 $\mathbf{A}$ 中\textbf{每个因子对应一块行、每个变量对应一块列}——$\mathbf{A}$ 的块结构\emph{镜像因子图}。于是 $\boldsymbol{\Lambda}=\mathbf{A}^\top\mathbf{A}$ 的稀疏模式恰是因子图诱导的\emph{无向图（马尔可夫随机场 MRF）}的邻接矩阵：两变量曾\emph{共现}于某因子 $\Rightarrow$ $\boldsymbol{\Lambda}$ 对应块非零；一个 $n$ 元因子在其全部变量间诱导一个\emph{团（clique）}（成对因子是其特例）。Handbook\cite{carlone2026handbook} 强调：$\mathbf{A}^\top\mathbf{A}$ \emph{不是真海森}，而是 GN 近似（由截断残差泰勒级数得到，呼应 \cref{thm:gn}）。

\paragraph{BA 的箭头结构与 Schur 消元。}\cite{gaoxiang2019slam14}
BA 状态分相机 $\mathbf{x}_c$（$6m$ 维，$m$ 个相机）与路标 $\mathbf{x}_p$（$3n$ 维，$n\gg m$）。因每条观测只连“一个相机 $+$ 一个路标”，$\boldsymbol{\Lambda}$ 有\emph{箭头（arrowhead）}结构：
\begin{equation}
  \begin{bmatrix}\mathbf{B}&\mathbf{E}\\\mathbf{E}^\top&\mathbf{C}\end{bmatrix}
  \begin{bmatrix}\Delta\mathbf{x}_c\\\Delta\mathbf{x}_p\end{bmatrix}
  =\begin{bmatrix}\mathbf{v}\\\mathbf{w}\end{bmatrix},
  \label{eq:ba-block}
\end{equation}
其中 $\mathbf{B}$（相机块）、$\mathbf{C}$（路标块）均为\emph{块对角}阵——$\mathbf{C}$ 由一个个 $3\times3$ 小块组成，求逆极廉价。左乘消元矩阵 $\big[\begin{smallmatrix}\mathbf{I}&-\mathbf{E}\mathbf{C}^{-1}\\\mathbf{0}&\mathbf{I}\end{smallmatrix}\big]$ 消去 $\mathbf{E}$，得只含相机增量的\emph{Schur 补}方程：
\begin{equation}
  \boxed{\ \underbrace{\big(\mathbf{B}-\mathbf{E}\mathbf{C}^{-1}\mathbf{E}^\top\big)}_{\textstyle\mathbf{S}}\,\Delta\mathbf{x}_c=\mathbf{v}-\mathbf{E}\mathbf{C}^{-1}\mathbf{w}\ },\qquad
  \Delta\mathbf{x}_p=\mathbf{C}^{-1}\big(\mathbf{w}-\mathbf{E}^\top\Delta\mathbf{x}_c\big).
  \label{eq:schur}
\end{equation}
此过程称\emph{边缘化}（marginalization）/ Schur 消元\cite{gaoxiang2019slam14}。

\begin{insight}[$O(\text{相机}^3)$ 而非 $O(n^3)$；且 $\mathbf{S}$ 比 $\mathbf{H}$ 更好条件]
直接分解全 $\boldsymbol{\Lambda}$ 是 $O\big((6m+3n)^3\big)$，被海量路标 $3n$ 拖垮。Schur 后，主成本是分解 $\mathbf{S}$（仅 $6m\times6m$），约 $O\big((6m)^3\big)$；$\mathbf{C}^{-1}$ 因块对角是 $O(n)$ 的廉价回代。\textbf{复杂度由“相机$+$路标总维立方”降到“相机维立方”}——这正是 BA 能在成百上千路标下实时的根本。\emph{即便没有速度收益，Schur 仍值得做}：消去路标块后的 Schur 补 $\mathbf{S}=\mathbf{B}-\mathbf{E}\mathbf{C}^{-1}\mathbf{E}^\top$ 的条件数通常\emph{小于}原 $\mathbf{H}$\cite{gaoxiang2019slam14}，求解更稳、迭代更快收敛——“消元同时改善了条件数”。Triggs 等\cite{triggs2000ba} 把这套“先消结构、得只含相机的约化系统”称作\emph{约化相机系统}（reduced camera system），是摄影测量 BA 的标准技巧。\textbf{对偶地，消去\emph{相机}变量（而非路标）也是 SLAM 中采用的手段}：当相机少而每相机看大量独占路标时，先消相机得只含路标的约化系统反而更省——消谁取决于谁的块对角更大、谁的分隔集更小（即下文“变量消元”的一般化视角）。
\end{insight}

\paragraph{Fill-in 与共视。}\cite{gaoxiang2019slam14,dellaert2017factor,carlone2026handbook}
$\mathbf{S}=\mathbf{B}-\mathbf{E}\mathbf{C}^{-1}\mathbf{E}^\top$ \emph{不再}保持 $\mathbf{B}$ 的对角稀疏：消元把“两相机经由公共路标的间接耦合”显式填进 $\mathbf{S}$ 的相机--相机块，称\emph{填入}（fill-in）——即\emph{超出原矩阵稀疏模式之外的额外非零元}\cite{carlone2026handbook}。其几何意义是\textbf{共视}（co-visibility）：$\mathbf{S}$ 的非对角非零块 $\iff$ 两相机看到\emph{共同路标}。
\emph{变量排序}极大影响 fill-in：任何排序给出\emph{相同}解，但决定 $\mathbf{R}$ 的非零元个数；\textbf{最小化 fill-in 的排序是 NP 难}\cite{dellaert2017factor,carlone2026handbook}，实践用 \emph{COLAMD}（column approximate minimum degree）等启发式。Handbook\cite{carlone2026handbook} 给一个真实例（图 1.9，$333$ 未知、$1126$ 行）：自然排序（位姿在前、路标在后）的 Cholesky 因子有 $9399$ 个非零元，COLAMD 重排后只剩 $4168$ 个——\emph{解完全相同}，只是分解快得多。

\paragraph{消元 $=$ 稀疏分解 $=$ 图模型（变量消元算法）。}\cite{dellaert2017factor,carlone2026handbook}
线性高斯下，“按某顺序逐变量消元”\emph{精确等价}于稀疏 Cholesky/QR 分解，也等价于把因子图转成有向的\emph{贝叶斯网}。消一个变量 $\mathbf{x}_j$：把所有与它相邻的因子相乘成 $\psi(\mathbf{x}_j,\mathbf{s}_j)$（$\mathbf{s}_j$ 为其\emph{分隔集}=消元后 $\mathbf{x}_j$ 所条件的变量集），再分解为条件 $p(\mathbf{x}_j\mid\mathbf{s}_j)$ 与 $\mathbf{s}_j$ 上的新因子 $\tau(\mathbf{s}_j)$。矩阵层面这一步用\emph{部分 QR}：把积因子的增广矩阵 $[\bar{\mathbf{A}}_j\mid\bar{\mathbf{b}}_j]$ 分解为 $\big[\begin{smallmatrix}\mathbf{R}_j&\mathbf{T}_j&\mathbf{d}_j\\\mathbf{0}&\tilde{\mathbf{A}}_\tau&\tilde{\mathbf{b}}_\tau\end{smallmatrix}\big]$，上行给条件（贡献一行上三角 $\mathbf{R}$），下行给新因子 $\tau$。\textbf{Schur 补正是“消去一个变量、更新其分隔集”的矩阵特例}：\cref{eq:schur} 里消去路标块 $\mathbf{C}$、更新相机块 $\mathbf{B}\to\mathbf{S}$，恰是变量消元一步。消完所有变量得贝叶斯网（弦图/三角化），其条件高斯
\begin{equation}
  \|\mathbf{R}_j\mathbf{x}_j+\mathbf{T}_j\mathbf{s}_j-\mathbf{d}_j\|_2^2=\|\mathbf{x}_j-\boldsymbol{\mu}_j\|_{\boldsymbol{\Sigma}_j}^2,\quad
  \boldsymbol{\mu}_j=\mathbf{R}_j^{-1}(\mathbf{d}_j-\mathbf{T}_j\mathbf{s}_j),\quad \boldsymbol{\Sigma}_j=(\mathbf{R}_j^\top\mathbf{R}_j)^{-1},
  \label{eq:nlopt-elim-cond}
\end{equation}
按\emph{逆消元序}回代即得各变量 MAP $\mathbf{x}_j^*=\mathbf{R}_j^{-1}(\mathbf{d}_j-\mathbf{T}_j\mathbf{s}_j^*)$。这条 $\boldsymbol{\Sigma}_j=(\mathbf{R}_j^\top\mathbf{R}_j)^{-1}$ 正是 \cref{sec:nlopt-cov} 协方差恢复的雏形。把这套增量化（Bayes 树、iSAM2：新观测只触及受影响的局部子树、并\emph{选择性重线性化}偏离工作点的变量）留给后端 / 增量 SLAM 专章（\cref{ch:vio} 及后端章），其库实现是 GTSAM\cite{dellaert2017factor}；本章只需记住：\textbf{稀疏 Schur/Cholesky/变量消元是 BA 可扩展的引擎，其结构由因子图直接给出}。

\begin{pitfall}[思维陷阱：以为变量排序会改变解]
不同排序给出\emph{完全相同}的 MAP 解（COLAMD 的 $4168$ 非零元版与自然序的 $9399$ 版回代结果一致）——排序只改变 fill-in（中间因子的非零元数），从而改变\emph{速度}，不改变\emph{答案}。所以可以放心用任何启发式排序加速；选错排序只会变慢、不会变错。反过来，\emph{别}为了“解更对”去手调排序，那是无用功。
\end{pitfall}

\begin{pitfall}[思维陷阱：把滤波的“边缘化致稠密”和平滑的“消元致 fill-in”当成一回事]
两者都叫“边缘化”，但后果不同。\emph{滤波}（EKF）边缘化掉\emph{历史全部}状态，使信息矩阵随时间\emph{不可逆地变稠密}\cite{dellaert2017factor}——信息一旦糊进当前块就拿不回来。\emph{平滑}（$\sqrt{}$SAM）保留整条轨迹，消元只在\emph{求解时}产生 fill-in，且信息矩阵\emph{始终稀疏}（这正是平滑优于滤波的计算理由）。SWF（滑动窗口滤波）则在窗口边界做边缘化，是二者的折中——其矩阵实质（边缘化 $=$ Schur 补、以及为何“新增帧平凡、删旧帧破坏稀疏”）见章末 \cref{der:nlopt-swf}，窗扩/窗缩的逐式递归块更新与滑动算法见 \cref{der:nlopt-swf-recursion,alg:nlopt-swf}。
\end{pitfall}

\begin{practice}
\begin{enumerate}
  \item 对一个“2 相机 $+$ 6 路标、各看 4 个点”的小 BA，画出 $\boldsymbol{\Lambda}$ 的箭头稀疏图案，并标出哪些块来自共视。
  \item 由 \cref{eq:ba-block} 推出 \cref{eq:schur}（左乘 $\big[\begin{smallmatrix}\mathbf{I}&-\mathbf{E}\mathbf{C}^{-1}\\\mathbf{0}&\mathbf{I}\end{smallmatrix}\big]$）。
  \item （概念）解释“边缘化路标”为何在相机间引入 fill-in，并联系 \cref{ch:eskf} 的“滤波边缘化历史致信息阵稠密”。
  \item 说明为什么“消去一个变量并更新其分隔集”就是 Schur 补的一步。
\end{enumerate}
\end{practice}

% ════════════════════════════════════════════════════════════════
\section{光束法平差：从重投影误差到稀疏求解}
\label{sec:nlopt-ba}

\paragraph{桥接：把前几节的“抽象残差”落到视觉 SLAM 最核心的问题上。}
到此为止，\cref{sec:nlopt-map,sec:nlopt-gn,sec:nlopt-lm} 给了求解\emph{任意}最小二乘的通法，\cref{sec:nlopt-solve,sec:nlopt-sparse} 给了数值与稀疏的工具。但抽象的 $\mathbf{e}(\mathbf{x})$ 在视觉 SLAM 里有一个最重要、最具体的化身\cite{gaoxiang2019slam14}——\emph{光束法平差}（Bundle Adjustment，BA）。本节把前几节的工具\emph{兑现}到 BA 上：先写出重投影误差与 BA 代价函数（\cref{sec:nlopt-map} 的 $J$ 的具体形式），再给出它对位姿/路标的雅可比块结构（\cref{thm:gn} 的 $\mathbf{J}$ 的具体形状），由此自然导出 \cref{sec:nlopt-sparse} 的箭头稀疏与 Schur，最后落到 Ceres/g2o 的 BA 实践。读完本节，\cref{eq:ba-block} 的抽象分块就有了血肉。

\paragraph{什么是 BA：让光束收束。}\cite{gaoxiang2019slam14}
BA 一词源自摄影测量：从特征点发出的几束光线（bundles of rays）在多个相机的成像平面上变成像素；\emph{调整}（adjustment）各相机位姿与各路标的三维位置，使这些光线尽量\emph{收束}回相机光心（即让预测像素与实测像素吻合），就是 BA。它同时优化\emph{相机}与\emph{结构}，是视觉三维重建（SfM）与 SLAM 后端的核心。

\paragraph{重投影的完整流水线（逐步，十四讲 §9.2.1 全录）。}\cite{gaoxiang2019slam14}
把世界点 $\mathbf{p}$ 投到像素，需四步，每步都是后面求雅可比的环节：
\begin{enumerate}
  \item \textbf{世界$\to$相机}（用外参 $\mathbf{R},\mathbf{t}$，即位姿 $\mathbf{T}$）：$\mathbf{P}'=\mathbf{R}\mathbf{p}+\mathbf{t}=[X',Y',Z']^\top$；
  \item \textbf{投到归一化平面}：$\mathbf{P}_c=[u_c,v_c,1]^\top=[X'/Z',\,Y'/Z',\,1]^\top$；
  \item \textbf{径向畸变}（系数 $k_1,k_2$，$r_c^2=u_c^2+v_c^2$）：$u_c'=u_c(1+k_1r_c^2+k_2r_c^4)$，$v_c'$ 同理；
  \item \textbf{内参成像}：$u_s=f_x u_c'+c_x$，$v_s=f_y v_c'+c_y$。
\end{enumerate}
整条流水线就是观测方程 $\mathbf{z}=h(\mathbf{x},\mathbf{y})=h(\mathbf{T},\mathbf{p})$（\cref{ch:camera}）。其中 $\mathbf{x}$ 是相机位姿（外参 $\mathbf{R},\mathbf{t}$，李群 $\mathbf{T}$、李代数 $\boldsymbol{\xi}$），$\mathbf{y}$ 是三维点 $\mathbf{p}$，观测是像素 $\mathbf{z}=[u_s,v_s]^\top$。

\paragraph{BA 代价函数：\cref{eq:nlopt-batch-ls} 在视觉下的具体形式。}\cite{gaoxiang2019slam14}
单次观测的重投影误差就是\emph{实测像素减预测像素}：
\begin{equation}
  \mathbf{e}_{ij}=\mathbf{z}_{ij}-h(\mathbf{T}_i,\mathbf{p}_j),
  \label{eq:nlopt-ba-reproj}
\end{equation}
$\mathbf{z}_{ij}$ 是在位姿 $\mathbf{T}_i$ 观测路标 $\mathbf{p}_j$ 产生的像素。把所有观测加权求和，得 BA 目标
\begin{equation}
  \min_{\mathbf{T},\mathbf{p}}\ \tfrac12\sum_{i=1}^{m}\sum_{j=1}^{n}\big\|\mathbf{z}_{ij}-h(\mathbf{T}_i,\mathbf{p}_j)\big\|_{\boldsymbol{\Sigma}_{ij}}^2,
  \label{eq:nlopt-ba-cost}
\end{equation}
这正是 \cref{eq:nlopt-batch-ls} 只保留观测项、且 $h$ 取重投影的特例——\textbf{对它求解就是“同时调整全部位姿与全部路标”，即 BA}。注意视觉里观测数远多于运动约束（一帧数百特征点），故 BA 由观测项主导。

\paragraph{雅可比的块结构：每条边只碰两个变量。}\cite{gaoxiang2019slam14}
把全部变量按 $\mathbf{x}=[\mathbf{T}_1,\dots,\mathbf{T}_m,\mathbf{p}_1,\dots,\mathbf{p}_n]^\top$ 堆叠，给增量后线性化（\cref{eq:gn-linearize} 的 BA 版）：
\begin{equation}
  \tfrac12\|\mathbf{e}(\mathbf{x}\boxplus\Delta\mathbf{x})\|^2\approx\tfrac12\sum_{i,j}\big\|\mathbf{e}_{ij}+\mathbf{F}_{ij}\,\Delta\boldsymbol{\xi}_i+\mathbf{E}_{ij}\,\Delta\mathbf{p}_j\big\|^2,
  \label{eq:nlopt-ba-lin}
\end{equation}
$\mathbf{F}_{ij}=\partial\mathbf{e}_{ij}/\partial\boldsymbol{\xi}_i$（$2\times6$，对位姿，用 \cref{ch:lie} 右扰动雅可比）、$\mathbf{E}_{ij}=\partial\mathbf{e}_{ij}/\partial\mathbf{p}_j$（$2\times3$，对路标，二者具体形式见 \cref{ch:vo}）。关键在于：\emph{$\mathbf{e}_{ij}$ 只描述“$\mathbf{T}_i$ 看到 $\mathbf{p}_j$”这一件事}，故其整体雅可比
\begin{equation}
  \mathbf{J}_{ij}=\Big(\mathbf{0},\dots,\underbrace{\tfrac{\partial\mathbf{e}_{ij}}{\partial\boldsymbol{\xi}_i}}_{2\times6},\dots,\mathbf{0},\dots,\underbrace{\tfrac{\partial\mathbf{e}_{ij}}{\partial\mathbf{p}_j}}_{2\times3},\dots,\mathbf{0}\Big)
  \label{eq:nlopt-ba-jij}
\end{equation}
\emph{除了 $i,j$ 两处全为零}。把全部 $\mathbf{F}_{ij},\mathbf{E}_{ij}$ 拼成大雅可比 $\mathbf{J}=[\mathbf{F}\ \mathbf{E}]$（$\mathbf{F}$ 为对位姿块、$\mathbf{E}$ 为对路标块），增量方程 $\mathbf{H}\Delta\mathbf{x}=\mathbf{g}$ 的 $\mathbf{H}=\mathbf{J}^\top\mathbf{W}^{-1}\mathbf{J}$ 即
\begin{equation}
  \mathbf{H}=\begin{bmatrix}\mathbf{F}^\top\mathbf{W}^{-1}\mathbf{F}&\mathbf{F}^\top\mathbf{W}^{-1}\mathbf{E}\\\mathbf{E}^\top\mathbf{W}^{-1}\mathbf{F}&\mathbf{E}^\top\mathbf{W}^{-1}\mathbf{E}\end{bmatrix}
  \;\xrightarrow{\text{记号}}\;\begin{bmatrix}\mathbf{B}&\mathbf{E}\\\mathbf{E}^\top&\mathbf{C}\end{bmatrix},
  \label{eq:nlopt-ba-H}
\end{equation}
\textbf{这正是 \cref{eq:ba-block} 的 $\mathbf{B},\mathbf{C},\mathbf{E}$ 的来源}（$\mathbf{B}=\mathbf{F}^\top\mathbf{W}^{-1}\mathbf{F}$ 相机块、$\mathbf{C}=\mathbf{E}^\top\mathbf{W}^{-1}\mathbf{E}$ 路标块）。LM 只是把 $\mathbf{H}$ 换成 $\mathbf{H}+\lambda\mathbf{I}$（\cref{eq:lm}），结构不变。

\begin{insight}[$\mathbf{H}$ 的稀疏 $=$ 因子图的邻接]
单条边 $\mathbf{e}_{ij}$ 对 $\mathbf{H}$ 的贡献 $\mathbf{J}_{ij}^\top\mathbf{W}^{-1}\mathbf{J}_{ij}$ \emph{只有四个非零块}：$(i,i),(i,j),(j,i),(j,j)$。累加 $\mathbf{H}=\sum_{ij}\mathbf{J}_{ij}^\top\mathbf{W}^{-1}\mathbf{J}_{ij}$ 后：$\mathbf{B}$（相机块）必为\emph{块对角}（不同相机间无直接边）、$\mathbf{C}$（路标块）也\emph{块对角}（不同路标间无直接边），只有 $\mathbf{E}$（相机--路标）按“谁看到谁”稀疏分布。于是 \textbf{$\mathbf{H}$ 的非对角非零块 $\iff$ 图 \cref{sec:nlopt-sparse} 中两变量间有边}——这把 \cref{sec:nlopt-sparse} 抽象的“稀疏镜像因子图”落成了可画的箭头矩阵（\cref{fig:nlopt-ba-sparsity}）。
\end{insight}

\paragraph{例：2 相机 6 路标的箭头矩阵（十四讲 §9.2.3 全录）。}\cite{gaoxiang2019slam14}
设相机 $C_1$ 看到 $P_1,P_2,P_3,P_4$、$C_2$ 看到 $P_3,P_4,P_5,P_6$，则
\begin{equation}
  J(\mathbf{x})=\tfrac12\big(\|\mathbf{e}_{11}\|^2+\|\mathbf{e}_{12}\|^2+\|\mathbf{e}_{13}\|^2+\|\mathbf{e}_{14}\|^2+\|\mathbf{e}_{23}\|^2+\|\mathbf{e}_{24}\|^2+\|\mathbf{e}_{25}\|^2+\|\mathbf{e}_{26}\|^2\big).
  \label{eq:nlopt-ba-26}
\end{equation}
变量顺序 $(\boldsymbol{\xi}_1,\boldsymbol{\xi}_2,\mathbf{p}_1,\dots,\mathbf{p}_6)$ 下，$\mathbf{e}_{11}$ 的雅可比只在 $\boldsymbol{\xi}_1$ 与 $\mathbf{p}_1$ 列非零，其余八列全零；八条边各类似。叠出的 $\mathbf{H}$ 是 $8\times8$ 块阵：左上 $2\times2$（相机--相机）由共视填充、右下 $6\times6$（路标--路标）严格块对角、右上/左下（相机--路标）按观测稀疏。

\begin{figure}[H]\centering
\begin{tikzpicture}[scale=0.42, every node/.style={font=\scriptsize}]
  % 8x8 grid; cell(col,rowFromTop): col 0..7, rowFromTop 0..7 -> y=7-row
  \def\cell#1#2{\fill[main!55] (#1,7-#2) rectangle (#1+1,8-#2);}   % 对角/相机块(深)
  \def\cellp#1#2{\fill[third!55] (#1,7-#2) rectangle (#1+1,8-#2);} % 相机-路标耦合(浅)
  % camera diagonal (C1C1,C2C2) + 共视耦合 (C1C2,C2C1)
  \cell{0}{0}\cell{1}{1}\cell{0}{1}\cell{1}{0}
  % landmark diagonal P1..P6 (idx 2..7)
  \cell{2}{2}\cell{3}{3}\cell{4}{4}\cell{5}{5}\cell{6}{6}\cell{7}{7}
  % C1(row/col 0) sees P1..P4 (idx 2..5)
  \cellp{2}{0}\cellp{3}{0}\cellp{4}{0}\cellp{5}{0}
  \cellp{0}{2}\cellp{0}{3}\cellp{0}{4}\cellp{0}{5}
  % C2(row/col 1) sees P3..P6 (idx 4..7)
  \cellp{4}{1}\cellp{5}{1}\cellp{6}{1}\cellp{7}{1}
  \cellp{1}{4}\cellp{1}{5}\cellp{1}{6}\cellp{1}{7}
  % grid lines
  \foreach \i in {0,...,8} { \draw[gray!40] (\i,0)--(\i,8); \draw[gray!40] (0,\i)--(8,\i); }
  % column labels (top) and row labels (left)
  \foreach \k/\t in {0/{C_1},1/{C_2},2/{P_1},3/{P_2},4/{P_3},5/{P_4},6/{P_5},7/{P_6}} {
    \node at (\k+0.5,8.5) {$\t$};
    \node at (-0.6,7.5-\k) {$\t$};
  }
\end{tikzpicture}
\caption{2 相机 6 路标的 $\mathbf{H}$ 矩阵箭头（镐形）稀疏：深色（相机块 $\mathbf{B}$）左上小而稠、浅色为相机--路标耦合 $\mathbf{E}$、右下路标块 $\mathbf{C}$ 严格块对角。$(C_1,C_2)$ 块非零因二者共视 $P_3,P_4$。}
\label{fig:nlopt-ba-sparsity}
\end{figure}

\noindent 当 $n\gg m$（路标远多于相机），$\mathbf{H}$ 左上角小、右下对角块占绝大部分，整体形如\emph{箭头}（arrow-like），亦俗称\emph{镐形}（pickaxe）矩阵\cite{gaoxiang2019slam14}。对它解 $\mathbf{H}\Delta\mathbf{x}=\mathbf{g}$ 时，正是 \cref{eq:schur} 的 Schur 消元（边缘化路标块 $\mathbf{C}$）把代价压到 $O(\text{相机}^3)$；消元后 $\mathbf{S}$ 的非对角非零块即\emph{共视}（co-visibility）（\cref{sec:nlopt-sparse}）。\textbf{无须重复推导——\cref{eq:schur} 已给出全过程，本节只是把它的 $\mathbf{B},\mathbf{C},\mathbf{E}$ 认成 BA 的相机/路标/观测块}。

\paragraph{图优化语言：顶点$=$变量，边$=$误差。}\cite{gaoxiang2019slam14}
目标函数 \cref{eq:nlopt-ba-cost} 只写了“有哪些变量、哪些误差项”，却没直观显示\emph{谁连谁}。\emph{图优化}（graph optimization）把它画成图：\textbf{顶点}（vertex）表优化变量（位姿、路标），\textbf{边}（edge）表误差项（一条边连它所依赖的全部顶点）。BA 的重投影边连\emph{一个相机$+$一个路标}（二元边）；曲线拟合（\cref{sec:nlopt-practice}）的边只连一个顶点（一元边）；位姿图（\cref{sec:nlopt-posegraph}）的边连两个位姿。一条边可连一个、两个或多个顶点，故又称\emph{超边}、整图称\emph{超图}（hypergraph）。这与 \cref{sec:nlopt-sparse} 的因子图是同一结构的两种叫法（边$\leftrightarrow$因子）。\textbf{图优化不是新理论，只是把 \cref{eq:nlopt-ba-H} 的稀疏结构画出来、便于建模与利用}——g2o 正是按此设计。

\paragraph{实践一：Ceres 解 BA（BAL 数据集 $+$ SPARSE\_SCHUR $+$ Huber，十四讲 §9.3）。}\cite{gaoxiang2019slam14}
BAL（Bundle Adjustment in the Large）数据集每个相机用 $9$ 维参数（$3$ 轴角旋转 $+3$ 平移 $+1$ 焦距 $+2$ 畸变 $k_1,k_2$），路标 $3$ 维；其投影平面在光心\emph{之后}，故投影后乘 $-1$。用前先 \texttt{Normalize}（路标中心置零、尺度归一，提升数值稳定）。残差就是 \cref{eq:nlopt-ba-reproj} 的两维重投影差。

\begin{codebox}[C++]{Ceres BA：重投影残差与求解（slambook2/ch9/SnavelyReprojectionError，关键体）}
struct SnavelyReprojectionError {
  SnavelyReprojectionError(double ox, double oy):observed_x(ox),observed_y(oy){}
  template<typename T>
  bool operator()(const T* camera, const T* point, T* residuals) const {
    // camera[0-2]=angle-axis, [3-5]=t, [6]=focal, [7-8]=k1,k2
    T p[3];
    ceres::AngleAxisRotatePoint(camera, point, p);      // R*point  (Rodrigues)
    p[0]+=camera[3]; p[1]+=camera[4]; p[2]+=camera[5];  // + t
    T xp=-p[0]/p[2], yp=-p[1]/p[2];                     // BAL: 投影面在光心后, 取负
    T r2=xp*xp+yp*yp;
    T distortion=T(1.0)+r2*(camera[7]+camera[8]*r2);    // 1 + k1 r^2 + k2 r^4
    T focal=camera[6];
    residuals[0]=focal*distortion*xp - T(observed_x);   // 预测 - 实测 (两维)
    residuals[1]=focal*distortion*yp - T(observed_y);
    return true;
  }
  double observed_x, observed_y;
};
// SolveBA: 每条观测加一个残差块, 连一个 camera(9维)+一个 point(3维)
for (int i=0;i<bal.num_observations();++i){
  auto* cost=new ceres::AutoDiffCostFunction<SnavelyReprojectionError,2,9,3>(
               new SnavelyReprojectionError(obs[2*i], obs[2*i+1]));
  ceres::LossFunction* loss=new ceres::HuberLoss(1.0); // 鲁棒核: IRLS+Huber 抑外点
  problem.AddResidualBlock(cost, loss, camera_ptr, point_ptr);
}
ceres::Solver::Options options;
options.linear_solver_type = ceres::SPARSE_SCHUR;  // 自动对路标做 Schur 边缘化
ceres::Solve(options, &problem, &summary);
// 16 相机 / 22106 点 / 83718 观测: 误差从 1.84e7 迭代下降; 输出 final.ply 可视化
\end{codebox}

\noindent 选 \texttt{SPARSE\_SCHUR} 即让 Ceres 按本节的 Schur（先边缘化路标）求解；但 Ceres \emph{自动}决定边缘化谁，用户不指定。\texttt{HuberLoss(1.0)} 即对每条边套 Huber 核（\cref{sec:nlopt-robust} 的 IRLS），自动降权误匹配。

\paragraph{实践二：g2o 解 BA（自定义顶点/边 $+$ 手动 setMarginalized，十四讲 §9.4）。}\cite{gaoxiang2019slam14}
g2o 把相机（位姿$+$内参，$9$ 维）与路标（$3$ 维）建成两类顶点、重投影建成二元边。\textbf{与 Ceres 的关键差别：g2o 用稀疏 Schur 时必须\emph{手动}对路标顶点 \texttt{setMarginalized(true)}}，否则运行报错。

\begin{codebox}[C++]{g2o BA：位姿内参顶点的右乘更新与边缘化设置（slambook2/ch9，关键体）}
// 9 维相机顶点: SO3 旋转 + t + focal + k1 + k2
class VertexPoseAndIntrinsics : public g2o::BaseVertex<9, PoseAndIntrinsics> {
public:
  virtual void oplusImpl(const double* update) override {
    // 位姿不在向量空间: 旋转用右扰动(此处源码左乘 Exp), 见 sec manifold
    _estimate.rotation = Sophus::SO3d::exp(
        Eigen::Vector3d(update[0],update[1],update[2])) * _estimate.rotation;
    _estimate.translation += Eigen::Vector3d(update[3],update[4],update[5]);
    _estimate.focal += update[6]; _estimate.k1 += update[7]; _estimate.k2 += update[8];
  }
};
class VertexPoint : public g2o::BaseVertex<3, Eigen::Vector3d> {   // 路标: 向量空间, 直接加
  virtual void oplusImpl(const double* u) override { _estimate += Eigen::Vector3d(u); }
};
// 二元边: 观测 2 维, 连 [相机, 路标]; computeError = project(point) - measurement
// 求解器: 块求解 <9,3> + LinearSolverCSparse + LM
for (auto* point_vertex : point_vertices)
  point_vertex->setMarginalized(true);     // !! 必须手动: 告诉 g2o 对路标做 Schur 边缘化
for (each observation) {
  auto* e = new EdgeProjection;
  e->setVertex(0, camera_vertex); e->setVertex(1, point_vertex);
  e->setMeasurement(Vector2d(u,v));
  e->setInformation(Matrix2d::Identity());           // 信息矩阵 = Sigma^-1
  e->setRobustKernel(new g2o::RobustKernelHuber());  // 鲁棒核 = IRLS + Huber
  optimizer.addEdge(e);
}
optimizer.optimize(40);
\end{codebox}

\noindent 顶点 \texttt{oplusImpl} 即 $\mathbf{x}\boxplus\Delta\mathbf{x}$：路标在向量空间故直接加，\emph{位姿不在向量空间故用 $\mathrm{Exp}$ 乘}（\cref{sec:nlopt-manifold}）。\texttt{setInformation} 即白化/信息加权（\cref{eq:nlopt-whiten}）；\texttt{setRobustKernel} 即 \cref{sec:nlopt-robust} 的鲁棒核。两库优化后点云一致、目标函数数值可不同（系数 $\tfrac12$ 与核处理差异），但 Meshlab 里效果相同。

\begin{pitfall}[g2o BA 忘了 \texttt{setMarginalized} $\to$ 运行时报错]
g2o 用稀疏 Schur 解 BA 时，\emph{必须}对要边缘化的路标顶点显式 \texttt{setMarginalized(true)}，否则求解器不知该消去谁、运行报错（与 Ceres 自动选边缘化不同）。这不是“可选优化”，是\emph{必需声明}。反过来，若把相机也误标成边缘化，则破坏 \cref{eq:schur} 的“先消路标块 $\mathbf{C}$”前提、得不到箭头加速。
\end{pitfall}

\begin{practice}
\begin{enumerate}
  \item 由重投影流水线（世界$\to$相机$\to$归一化$\to$畸变$\to$像素）链式求出 $\mathbf{F}_{ij}=\partial\mathbf{e}_{ij}/\partial\boldsymbol{\xi}_i$ 的骨架（右扰动，细节见 \cref{ch:vo}），说明它为何是 $2\times6$。
  \item 对 \cref{fig:nlopt-ba-sparsity} 的 2 相机 6 路标例，写出 $\mathbf{B},\mathbf{C},\mathbf{E}$ 各自的块维数，并指出 $(C_1,C_2)$ 块非零的物理原因。
  \item （编程）在 Ceres BA 上把 \texttt{SPARSE\_SCHUR} 换成 \texttt{DENSE\_SCHUR} 或 \texttt{SPARSE\_NORMAL\_CHOLESKY}，对比 BAL \texttt{problem-16-22106} 的耗时。
  \item （概念）解释为什么 g2o BA 必须 \texttt{setMarginalized(true)} 而 Ceres 不必，联系 \cref{eq:schur}。
\end{enumerate}
\end{practice}

% ════════════════════════════════════════════════════════════════
\section{位姿图优化}
\label{sec:nlopt-posegraph}

\paragraph{桥接：当路标收敛后，BA 还能再瘦身。}\cite{gaoxiang2019slam14}
\cref{sec:nlopt-ba} 的 BA 同时优化位姿与路标，但一个常见观察是：经过若干次观测后，\emph{收敛的路标位置几乎不再变化}，对它们反复优化“费力不讨好”；而路标数又远多于相机（实时 BA 即便用稀疏也只能到几万点）\cite{gaoxiang2019slam14}。于是有了更省的近似——\emph{位姿图}（pose graph）：优化\emph{几次后把路标固定}，只当作位姿间的约束，从此\textbf{只优化相机轨迹、不再优化路标}。本节给出位姿图的误差、雅可比与实践，它是 \cref{sec:nlopt-ba} 的“去路标”简化，也是大尺度建图与回环（\cref{ch:vo}）、激光 SLAM（\cref{ch:lidar_slam}）的主力。

\paragraph{位姿图的结构。}\cite{gaoxiang2019slam14}
顶点$=$相机位姿 $\mathbf{T}_1,\dots,\mathbf{T}_n$；边$=$两位姿间的\emph{相对运动}估计 $\Delta\mathbf{T}_{ij}$，可来自特征点法/直接法/GPS/IMU 积分。记观测到的相对位姿
\begin{equation}
  \mathbf{T}_{ij}=\mathbf{T}_i^{-1}\mathbf{T}_j\quad(\text{即 }\Delta\boldsymbol{\xi}_{ij}=\ln(\mathbf{T}_i^{-1}\mathbf{T}_j)^\vee),
  \label{eq:nlopt-pg-meas}
\end{equation}
因噪声此式不精确成立，故构造误差（把观测 $\mathbf{T}_{ij}$ 移到右侧）
\begin{equation}
  \mathbf{e}_{ij}=\ln\!\big(\mathbf{T}_{ij}^{-1}\,\mathbf{T}_i^{-1}\mathbf{T}_j\big)^\vee.
  \label{eq:nlopt-pg-err}
\end{equation}
误差为零 $\iff$ 估计的相对位姿与观测吻合。

\paragraph{雅可比：用伴随把扰动“搬”到一侧（逐步）。}\cite{gaoxiang2019slam14}
对 $\mathbf{T}_i,\mathbf{T}_j$ 各加扰动求导。本书主线用右扰动；为与常见实现（如不少 g2o 位姿图代码）对照，这里先给\emph{左}扰动版 $\delta\boldsymbol{\xi}_i,\delta\boldsymbol{\xi}_j$（两套经伴随 $\mathrm{Ad}$ 互转，\cref{ch:lie}），误差变
\begin{equation}
  \hat{\mathbf{e}}_{ij}=\ln\!\big(\mathbf{T}_{ij}^{-1}\mathbf{T}_i^{-1}\exp((-\delta\boldsymbol{\xi}_i)^\wedge)\exp(\delta\boldsymbol{\xi}_j^\wedge)\mathbf{T}_j\big)^\vee.
  \label{eq:nlopt-pg-pert}
\end{equation}
两个扰动夹在中间，无法直接用 BCH。\emph{关键技巧}：用伴随性质 $\exp(\boldsymbol{\xi}^\wedge)\mathbf{T}=\mathbf{T}\exp\big((\mathrm{Ad}(\mathbf{T}^{-1})\boldsymbol{\xi})^\wedge\big)$（\cref{ch:lie}）把扰动\emph{挪到最右}，再用 $\ln(\cdot)\approx\mathbf{e}_{ij}+\mathcal{J}_r^{-1}(\cdot)$ 一阶展开，得两个雅可比：
\begin{equation}
  \frac{\partial\mathbf{e}_{ij}}{\partial\delta\boldsymbol{\xi}_i}=-\mathcal{J}_r^{-1}(\mathbf{e}_{ij})\,\mathrm{Ad}(\mathbf{T}_j^{-1}),\qquad
  \frac{\partial\mathbf{e}_{ij}}{\partial\delta\boldsymbol{\xi}_j}=\mathcal{J}_r^{-1}(\mathbf{e}_{ij})\,\mathrm{Ad}(\mathbf{T}_j^{-1}).
  \label{eq:nlopt-pg-jac}
\end{equation}
$\mathfrak{se}(3)$ 上的 $\mathcal{J}_r^{-1}$ 形式繁琐，工程上常用近似（误差接近零时 $\mathcal{J}_r^{-1}\approx\mathbf{I}$，或保留一阶项 $\mathcal{J}_r^{-1}\approx\mathbf{I}+\tfrac12\big[\begin{smallmatrix}\boldsymbol{\phi}_e^\wedge&\boldsymbol{\rho}_e^\wedge\\\mathbf{0}&\boldsymbol{\phi}_e^\wedge\end{smallmatrix}\big]$，$\mathbf{e}=[\boldsymbol{\rho}_e;\boldsymbol{\phi}_e]$）。

\begin{insight}[左扰动雅可比要配左乘更新——本书主线用右扰动]
\cref{eq:nlopt-pg-jac} 是\emph{左}扰动版（配 \texttt{oplusImpl} 里的左乘 $\mathrm{Exp}(\Delta\boldsymbol{\xi})\mathbf{T}$），文献中常见此形\cite{gaoxiang2019slam14}。本书主线用\emph{右}扰动（\cref{sec:nlopt-manifold} 铁律：雅可比侧$=$更新侧）；引用本式时须经伴随 $\mathrm{Ad}$ 转成右扰动（\cref{ch:lie}）再用，或干脆按右扰动重推（习题）。无论左右，\textbf{误差接近零时 $\mathcal{J}_r^{-1}\approx\mathbf{I}$，故两套实现收敛到几乎相同结果}——这解释了下文“近似为 $\mathbf{I}$ 也能收敛”的现象。
\end{insight}

\paragraph{位姿图目标函数与求解。}\cite{gaoxiang2019slam14}
记 $\mathcal{E}$ 为全部边，目标
\begin{equation}
  \min\ \tfrac12\sum_{(i,j)\in\mathcal{E}}\mathbf{e}_{ij}^\top\boldsymbol{\Sigma}_{ij}^{-1}\mathbf{e}_{ij}.
  \label{eq:nlopt-pg-cost}
\end{equation}
这又是一个最小二乘，照旧用 GN/LM（\cref{sec:nlopt-gn,sec:nlopt-lm}）求解，唯一特别处是优化变量为李群位姿、用 \cref{eq:nlopt-pg-jac} 的雅可比。\textbf{位姿图比 BA 简单（无路标），但其稀疏性取决于运动}：直线前进得带状（稀疏）位姿图、来回环绕（loopy）得稠密位姿图——无运动先验时难再深挖结构。

\paragraph{实践：g2o 优化“球”位姿图。}\cite{gaoxiang2019slam14}
一个经典演示用 g2o 自带的 \texttt{create\_sphere} 仿真一个球形轨迹（$2500$ 位姿、$9799$ 边：相邻为\emph{里程计边}、层间为\emph{回环边}），加噪后初值偏离球面，优化目标是“把它掰回圆球”\cite{gaoxiang2019slam14}。\texttt{sphere.g2o} 文本里顶点为 \texttt{VERTEX\_SE3:QUAT}（四元数$+$平移）、边为 \texttt{EDGE\_SE3:QUAT}（相对位姿$+$上三角信息矩阵）。两种实现：

\begin{codebox}[C++]{g2o 位姿图：内置 SE3 顶点 vs 李代数自定义顶点（slambook2/ch10，关键体）}
// (A) 用 g2o 内置 SE3 顶点/边 (四元数表示), 仅配置求解器即可
typedef g2o::BlockSolver<g2o::BlockSolverTraits<6,6>> BlockSolverType;
auto solver = new g2o::OptimizationAlgorithmLevenberg(/* ... LinearSolverEigen ... */);
// 读 sphere.g2o: VERTEX_SE3:QUAT -> g2o::VertexSE3; EDGE_SE3:QUAT -> g2o::EdgeSE3
v->setFixed(true);            // 固定 0 号顶点作基准 (消 gauge 自由度)
optimizer.optimize(30);       // LM 30 次, chi2 从 1.0e9 降到 ~4.5e4, 形状回到圆球

// (B) 用 Sophus 李代数自定义顶点/边, 误差与雅可比按本节推导
class VertexSE3LieAlgebra : public g2o::BaseVertex<6, Sophus::SE3d> {
  virtual void oplusImpl(const double* update) override {        // 左乘更新
    Vector6d u; u << update[0],...,update[5];
    _estimate = Sophus::SE3d::exp(u) * _estimate;                // Exp(dxi) * T
  }
};
class EdgeSE3LieAlgebra : public g2o::BaseBinaryEdge<6,Sophus::SE3d,
                              VertexSE3LieAlgebra,VertexSE3LieAlgebra> {
  virtual void computeError() override {
    auto v1=...; auto v2=...;                                    // T_i, T_j
    _error = (_measurement.inverse()*v1.inverse()*v2).log();     // e_ij = ln(T_ij^-1 T_i^-1 T_j)
  }
  virtual void linearizeOplus() override {
    Matrix6d J = JRInv(Sophus::SE3d::exp(_error));               // J_r^{-1}(e) 近似
    _jacobianOplusXi = -J * v2.inverse().Adj();                  // 对 T_i  (eq pg-jac)
    _jacobianOplusXj =  J * v2.inverse().Adj();                  // 对 T_j
  }
};
// JRInv: J = I + 0.5 * [[phi^, rho^],[0, phi^]], 即 eq nlopt-pg-jac 的一阶近似
\end{codebox}

\noindent 结果：李代数版迭代约 $23$ 次误差稳定在 $\approx44360$，略优于内置 SE3 版 $30$ 次的 $\approx44811$——\textbf{李代数参数化用更少迭代得更好结果}；即便把 $\mathcal{J}_r^{-1}$ 近似为 $\mathbf{I}$ 也收敛到相近值（误差接近零时雅可比本就近 $\mathbf{I}$）。注意：后端位姿图\emph{无须实时}（PTAM 以来前后端分线程，前端 $30$ fps、后端慢跑即可），故 $2500$ 顶点跑数十秒可接受。

\begin{pitfall}[动手体会雅可比的作用：把 $\mathcal{J}_r^{-1}$ 近似为 $\mathbf{I}$ 会怎样]
位姿图收敛时残差趋零、$\mathcal{J}_r^{-1}\approx\mathbf{I}$，故\emph{把 \cref{eq:nlopt-pg-jac} 里的 $\mathcal{J}_r^{-1}$ 直接换成 $\mathbf{I}$}（甚至索性\emph{注释掉 \texttt{linearizeOplus}}、让 g2o 退回数值求导）通常\emph{仍能收敛}到相近结果。但这不等于雅可比无用：近似为 $\mathbf{I}$ 会让\emph{远离最优}时的下降方向不准、迭代次数变多、对大初值误差更脆弱。\textbf{建议亲手对比}三种实现——(a) 解析 $\mathcal{J}_r^{-1}$ 一阶近似（\cref{eq:nlopt-pg-jac}）、(b) $\mathcal{J}_r^{-1}\equiv\mathbf{I}$、(c) 注释 \texttt{linearizeOplus} 用数值求导——看迭代次数与终值差异，便能直观体会“雅可比准不准”如何换取“收敛快不快、稳不稳”\cite{gaoxiang2019slam14}。误区：因为“近似 $\mathbf{I}$ 也能收敛”就以为可永远偷懒——一旦初值差或问题强非线性（如大尺度回环），准确雅可比的价值立刻显现。
\end{pitfall}

\begin{pitfall}[位姿图不固定基准帧 $\to$ gauge 自由度致 $\mathbf{H}$ 奇异]
\cref{eq:nlopt-pg-cost} 只含\emph{相对}位姿约束，整张图可整体平移旋转而代价不变（gauge 自由度），$\mathbf{H}$ 因此奇异、GN/LM 解不稳。对策：\texttt{setFixed(true)} 固定一个顶点（或加一个绝对位姿先验边）作基准——这正是 \cref{sec:nlopt-why} pitfall“初值/可观测”与故障排查“$\boldsymbol{\Lambda}$ 奇异”在位姿图上的具体表现。
\end{pitfall}

\begin{practice}
\begin{enumerate}
  \item 用\emph{右}扰动重推 \cref{eq:nlopt-pg-jac}（更新 $\mathbf{T}\,\mathrm{Exp}(\delta\boldsymbol{\xi})$），与左扰动版对比（十四讲第 10 讲习题 1--2）。
  \item （概念）位姿图相比 \cref{sec:nlopt-ba} 的 BA 省了什么、代价是什么？何时该用位姿图、何时该用全 BA？
  \item 解释为什么误差接近零时 $\mathcal{J}_r^{-1}\approx\mathbf{I}$，并由此说明“近似为 $\mathbf{I}$ 也收敛”。
  \item （编程）参照 g2o 版在 Ceres 里实现“球”位姿图优化，对比迭代次数与终值。
\end{enumerate}
\end{practice}

% ════════════════════════════════════════════════════════════════
\section{流形上的优化（右扰动）}
\label{sec:nlopt-manifold}

SLAM 的变量含位姿 $\mathbf{T}\in\SE(3)$，不在向量空间，普通加法 $\mathbf{x}+\Delta\mathbf{x}$ 非法（破坏正交约束，\cref{thm:nlopt-map} 结构特点 2）。\cref{ch:lie} 已给出出路：在流形上用\emph{右扰动}定义增量。本节把 GN/LM 干净地搬到流形——其余推导\emph{完全照旧}，只换“增量定义”与“雅可比”。

\paragraph{$\boxplus/\boxminus$ 与 retraction。}\cite{hertzberg2013integrating}
用 \emph{$\boxplus$}（retraction）把局部向量增量“贴”回流形：对位姿
\begin{equation}
  \mathbf{T}\boxplus\delta\boldsymbol{\xi}:=\mathbf{T}\,\mathrm{Exp}(\delta\boldsymbol{\xi}),\qquad
  \mathbf{T}_2\boxminus\mathbf{T}_1:=\mathrm{Log}(\mathbf{T}_1^{-1}\mathbf{T}_2),
  \label{eq:retraction}
\end{equation}
即\emph{右乘}一个小扰动 $\mathrm{Exp}(\delta\boldsymbol{\xi})$（本书右扰动约定，$\delta\boldsymbol{\xi}=[\delta\boldsymbol{\rho};\delta\boldsymbol{\phi}]$）。普通向量分量的 $\boxplus$ 就是加法。Hertzberg 等\cite{hertzberg2013integrating} 用 $\boxplus/\boxminus$ 把“流形上的优化”封装成“切空间的欧氏优化 $+$ retraction”，使通用优化器无需懂流形细节。

\paragraph{流形高斯牛顿（GN shortcut 在李群上的落地）。}\cite{barfoot2024state,gaoxiang2019slam14}
优化变量改为“当前估计 $\hat{\mathbf{x}}$ 上的局部增量 $\delta\mathbf{x}$”。用 GN 的 shortcut（\cref{sec:nlopt-gn}）：先对残差关于 $\delta\mathbf{x}$ 线性化，
\begin{equation}
  \mathbf{e}(\hat{\mathbf{x}}\boxplus\delta\mathbf{x})\approx\mathbf{e}(\hat{\mathbf{x}})+\mathbf{J}\,\delta\mathbf{x},\qquad
  \mathbf{J}=\frac{\partial\,\mathbf{e}(\hat{\mathbf{x}}\boxplus\delta\mathbf{x})}{\partial\delta\mathbf{x}}\bigg|_{\delta\mathbf{x}=\mathbf{0}},
  \label{eq:manifold-jac}
\end{equation}
$\mathbf{J}$ 即\emph{右扰动雅可比}（如 $\partial(\mathbf{R}\mathbf{p})/\partial\delta\boldsymbol{\phi}=-\mathbf{R}\mathbf{p}^\wedge$，见 \cref{ch:lie}）。其余完全照旧：解 \cref{eq:gn-normal}（或 LM \cref{eq:lm}）得 $\delta\mathbf{x}^*$，再\emph{右乘}更新 $\hat{\mathbf{x}}\leftarrow\hat{\mathbf{x}}\boxplus\delta\mathbf{x}^*$。增量 $\delta\boldsymbol{\xi}$ 是无约束的 $\mathbb{R}^6$ 向量，正规方程因此是\emph{无约束}最小二乘——这正是 \cref{ch:lie} 强调的“李代数把带约束优化变无约束”。

\begin{insight}
求解器里的 “local parameterization / manifold / \texttt{oplusImpl}” 就是 \cref{eq:retraction} 的右乘更新。GTSAM 的 \texttt{Pose3}、Ceres 的 \texttt{Manifold}（旧称 \texttt{LocalParameterization}）、g2o 顶点的 \texttt{oplusImpl} 都要你提供 $\boxplus$；理解了右扰动，这些不再是黑话，而是“在流形切空间做 GN、再 retract 回去”。Barfoot\cite{barfoot2024state} 的“批量误差线性化 $\to$ 解正规方程”推导对欧氏与流形\emph{同一套}，区别只在雅可比是普通导数还是右扰动雅可比。
\end{insight}

\begin{pitfall}[思维陷阱：雅可比侧与更新侧不一致]
若雅可比按\emph{右}扰动算（$\partial/\partial\delta\boldsymbol{\xi}$，$\mathbf{T}\,\mathrm{Exp}(\delta\boldsymbol{\xi})$），更新却写成\emph{左}乘 $\mathrm{Exp}(\delta\boldsymbol{\xi})\,\mathbf{T}$（或反之），增量方向与几何不符，迭代不收敛甚至发散。\textbf{铁律：雅可比侧 $=$ 更新侧 $=$ 全程同侧}（本书统一右扰动）。引用 Barfoot/十四讲的左扰动公式时，须经伴随 $\mathrm{Ad}$ 转换到右扰动（\cref{ch:lie}）再用。
\end{pitfall}

\begin{practice}
\begin{enumerate}
  \item 把 \cref{sec:nlopt-gn} 的曲线拟合 GN 改成位姿图 GN：写出残差 $\mathbf{e}=\mathrm{Log}(\mathbf{T}_{ij}^{-1}\mathbf{T}_i^{-1}\mathbf{T}_j)$ 对 $\mathbf{T}_i,\mathbf{T}_j$ 的右扰动雅可比骨架。
  \item （概念）为什么 $\delta\boldsymbol{\xi}\in\mathbb{R}^6$ 上的正规方程是“无约束”的，而直接优化 $\mathbf{T}$ 的 16 个元素不是？
  \item 解释 Ceres 的 \texttt{Manifold::Plus} 与本节 $\boxplus$ 的对应关系。
\end{enumerate}
\end{practice}

% ════════════════════════════════════════════════════════════════
\section{鲁棒化：前端剔除与后端鲁棒核}
\label{sec:nlopt-robust}

真实数据有\emph{外点}（outlier）：按噪声模型\emph{极不可能}的观测（常见判据：偏离均值超过\emph{三个标准差}）\cite{barfoot2024state}。本节从“外点从哪来”讲起，给出抑制外点的\emph{两道防线}——\textbf{前端几何剔除}（\cref{sec:nlopt-robust-frontend}：RANSAC 与成对一致性最大化）与\textbf{后端鲁棒化}（\cref{sec:nlopt-robust-mest} 起：M-估计/鲁棒核/IRLS），并\emph{证明}“鲁棒 $=$ 逆 Wishart 先验下的 MAP”、点明其为 Black--Rangarajan 对偶的一个实例，再扩到渐进非凸 GNC。这套后端鲁棒化\emph{对滤波与优化通用}：\cref{ch:eskf} 已在滤波语境给过同一族鲁棒核（Huber/Cauchy/GM）、IRLS 膨胀协方差与“鲁棒 $=$ 逆 Wishart MAP”的要点（\cref{eq:robust-kernels,eq:irls-weight,thm:robust-iw}），本节是其\emph{完整}版（含 Tukey 全家福 \cref{tab:robust} 与 Sherman--Morrison 全步证明），两处字母（$\boldsymbol{\Sigma}_i$ 基协方差、$\mathbf{Y}_i$ 膨胀协方差）已统一，读者可视作同一内容的“要点版 + 完整版”。

\paragraph{外点从哪来：数据关联错误与模型假设违背。}\cite{carlone2026handbook}
上一章把 SLAM 钉成最小二乘时假设每条观测都服从零均值高斯。实践中此假设常被两类原因破坏。其一是\emph{数据关联}（data association）错误：观测 $\mathbf{z}_{ij}$ 由前端预处理原始传感数据得到（如对图像做特征检测与匹配、再算出对路标的方位），而检测/匹配会出错——被标记为路标 $\ell_j$ 的检测实际可能是另一个路标，使 $\mathbf{z}_{ij}$ 大幅偏离模型。地标 SLAM 里这表现为误匹配，位姿图回环里表现为\emph{回环误检}：回环靠地点识别（place recognition）判断两位姿是否看到同一场景，而当前方法易受\emph{感知混叠}（perceptual aliasing）所误——两处长得像的地方（办公楼里两间相似的隔间、两条相似的走廊）被错判为同一处，把回环约束接到错误的位姿对上。其二是\emph{模型假设违背}：多数 SLAM 假设路标静止，故对动态物体的检测（即便关联正确）也会成为大残差外点；传感器退化（轮速计打滑、镜头积灰）同理。无论哪类，外点都是 SLAM 里\emph{不可避免}的。

\paragraph{为何外点毁掉最小二乘：二次代价放大大残差。}\cite{gaoxiang2019slam14,barfoot2024state,carlone2026handbook}
高斯噪声是\emph{轻尾}的，本质上排除了“大误差”的可能；而二范数对大误差\emph{增长太快}（$\rho=\tfrac12u^2$，影响函数 $\psi=\rho'=u$ 无界）。残差被平方后，外点项在目标 $\sum_i\mathbf{r}_i^2$ 里有\emph{不成比例}的权重：一条误匹配边误差大、梯度大，优化会优先迁就这条“无理的边”、抹平其他正确边，整条轨迹被拉飞。在 M3500、SubT、Victoria Park 三个基准上即可直观看到：仅 $15\%$ 随机外点就使位姿图/地标 SLAM 的估计轨迹完全错乱，且外点常\emph{暴露}环境的感知混叠（如把两条相似走廊误并成一条）\cite{carlone2026handbook}。这正是 \cref{thm:nlopt-map} 结构特点 3 埋下的伏笔。出路有二，互补使用：\textbf{前端}在送进后端前用几何约束\emph{硬剔除}外点（\cref{sec:nlopt-robust-frontend}）；\textbf{后端}把残余外点\emph{软降权}（\cref{sec:nlopt-robust-mest}）。前端剔除是\emph{必要非充分}——几何/一致性约束即便在真值处成立也只是内点的必要条件，且近似求解会漏检——故残余外点仍会毒化二次代价，后端鲁棒化不可或缺。

% ----------------------------------------------------------------
\subsection{前端外点剔除：RANSAC 与成对一致性最大化（PCM）}
\label{sec:nlopt-robust-frontend}

前端的职责是从原始传感数据里抽出中间表示或\emph{伪观测}（再转成后端的因子）：通常先算一批初始观测（含大量外点），再后处理剔除外点\cite{carlone2026handbook}。本子节给两种前端剔除法。

\paragraph{RANSAC：把剔除看成共识最大化。}\cite{fischler1981ransac,carlone2026handbook}
RANSAC（random sample consensus，Fischler--Bolles 1981）是地标 SLAM 系统的关键部件，其原理立在两条洞见上。\emph{洞见一：内点须满足几何约束}。以双目/单目视觉 SLAM 的 2D-2D 对应为例，静止 3D 点随相机运动在两帧间的像素位移不能任意，须满足\emph{对极约束}（\cref{ch:vo}）：对标定相机，$\mathbf{z}_i(k{-}1)^\top\big([\mathbf{t}_k^{k-1}]_\times\mathbf{R}_k^{k-1}\big)\mathbf{z}_i(k)=0$。一般地，第 $i$ 个对应 $\mathbf{z}_i$ 须满足
\begin{equation}
  C(\mathbf{z}_i,\mathbf{x})\le\gamma,
  \label{eq:nlopt-geomconstraint}
\end{equation}
$C$ 为几何约束（可含未知状态 $\mathbf{x}$），$\gamma$ 是吸收噪声的松弛阈值。\emph{洞见二：内点是“彼此认同同一 $\mathbf{x}$”的最大子集}。于是剔除外点可写成\emph{共识最大化}（consensus maximization）
\begin{equation}
  \mathsf{S}^*_{\mathrm{CM}}=\arg\max_{\mathbf{x},\,\mathsf{S}\subset\mathsf{M}}|\mathsf{S}|
  \quad\text{s.t.}\quad C(\mathbf{z}_i,\mathbf{x})\le\gamma,\ \forall i\in\mathsf{S},
  \label{eq:nlopt-consensusmax}
\end{equation}
$\mathsf{M}$ 为初始对应集、$|\mathsf{S}|$ 为子集势。\cref{eq:nlopt-consensusmax} 只涉及一小块状态（对极约束只用两帧相对位姿，非整条轨迹），但仍是\emph{组合难}问题，与前端的实时要求冲突，故实践用启发式近似解。RANSAC 即最知名的近似：它假设 $\mathbf{x}$ 低维、可由一个\emph{最小集}（minimal set）经\emph{最小求解器}（minimal solver）估出（如标定相机的相对运动可由 Nistér 五点法\cite{nister2004fivepoint} 从 $5$ 对像素求得，至多相差一个尺度），于是不必穷举子集，只需迭代：（1）随机采 $n$ 个对应（$n=$ 最小集大小）；（2）用最小求解器算 $\hat{\mathbf{x}}$；（3）取满足 $C(\mathbf{z}_i,\hat{\mathbf{x}})\le\gamma$ 的对应组成\emph{共识集} $\mathsf{S}$，比历史更大则留存。重复至共识集足够大或达最大迭代数。设单点为内点的概率为 $w$，找到一个全内点最小集的期望迭代数是 $1/w^n$；要以概率 $p$ 选出全内点子集，需迭代数
\begin{equation}
  k=\frac{\ln(1-p)}{\ln(1-w^n)},
  \label{eq:nlopt-ransac}
\end{equation}
（由 $1-p=(1-w^n)^k$ 解得）。$n{=}5,w{=}0.7$ 时期望迭代数不到 $10$，加上最小求解器极快，使 RANSAC 在“外点少、最小集小”时极具吸引力，且顺带给出可作后端初值的 $\hat{\mathbf{x}}$。\textbf{但其期望迭代数随外点增多或最小集增大而爆炸}：$n{=}10,w{=}0.1$ 时需 $1/0.1^{10}=10^{10}$ 次，提前终止则多半返回错误解。位姿图 SLAM 正是反例——最小集须含至少 $N{-}1$ 条边（一棵生成树），$N$ 动辄上千。

\paragraph{成对一致性最大化（PCM）：用最大团剔除外点。}\cite{mangelson2018pcm,carlone2026handbook}
感知混叠严重的环境外点率可达 $90\%$，且并非所有问题都有小最小集的快速求解器。一条替代思路\emph{不}采样最小集，而是用图论找出\emph{彼此内部一致}的最大观测子集。其关键是：许多问题可定义\emph{一致性函数} $F$，判断\emph{一对}观测是否相容，且 $F$ \emph{与状态无关}、可预计算。两例：
\begin{itemize}
  \item \textbf{3D-3D 对应（RGB-D）}：两帧的 3D 特征点理应是同一组静止点的不同视角，故一对对应点的\emph{距离应随时间不变}：
  \begin{equation}
    \big|\,\|\mathbf{z}_i(k{-}1)-\mathbf{z}_j(k{-}1)\|-\|\mathbf{z}_i(k)-\mathbf{z}_j(k)\|\,\big|\le\gamma.
    \label{eq:nlopt-pcm-3d3d}
  \end{equation}
  \item \textbf{位姿图回环}（设里程计可靠、回环可能含外点）：无噪声时沿图中任一环路复合位姿应得单位阵，故一对回环 $\mathbf{T}_b^a$、$\mathbf{T}_d^c$ 须满足
  \begin{equation}
    \operatorname{dist}\!\big(\mathbf{T}_b^a\,\bar{\mathbf{T}}_c^b\,\mathbf{T}_d^c\,\bar{\mathbf{T}}_a^d,\ \mathbf{I}\big)\le\gamma,
    \label{eq:nlopt-pcm-loop}
  \end{equation}
  $\bar{\mathbf{T}}_c^b,\bar{\mathbf{T}}_a^d$ 为里程计链、$\operatorname{dist}$ 度量与单位阵的偏离，$\gamma$ 按环路长度取（长环累积更多噪声）。
\end{itemize}
一般地，成对一致性约束写作 $F(\mathbf{z}_i,\mathbf{z}_j)\le\gamma$；它\emph{不依赖} $\mathbf{x}$，故可对每对 $(i,j)$ 预先核验、无需最小求解器。于是有
\begin{equation}
  \mathsf{S}^*_{\mathrm{PCM}}=\arg\max_{\mathsf{S}\subset\mathsf{M}}|\mathsf{S}|
  \quad\text{s.t.}\quad F(\mathbf{z}_i,\mathbf{z}_j)\le\gamma,\ \forall i,j\in\mathsf{S}.
  \label{eq:nlopt-pcm}
\end{equation}
此即\emph{成对一致性最大化}（pairwise consistency maximization，PCM）。它仍是组合问题，但比 \cref{eq:nlopt-consensusmax} 略易——不含 $\mathbf{x}$、约束可预计算，且\emph{有图论解法}。把每个观测当一个节点、当 $F(\mathbf{z}_i,\mathbf{z}_j)\le\gamma$ 时在 $i,j$ 间连边，得\emph{一致性图}；\cref{eq:nlopt-pcm} 即求“每对节点都相邻的最大子集”——恰是图论的\emph{最大团}（maximum clique）：
\begin{equation}
  \mathsf{S}^*_{\mathrm{PCM}}=\text{一致性图的最大团}.
  \label{eq:nlopt-pcm-clique}
\end{equation}
最大团问题 NP 难且难近似，但已有大量精确（分支定界，最坏指数）与启发式（利用结构、快但不保证最优）算法，部分可在大稀疏图上并行快速求解。\textbf{PCM 与 RANSAC 的权衡}：RANSAC 用最小求解器的估计核验\emph{单个}观测的一致性，PCM 核验观测\emph{两两之间}的一致性。位姿图回环里，一对回环各自“看似合理”不等于二者\emph{互相}一致，故 PCM 远胜 RANSAC（多机 SLAM 尤甚）；但 3D-3D 这类问题里 \cref{eq:nlopt-pcm-3d3d} 比 RANSAC 更宽松、可能漏判外点。一致性图稠密时精确最大团慢，宜用启发式；而 2D-2D 对应因缺乏合适不变量，难设计通用一致性函数。

\begin{insight}[前端硬剔除是必要非充分，后端鲁棒化兜底]
\cref{eq:nlopt-geomconstraint,eq:nlopt-pcm} 的几何/一致性约束，即便在真值 $\mathbf{x}$ 处也只是内点的\emph{必要}条件（外点偶尔也能碰巧满足）；加之共识最大化与 PCM 多用近似算法、不保证选出全部内点——前端剔除后\emph{仍可能残留外点}。故工程上前端剔除后\emph{必接}后端鲁棒化（下文），二者一硬一软、互补。这也呼应 \cref{eq:nlopt-ransac} 旁的结论：RANSAC \emph{剔除}外点、鲁棒核\emph{降权}外点。
\end{insight}

% ----------------------------------------------------------------
\subsection{后端鲁棒化：M-估计与鲁棒核族}
\label{sec:nlopt-robust-mest}

\paragraph{M-估计与鲁棒核族。}\cite{barfoot2024state,gaoxiang2019slam14,carlone2026handbook}
M-估计（“极大似然型”估计）把目标改为
\begin{equation}
  J'(\mathbf{x})=\sum_i\alpha_i\,\rho\big(u_i(\mathbf{x})\big),\qquad u_i(\mathbf{x})=\sqrt{\mathbf{e}_i(\mathbf{x})^\top\boldsymbol{\Sigma}_i^{-1}\mathbf{e}_i(\mathbf{x})},
  \label{eq:nlopt-mestim}
\end{equation}
$\alpha_i>0$，$u_i$ 为标量马氏范数，$\rho$ 须有界增长慢、$u=0$ 唯一零点、$u>0$ 单调增。常见核（影响函数 $\psi=\rho'$ 决定大残差的“话语权”）：

\begin{table}[H]\centering\small
\caption{常见鲁棒核全家福（标量马氏范数 $u$；影响函数 $\psi=\rho'$ 决定大残差的话语权，IRLS 权 $w(u)\propto\psi(u)/u$，见 \cref{eq:irls}）}
\label{tab:robust}
\begin{tabular}{p{2.4cm}p{4.4cm}p{6.2cm}}
\toprule
核 & $\rho(u)$ & 大残差行为（$\psi=\rho'$、权 $w\propto\psi/u$） \\
\midrule
二范数（L2，非鲁棒） & $\tfrac12 u^2$ & $\psi=u$ 无界、$w=1$ 恒定，外点主导 \\
Huber & $\tfrac12u^2\ (|u|\le\delta)$；$\delta(|u|-\tfrac12\delta)\ (|u|>\delta)$ & $\psi=\delta\,\mathrm{sgn}(u)$、$w=\delta/|u|$（外点权 $\propto1/|u|$ 衰减），内点二次外点线性、\emph{凸} \\
Cauchy（Lorentzian） & $\tfrac{c^2}{2}\ln(1+u^2/c^2)$ & $\psi=\dfrac{u}{1+u^2/c^2}\to0$、$w=\dfrac{1}{1+u^2/c^2}$，强抑制但\emph{永不为零}，非凸 \\
Geman--McClure & $\tfrac12\dfrac{u^2}{1+u^2}$ & $\psi=\dfrac{u}{(1+u^2)^2}$、$w=\dfrac{1}{(1+u^2)^2}$，$\rho$ 有界（饱和），重外点近乎零权，非凸 \\
截断二次（truncated quadratic） & $\min\{u^2,\beta^2\}$：$u^2\ (|u|\le\beta)$；$\beta^2\ (|u|>\beta)$ & $|u|\le\beta$ 时 $\psi=2u,w=1$（内点照常）；$|u|>\beta$ 时 $\psi=w=0$（硬拒绝、代价封顶 $\beta^2$），非凸；阈值 $\beta$ 已知时优选（\cref{eq:nlopt-truncquad}） \\
Tukey（biweight/ bisquare） & $\dfrac{c^2}{6}\Big[1-\big(1-(u/c)^2\big)^3\Big]\ (|u|\le c)$；\ $\dfrac{c^2}{6}\ (|u|>c)$ & $\psi=u\big(1-(u/c)^2\big)^2$、$w=\big(1-(u/c)^2\big)^2\ (|u|\le c)$，\textbf{$|u|>c$ 时 $\psi=w=0$（硬拒绝）}，非凸 \\
最大共识（max-consensus） & 大残差为常数（计 $1$）、小残差为 $0$ & “数外点”：最小化即令外点最少 $\Leftrightarrow$ 共识最大化 \cref{eq:nlopt-consensusmax}（RANSAC 采样优化的就是此损失）\\
\bottomrule
\end{tabular}
\end{table}

\noindent Huber 是“内点保二次、外点转线性”的折中（凸、收敛好）；Cauchy/Geman--McClure 对外点压制更强但\emph{非凸}、易陷局部极小\cite{barfoot2024state}。

\paragraph{Tukey biweight：唯一“硬拒绝”外点的核。}
上表末行的 Tukey（双权/biweight，亦称 bisquare）值得单列：它与前几款“软降权”核的本质差别在\emph{权函数在 $|u|>c$ 处\textbf{严格归零}}——残差超过阈值 $c$ 的观测被\emph{彻底剔除}（$\psi=w=0$），对解不再有任何话语权（梯度贡献为零）；而 Cauchy/Geman--McClure 的权只\emph{渐近}趋零、永远保留一丝影响。这把 Tukey 摆在“M-估计软降权”与“RANSAC 硬剔除”（\cref{eq:nlopt-ransac}）之间：它\emph{在优化迭代内部}做硬剔除，可看作\emph{可微化、自适应阈值}的 RANSAC。代价是 $\rho$ 在 $|u|>c$ 处\emph{完全平坦}（$\rho\equiv c^2/6$ 饱和），故\emph{强非凸}、对初值与外点比例敏感，且一旦某内点被误判出阈就再难“拉回”——实践中常先用 Huber/Cauchy 暖启、或配渐进非凸 GNC（本节后文 \cref{eq:nlopt-robust-cauchy} 之 insight）逐步收紧 $c$。调参 $c$ 即“判为外点的残差阈值（按 $\sigma$ 计）”：统计上取 $c=4.685$ 可在标准高斯噪声下达 $95\%$ 渐近效率（即相对克拉默--拉奥下界只损失 $5\%$，\cref{sec:nlopt-cov}）\cite{zhang1997parameter}。其影响函数\emph{先升后\textbf{降回零}}（redescending，重下降型），是“重下降 M-估计”的典型代表，比影响函数单调饱和的核压制外点更彻底。
（\textbf{形式由来}：把 $\psi(u)=u(1-(u/c)^2)^2$ 对 $u$ 积分即得上表 $\rho$；其 IRLS 权 $w(u)=\psi(u)/u=(1-(u/c)^2)^2$ 恰是 \cref{eq:irls} 中 $\tfrac{\alpha_i}{u_i}\partial\rho/\partial u_i$ 给出的因子，故 Tukey 与其他核共用同一套 IRLS 机器，只是权在阈外被钳为零。）

\paragraph{截断二次核：硬阈值的“计外点”损失。}\cite{carlone2026handbook}
表中另一个硬截断核是\emph{截断二次}（truncated quadratic）损失
\begin{equation}
  \rho(u)=\min\{u^2,\beta^2\}=
  \begin{cases}u^2,&|u|\le\beta\ (\text{内点}),\\ \beta^2,&|u|>\beta\ (\text{外点}),\end{cases}
  \label{eq:nlopt-truncquad}
\end{equation}
$\beta$ 是\emph{内点最大容许残差}：残差在阈内照常按二次计，超阈则代价\emph{封顶} $\beta^2$、梯度归零（与 Tukey 同为硬拒绝、重下降）。它与 Cauchy/G-M 的“软降权”相反，在阈值已知（即知道内点的最大期望误差）时优选；代价是阈处\emph{突变}使其强非凸、需专用解法（下文 GNC）。截断二次还把后端与前端的两条线\emph{接上}：把它推到极端、对所有大残差只罚常数 $1$、小残差罚 $0$，就是表末的\emph{最大共识损失}——最小化它即令外点数最少，恰是 \cref{eq:nlopt-consensusmax} 的共识最大化（RANSAC 采样优化的正是这个损失）。这条“鲁棒核 $\leftrightarrow$ 共识最大化”的桥与 \cref{sec:nlopt-robust-frontend} 闭环。它也可由 MAP 导出：设噪声服从“高斯（内点）$+$ 均匀（外点）”的最大混合分布，其负对数即截断二次\cite{carlone2026handbook}——这与下文“鲁棒 $=$ 协方差先验 MAP”同一精神。

一条更完整的核列表见 Zhang (1997)\cite{zhang1997parameter}；比较研究见 MacTavish--Barfoot (2015)。

\paragraph{IRLS：把鲁棒核变回加权最小二乘（完整推导）。}\cite{barfoot2024state,carlone2026handbook}
鲁棒核如何进入正规方程？答案是\emph{迭代重加权最小二乘}（IRLS）。对 \cref{eq:nlopt-mestim} 用链式法则求梯度：
\begin{equation}
  \frac{\partial J'}{\partial\mathbf{x}}=\sum_i\alpha_i\frac{\partial\rho}{\partial u_i}\frac{\partial u_i}{\partial\mathbf{e}_i}\frac{\partial\mathbf{e}_i}{\partial\mathbf{x}},
  \qquad \frac{\partial u_i}{\partial\mathbf{e}_i}=\frac{1}{u_i}\mathbf{e}_i^\top\boldsymbol{\Sigma}_i^{-1}.
  \label{eq:nlopt-irls-grad}
\end{equation}
代入得 $\partial J'/\partial\mathbf{x}=\sum_i\mathbf{e}_i^\top\mathbf{Y}_i^{-1}\,\partial\mathbf{e}_i/\partial\mathbf{x}$，其中\emph{依赖 $\mathbf{x}$ 的新（逆）协方差}
\begin{equation}
  \mathbf{Y}_i(\mathbf{x})^{-1}=\frac{\alpha_i}{u_i(\mathbf{x})}\frac{\partial\rho}{\partial u_i}\Big|_{u_i(\mathbf{x})}\boldsymbol{\Sigma}_i^{-1}.
  \label{eq:irls}
\end{equation}
此梯度与普通加权最小二乘 $\partial(\tfrac12\sum\mathbf{e}_i^\top\mathbf{Y}_i^{-1}\mathbf{e}_i)/\partial\mathbf{x}$ \emph{形式完全相同}，只是 $\boldsymbol{\Sigma}_i\to\mathbf{Y}_i$。因 $\mathbf{e}_i$ 本就非线性、要迭代，故在上一步工作点 $\mathbf{x}_{\mathrm{op}}$ 处\emph{冻结}权 $\mathbf{Y}_i(\mathbf{x}_{\mathrm{op}})$，每次迭代解一个普通加权最小二乘 $J''(\mathbf{x})=\tfrac12\sum_i\mathbf{e}_i^\top\mathbf{Y}_i(\mathbf{x}_{\mathrm{op}})^{-1}\mathbf{e}_i$，再更新 $\mathbf{x}_{\mathrm{op}}$、重算权。\textbf{合法性}：收敛时 $\hat{\mathbf{x}}=\mathbf{x}_{\mathrm{op}}$，故 $\partial J'/\partial\mathbf{x}|_{\hat{\mathbf{x}}}=\partial J''/\partial\mathbf{x}|_{\hat{\mathbf{x}}}=\mathbf{0}$——\emph{两者同极小、同解}，只是到达路径不同。在 Ceres/g2o 里，给残差块设一个 \texttt{LossFunction}/\texttt{RobustKernel} 即自动启用此机制。

\paragraph{例：Cauchy 核的 IRLS 权。}\cite{barfoot2024state}
对 Cauchy $\rho=\tfrac12\ln(1+u^2)$，$\partial\rho/\partial u=u/(1+u^2)$，代入 \cref{eq:irls}：
\begin{equation}
  \mathbf{Y}_i(\mathbf{x}_{\mathrm{op}})^{-1}=\frac{\alpha_i}{u_i}\cdot\frac{u_i}{1+u_i^2}\boldsymbol{\Sigma}_i^{-1}
  \ \Rightarrow\ \mathbf{Y}_i(\mathbf{x}_{\mathrm{op}})=\frac{1}{\alpha_i}\big(1+\mathbf{e}_i^\top\boldsymbol{\Sigma}_i^{-1}\mathbf{e}_i\big)\boldsymbol{\Sigma}_i,
  \label{eq:nlopt-cauchy-Y}
\end{equation}
即原协方差 $\boldsymbol{\Sigma}_i$ 的\emph{膨胀版}——随二次误差增大而变大，对外点自动\emph{少信任}。这与下文 MAP 推导的最优协方差（\cref{eq:nlopt-invwishart-M}）形式呼应。

\paragraph{（选读）鲁棒 $=$ 逆 Wishart 先验下的 MAP（完整证明）。}\cite{barfoot2024state}
鲁棒核不是 ad-hoc 补丁，而能从 MAP 严格导出。把每个因子的协方差 $\mathbf{M}_i$ 也当\emph{待估隐变量}，并给它一个\emph{逆 Wishart 先验} $\mathbf{M}_i\sim\mathcal{W}^{-1}(\boldsymbol{\Psi}_i,\nu_i)$（定义在正定矩阵上，$\boldsymbol{\Psi}_i$ 为 scale 矩阵、$\nu_i$ 为自由度，$M_i=\dim\mathbf{M}_i$）：
\begin{equation}
  p(\mathbf{M}_i)\propto\det(\mathbf{M}_i)^{-\frac{\nu_i+M_i+1}{2}}\exp\!\Big(-\tfrac12\operatorname{tr}(\boldsymbol{\Psi}_i\mathbf{M}_i^{-1})\Big).
  \label{eq:nlopt-invwishart}
\end{equation}
联合 MAP $\{\hat{\mathbf{x}},\hat{\mathbf{M}}\}=\arg\min J'(\mathbf{x},\mathbf{M})$，$J'=-\ln p(\mathbf{x},\mathbf{M}\mid\mathbf{z})$。后验因式分解 $p(\mathbf{x},\mathbf{M}\mid\mathbf{z})=\prod_i p(\mathbf{x}\mid\mathbf{z}_i,\mathbf{M}_i)p(\mathbf{M}_i)$，代入高斯似然与 \cref{eq:nlopt-invwishart}（丢与 $\mathbf{x},\mathbf{M}$ 无关项）得
\begin{equation}
  J'(\mathbf{x},\mathbf{M})=\tfrac12\sum_i\Big(\mathbf{e}_i^\top\mathbf{M}_i^{-1}\mathbf{e}_i-\alpha_i\ln\det(\mathbf{M}_i^{-1})+\operatorname{tr}(\boldsymbol{\Psi}_i\mathbf{M}_i^{-1})\Big),\quad \alpha_i=\nu_i+M_i+2.
  \label{eq:nlopt-invwishart-J}
\end{equation}
\emph{先消去 $\mathbf{M}_i$}：对 $\mathbf{M}_i^{-1}$ 求导（用 $\partial(\mathbf{e}^\top\mathbf{M}^{-1}\mathbf{e})/\partial\mathbf{M}^{-1}=\mathbf{e}\mathbf{e}^\top$、$\partial\ln\det(\mathbf{M}^{-1})/\partial\mathbf{M}^{-1}=\mathbf{M}$、$\partial\operatorname{tr}(\boldsymbol{\Psi}\mathbf{M}^{-1})/\partial\mathbf{M}^{-1}=\boldsymbol{\Psi}$）置零：
\begin{equation}
  \tfrac12\mathbf{e}_i\mathbf{e}_i^\top-\tfrac12\alpha_i\mathbf{M}_i+\tfrac12\boldsymbol{\Psi}_i=\mathbf{0}
  \ \Rightarrow\ \mathbf{M}_i(\mathbf{x})=\underbrace{\tfrac1{\alpha_i}\boldsymbol{\Psi}_i}_{\text{常数}}+\underbrace{\tfrac1{\alpha_i}\mathbf{e}_i(\mathbf{x})\mathbf{e}_i(\mathbf{x})^\top}_{\text{膨胀}}.
  \label{eq:nlopt-invwishart-M}
\end{equation}
最优协方差在残差大处从常数被\emph{膨胀}（与 IRLS 的 \cref{eq:nlopt-cauchy-Y} 神似）。把 $\mathbf{M}_i(\mathbf{x})$ 代回 \cref{eq:nlopt-invwishart-J}（丢与 $\mathbf{x}$ 无关因子）得
\begin{equation}
  J'(\mathbf{x})=\tfrac12\sum_i\alpha_i\ln\!\big(1+\mathbf{e}_i(\mathbf{x})^\top\boldsymbol{\Psi}_i^{-1}\mathbf{e}_i(\mathbf{x})\big),
  \label{eq:nlopt-robust-cauchy}
\end{equation}
当 scale 取通常协方差 $\boldsymbol{\Psi}_i=\boldsymbol{\Sigma}_i$ 时，\textbf{这正是加权 Cauchy 鲁棒核}（\cref{tab:robust}）。

\begin{insight}[鲁棒核 $=$ 某协方差先验下的 MAP]
\cref{eq:nlopt-robust-cauchy} 把“鲁棒估计”从工程补丁提升为\emph{有概率依据的 MAP}：大残差对应“这条观测的真实噪声其实更大”，故自动降权——而\emph{降多少}由逆 Wishart 先验定量给出。Barfoot\cite{barfoot2024state} 推测\emph{所有}鲁棒核都源自对协方差先验 $p(\mathbf{M})$ 的某种选择（Black--Rangarajan 1996 是起点）。现代做法还有：自适应单参数核族（Barron 2019，估计中\emph{自选}最佳核）、渐进非凸 GNC（逐步施加非凸核以避开劣质局部极小，下文）。
\end{insight}

\paragraph{（选读）一般化：Black--Rangarajan 对偶与外点过程。}\cite{black1996rangarajan,carlone2026handbook}
上面“鲁棒 $=$ 协方差先验 MAP”的推导给的是 Cauchy 这一\emph{特例}（取逆 Wishart scale $\boldsymbol{\Psi}_i=\boldsymbol{\Sigma}_i$）。它其实是一个更一般定理——\emph{Black--Rangarajan（B-R）对偶}——的一个实例。B-R 对偶给 IRLS 的权更新一个比 \cref{eq:irls} 启发式推导更严谨的根据，也修好了 \cref{eq:irls} 在 $\rho$ 不可导处（如截断二次的阈点）失定义的毛病。先看截断二次 \cref{eq:nlopt-truncquad} 的直觉：引入辅助权 $w_i\in[0,1]$，可把 $\min\{r_i^2,\beta_i^2\}$ 等价写成 $\min_{w_i\in[0,1]}w_ir_i^2+(1-w_i)\beta_i^2$（$r_i^2\le\beta_i^2$ 时最优 $w_i{=}1$ 取 $r_i^2$、否则 $w_i{=}0$ 取 $\beta_i^2$），第一项是加权最小二乘、第二项只惩罚权而\emph{不含} $\mathbf{x}$。一般地：

\begin{theorem}[Black--Rangarajan 对偶]\label{thm:nlopt-br}
给鲁棒核 $\rho(\cdot)$，记 $\phi(z):=\rho(\sqrt{z})$。若 $\phi$ 满足 $\lim_{z\to0}\phi'(z)=c>0$（$c$ 为归一化常数，取 $1$ 或与本书 $\tfrac12$-loss 一致的 $\tfrac12$ 均可，不影响结论）、$\lim_{z\to\infty}\phi'(z)=0$、$\phi''(z)<0$，则 M-估计 \cref{eq:nlopt-mestim} 等价于联合最小化
\begin{equation}
  \min_{\mathbf{x},\,w_i\in[0,1]}\ \sum_i\Big[\,w_i\,r_i^2(\mathbf{x})+\Phi_\rho(w_i)\,\Big],
  \label{eq:nlopt-br}
\end{equation}
其中 $r_i=u_i$ 为标量马氏范数，$w_i$ 为每条残差的权变量，$\Phi_\rho(w_i)$ 称\emph{外点过程}（outlier process），其形式由 $\rho$ 决定。
\end{theorem}

\noindent 外点过程把“核选得对不对”编码进对权的惩罚。两个实例：截断二次给 $\Phi_\rho(w_i)=(1-w_i)\beta_i^2$；Geman--McClure（$\rho=\beta_i^2 r_i^2/(\beta_i^2+r_i^2)$）给 $\Phi_\rho(w_i)=\beta_i^2(\sqrt{w_i}-1)^2$（验证：对 \cref{eq:nlopt-br} 的 $w_i$ 求导置零得最优权 $w_i^*=\big(\beta_i^2/(r_i^2+\beta_i^2)\big)^2$，代回即还原 G-M 核）。

有了 \cref{thm:nlopt-br}，IRLS 自然现身为\emph{交替最小化}：\cref{eq:nlopt-br} 难对 $\mathbf{x},w_i$ \emph{联合}优化，但固定一方优另一方都简单——固定 $w_i$ 是加权最小二乘、固定 $\mathbf{x}$ 则按因子拆成 $N$ 个标量子问题 $\min_{w_i\in[0,1]}\Phi_\rho(w_i)+w_ir_i^2$（常有闭式）。于是第 $t$ 步交替：
\begin{equation}
  \textbf{变量更新：}\ \mathbf{x}^{(t)}\in\arg\min_{\mathbf{x}}\sum_i w_i^{(t)}r_i^2(\mathbf{x});
  \qquad
  \textbf{权更新：}\ w_i^{(t+1)}\in\arg\min_{w_i\in[0,1]}\Phi_\rho(w_i)+w_i\,r_i^2(\mathbf{x}^{(t)}).
  \label{eq:nlopt-br-altmin}
\end{equation}
由 B-R 对偶得到的权更新恰与 \cref{eq:irls} 那条 $w_i\propto\psi(r_i)/r_i$ 吻合，但定义更干净（阈点也良定义）。\textbf{一个有趣落点}：在 SLAM 里用 G-M 核跑这套 IRLS，就得到鲁棒 SLAM 中的\emph{动态协方差缩放}（dynamic covariance scaling）算法——可见它并非另起炉灶，而是 B-R 对偶 $+$ G-M 核的具体化（与 \cref{sec:nlopt-robust-mest} 的“开关约束 \cite{sunderhauf2012switchable}”同族）。

% ----------------------------------------------------------------
\subsection{渐进非凸（GNC）：让 IRLS 不再怕坏初值}
\label{sec:nlopt-robust-gnc}

\cref{eq:nlopt-br-altmin} 的 IRLS 因常见鲁棒核\emph{非凸}，其收敛\emph{高度依赖初值}——$\mathbf{x}^{(0)}$ 离最优远、或初始权 $w_i^{(0)}$ 没反映内外点归属，就易陷劣质局部极小（截断二次/G-M 在仅 $10\%$ 外点时就可能失败\cite{carlone2026handbook}）。\emph{渐进非凸}（graduated non-convexity，GNC）\cite{yang2020gnc} 让 IRLS 大大减弱对初值的敏感：给核 $\rho$ 造一个带标量\emph{光滑因子} $\mu$ 的光滑族 $\rho_\mu$，$\mu$ 控制非凸程度——一端 $\rho_\mu$ \emph{凸}、另一端还原原核 $\rho$。两例：
\begin{align}
  \text{GNC 截断二次：}\ &\rho_\mu(r)=
  \begin{cases}
  r^2,& r^2\in[0,\tfrac{\mu}{\mu+1}\beta^2],\\
  2\beta|r|\sqrt{\mu(\mu+1)}-\mu(\beta^2+r^2),& r^2\in[\tfrac{\mu}{\mu+1}\beta^2,\tfrac{\mu+1}{\mu}\beta^2],\\
  \beta^2,& r^2\in[\tfrac{\mu+1}{\mu}\beta^2,+\infty),
  \end{cases}
  \label{eq:nlopt-gnc-tq}\\
  \text{GNC Geman--McClure：}\ &\rho_\mu(r)=\frac{\mu\beta^2 r^2}{\mu\beta^2+r^2}.
  \label{eq:nlopt-gnc-gm}
\end{align}
\cref{eq:nlopt-gnc-tq} 在 $\mu\to0$ 时凸、$\mu\to\infty$ 时还原截断二次；\cref{eq:nlopt-gnc-gm} 在 $\mu\to\infty$ 时凸、$\mu{=}1$ 时还原 G-M。GNC 算法在 \cref{eq:nlopt-br-altmin} 两步上再加一步\emph{控制参数更新}：
\begin{algo}[渐进非凸（GNC）]\label{alg:nlopt-gnc}
给光滑核 $\rho_\mu$、初值 $\mathbf{x}^{(0)}$、初始 $\mu$（取使 $\rho_\mu$ 凸的一端）。重复：
(1) \textbf{变量更新}：以当前权解加权最小二乘 $\mathbf{x}^{(t)}\in\arg\min_{\mathbf{x}}\sum_i w_i^{(t)}r_i^2(\mathbf{x})$；
(2) \textbf{权更新}：$w_i^{(t+1)}\in\arg\min_{w_i\in[0,1]}\Phi_{\rho_\mu}(w_i)+w_i r_i^2(\mathbf{x}^{(t)})$；
(3) \textbf{控制参数更新}：调 $\mu$ 朝原核方向加大非凸性（截断二次 $\mu\leftarrow\gamma\mu$、G-M $\mu\leftarrow\mu/\gamma$，$\gamma>1$）。
直至 $\mu$ 还原原核且收敛。
\end{algo}
\noindent 关键：\textbf{B-R 对偶对光滑核 $\rho_\mu$ 仍成立}，故 GNC 与 IRLS \cref{eq:nlopt-br-altmin} 共用同一台“交替最小化”机器，只是从凸的 $\rho_\mu$ 起步、逐步收紧到原核——“先解一个好解的凸问题、再缓慢逼近真问题”，这正是 \cref{sec:nlopt-why} pitfall“初值决定成败”在鲁棒估计里的系统化对策。GNC 经验上对位姿图能耐受 $80\text{--}90\%$ 错误回环，但\emph{不}提供收敛保证、且性能因问题而异（位姿图极佳、别处可能退化）。

\paragraph{工程直觉：何时够、何时崩。}\cite{carlone2026handbook}
把上述方法在 M3500、SubT、Victoria Park 上对比可得若干诚实结论\cite{carlone2026handbook}：(i) 有可靠\emph{里程计骨架}（odometry backbone，连续帧视角相近、数据关联简单，可作可信约束与初值）时，鲁棒核 $+$ IRLS 已能在 M3500/SubT 上夺回鲁棒性，但朴素梯度下降在强非凸的截断二次上仍可能在 Victoria Park 失败；GNC 收敛更好、三个数据集都拿到较准轨迹。(ii) 一旦\emph{连里程计也含外点}（轮速打滑、激光误配、或多机 SLAM 无统一骨架），问题极难：此时 GNC 因坏初值会把\emph{所有}观测（含内点）误判成外点而崩；\emph{前端} PCM 不依赖初值，反而能在 SubT/Victoria Park 上恢复，而\emph{PCM 前端 $+$ GNC 后端}组合最稳——再次印证“前端硬剔除 $+$ 后端软降权”的互补。(iii) 即便如此，三法在 M3500 上仍可能全崩，说明无骨架的鲁棒 SLAM 仍是开放难题。这也引出下一节的\emph{可证最优}求解器（\cref{sec:nlopt-certifiable}）：用半定松弛把鲁棒估计松弛成凸问题，附带最优性证书。

\begin{pitfall}[概念误区：鲁棒核是“工程黑魔法”]
“加个 Huber 就行”常被当成玄学。其实鲁棒核 $=$ IRLS 重加权（\cref{eq:irls}）$=$ 逆 Wishart 先验下的 MAP（\cref{eq:nlopt-robust-cauchy}）——三者等价、有严格概率根。误区后果：乱设核参数 $c$/$\delta$ 而不理解其“残差多大才算外点”的物理含义。正确做法：核参数对应“开始降权的残差阈值”（约几个 $\sigma$），按噪声尺度设。另注意非凸核（Cauchy/GM）需好初值或 GNC，否则陷局部极小。
\end{pitfall}

\begin{practice}
\begin{enumerate}
  \item 画出 L2、Huber、Cauchy 的 $\rho(u)$ 与影响函数 $\psi(u)=\rho'(u)$，比较大 $u$ 处的话语权。
  \item 对 Cauchy 核推出 IRLS 权 \cref{eq:nlopt-cauchy-Y} 的具体形式。
  \item （推导，草稿纸）完整重做 \cref{eq:nlopt-invwishart-J}$\to$\cref{eq:nlopt-invwishart-M}$\to$\cref{eq:nlopt-robust-cauchy} 的“消去 $\mathbf{M}_i$ 得 Cauchy 核”。
  \item （概念）为什么 Huber 比 Cauchy 收敛更稳（凸），但抑制外点更弱？用 \cref{eq:nlopt-ransac} 估算 $w=0.5,n=2,p=0.99$ 需多少 RANSAC 迭代。
  \item （概念）位姿图回环为何更适合 PCM（\cref{eq:nlopt-pcm}）而非 RANSAC？由“最小集须含生成树（$\ge N-1$ 条边）”与 \cref{eq:nlopt-ransac} 的 $1/w^n$ 论证之。
  \item （推导）对截断二次 \cref{eq:nlopt-truncquad}，由 B-R 对偶 \cref{eq:nlopt-br} 验证外点过程 $\Phi_\rho(w_i)=(1-w_i)\beta^2$（对 $w_i$ 求最优给出 $w_i^*\in\{0,1\}$ 的硬判别）。对 G-M 核重做，得 $\Phi_\rho=\beta^2(\sqrt{w_i}-1)^2$ 与 \cref{eq:nlopt-br-altmin} 的权。
  \item （概念）GNC（\cref{alg:nlopt-gnc}）为什么能减弱对初值的敏感？联系 \cref{sec:nlopt-why} 的“初值决定成败”与“先解凸问题再逼近”。并说明“连里程计含外点时 GNC 反而崩、PCM$+$GNC 最稳”的原因。
\end{enumerate}
\end{practice}

% ════════════════════════════════════════════════════════════════
\section{收敛之后：协方差、偏差与一致性}
\label{sec:nlopt-cov}

\paragraph{拉普拉斯近似：逆海森即协方差。}\cite{barfoot2024state}
MAP/GN 只给点估计 $\hat{\mathbf{x}}$，但常需其不确定度。\emph{拉普拉斯近似}用“以 $\hat{\mathbf{x}}$ 为中心、逆海森为协方差”的高斯近似后验：MAP 收敛后令 $\hat{\mathbf{x}}=\mathbf{x}_{\mathrm{op}}$，而 $\hat{\boldsymbol{\Sigma}}$ 取为代价函数海森在 $\hat{\mathbf{x}}$ 处的逆。GN 情形下，正规方程左端\emph{恰是逆协方差}：
\begin{equation}
  \hat{\boldsymbol{\Sigma}}^{-1}=\mathbf{J}^\top\mathbf{W}^{-1}\mathbf{J}=\boldsymbol{\Lambda}\quad(\text{在 }\hat{\mathbf{x}}\text{ 处求值}).
  \label{eq:laplace}
\end{equation}
即\textbf{信息矩阵 $=$ 逆协方差}（Barfoot\cite{barfoot2024state} §4.3.1 把 GN 左端直接取作 Laplace 逆协方差 $\hat{\mathbf{P}}^{-1}=\mathbf{H}^\top\mathbf{W}^{-1}\mathbf{H}$）。这与 \cref{ch:slam_est}、\cref{ch:eskf} 一脉相承，也和变量消元的 $\boldsymbol{\Sigma}_j=(\mathbf{R}_j^\top\mathbf{R}_j)^{-1}$（\cref{eq:nlopt-elim-cond}）一致——回代 $\boldsymbol{\Lambda}^{-1}=\mathbf{R}^{-1}\mathbf{R}^{-\top}$ 即得后验协方差。需要某变量的\emph{边缘}协方差时，回代 $\boldsymbol{\Lambda}^{-1}$ 的对应块（勿显式整体求逆）。

\paragraph{与 CRLB 的关系。}
\cref{eq:laplace} 的信息矩阵 $\boldsymbol{\Lambda}=\mathbf{J}^\top\mathbf{W}^{-1}\mathbf{J}$ 正是高斯观测下的\emph{费舍尔信息矩阵}（Fisher information），其逆即\emph{克拉默--拉奥下界}（CRLB）——任何无偏估计的协方差不可能小于 $\boldsymbol{\Lambda}^{-1}$。\footnote{“CRLB $=$ 费舍尔信息之逆”是经典统计结论：对数似然的二阶导期望（费舍尔信息）之逆给出无偏估计协方差的下界，且极大似然在大样本下渐近达到该界（统计高效）。高斯观测下负对数似然恰为 $\tfrac12\mathbf{e}^\top\mathbf{W}^{-1}\mathbf{e}$，其费舍尔信息即 $\boldsymbol{\Lambda}=\mathbf{J}^\top\mathbf{W}^{-1}\mathbf{J}$（GN 近似海森），故此处三者重合。本书 \cref{ch:slam_est} 已就线性高斯情形给出同一等号；非线性时它只在 $\hat{\mathbf{x}}$ 邻域近似成立（Laplace 近似），且 ML 偏差（\cref{eq:nlopt-mlbias}）使“无偏”前提本身受损。}换言之，GN 收敛后的逆海森\emph{同时}是“后验协方差（Laplace）”“费舍尔信息之逆”“CRLB”——三位一体，这是“最小二乘解在弱非线性下\emph{统计高效}”的根。

\paragraph{（选读）最大似然/MAP 有偏（Box 1971 闭式）。}\cite{barfoot2024state}
非线性观测下，ML/MAP \emph{有偏}（除非模型线性）——因后验众数 $\neq$ 均值（前述 pitfall）。Box (1971) 给近似偏差闭式：设 ML $\hat{\mathbf{x}}=\arg\min\tfrac12\sum_k(\mathbf{y}_k-\mathbf{g}_k(\mathbf{x}))^\top\boldsymbol{\Sigma}_{v,k}^{-1}(\mathbf{y}_k-\mathbf{g}_k(\mathbf{x}))$，记 $\mathbf{G}_k=\partial\mathbf{g}_k/\partial\mathbf{x}$、$\boldsymbol{\mathcal{G}}_{jk}=\partial^2 g_{jk}/\partial\mathbf{x}\partial\mathbf{x}^\top$、$\mathbf{W}(\mathbf{x})=\sum_k\mathbf{G}_k^\top\boldsymbol{\Sigma}_{v,k}^{-1}\mathbf{G}_k$，$\mathbf{1}_j$ 为单位阵第 $j$ 列，则系统偏差
\begin{equation}
  \mathbb{E}[\delta\mathbf{x}]=-\tfrac12\mathbf{W}(\mathbf{x})^{-1}\sum_k\mathbf{G}_k^\top\boldsymbol{\Sigma}_{v,k}^{-1}\sum_j\mathbf{1}_j\operatorname{tr}\!\big(\boldsymbol{\mathcal{G}}_{jk}\,\mathbf{W}(\mathbf{x})^{-1}\big).
  \label{eq:nlopt-mlbias}
\end{equation}
\emph{推导骨架}（Barfoot\cite{barfoot2024state} §4.3.3，完整步在 \cref{sec:nlopt-appendix}）：把最优条件 $\sum_k\mathbf{G}_k(\hat{\mathbf{x}})^\top\boldsymbol{\Sigma}_{v,k}^{-1}(\mathbf{y}_k-\mathbf{g}_k(\hat{\mathbf{x}}))=\mathbf{0}$ 在真值处二阶展开，设 $\delta\mathbf{x}=\mathbf{A}\mathbf{n}+\mathbf{b}(\mathbf{n})$（线性 $+$ 二次依赖噪声），用奇偶对消（$\mathbf{n}\to-\mathbf{n}$）定出线性系数 $\mathbf{A}=\mathbf{W}^{-1}\sum_k\mathbf{G}_k^\top\boldsymbol{\Sigma}_{v,k}^{-1}\mathbf{P}_k$ 使线性项为零，再对二次项取期望（用 $\operatorname{tr}$ 恒等式 $\mathbb{E}[\mathbf{n}^\top\mathbf{A}^\top\boldsymbol{\mathcal{G}}\mathbf{A}\mathbf{n}]=\operatorname{tr}(\boldsymbol{\mathcal{G}}\mathbf{A}\boldsymbol{\Sigma}\mathbf{A}^\top)$）即得 \cref{eq:nlopt-mlbias}。偏差正比于观测模型的\emph{曲率}（二阶导 $\boldsymbol{\mathcal{G}}$）乘协方差量级；弱非线性时可近似\emph{去偏} $\hat{\mathbf{x}}\leftarrow\hat{\mathbf{x}}-\mathbb{E}[\delta\mathbf{x}]$。强非线性时这是 GN/EKF 系统误差之源（\cref{ch:eskf} 的“EKF 有偏不一致”同源）。

\paragraph{一致性检验：NEES 与 NIS。}\cite{barfoot2024state}
估出协方差后，怎么知道它\emph{可信}（不过自信/欠自信）？两个统计量。
\emph{NEES}（归一化估计误差平方，需真值 $\mathbf{x}_{\text{true}}$）：
\begin{equation}
  \epsilon_{\text{nees},k}=\hat{\mathbf{e}}_k^\top\hat{\boldsymbol{\Sigma}}_k^{-1}\hat{\mathbf{e}}_k,\quad \hat{\mathbf{e}}_k=\hat{\mathbf{x}}_k-\mathbf{x}_{\text{true},k},\qquad \mathbb{E}[\epsilon_{\text{nees},k}]=N,
  \label{eq:nlopt-nees}
\end{equation}
$N$ 为状态维。$<N$ 为\emph{欠自信}、$>N$ 为\emph{过自信}。\emph{NIS}（归一化创新平方，无需真值）：
\begin{equation}
  \epsilon_{\text{nis},k}=\mathbf{e}_{y,k}^\top\mathbf{S}_{y,k}^{-1}\mathbf{e}_{y,k},\quad \mathbf{e}_{y,k}=\mathbf{y}_k-\mathbf{g}(\check{\mathbf{x}}_k),\quad \mathbf{S}_{y,k}=\mathbf{G}_k\check{\boldsymbol{\Sigma}}_k\mathbf{G}_k^\top+\boldsymbol{\Sigma}_{v,k},\qquad \mathbb{E}[\epsilon_{\text{nis},k}]=M,
  \label{eq:nlopt-nis}
\end{equation}
$M$ 为观测维。两者在多次试验或（遍历假设下）多时刻上应分别近 $N,M$；统计检验用卡方分位数 $Q_{\chi^2(NK)}(\ell)\le\sum_k\epsilon_{\text{nees},k}\le Q_{\chi^2(NK)}(u)$（NIS 同理）。\textbf{警示}：观测/运动非线性时（线性化推导），NEES/NIS 永不\emph{精确}满足，只作近似诊断\cite{barfoot2024state}。\cref{sec:nlopt-appendix} 给“自适应协方差估计与 NIS 的联系”。

\begin{pitfall}[概念误区：把 GN 的逆海森协方差当“真协方差”]
$\hat{\boldsymbol{\Sigma}}=\boldsymbol{\Lambda}^{-1}$ 只是\emph{Laplace 近似}（把后验当以众数为心的高斯），且 $\boldsymbol{\Lambda}=\mathbf{J}^\top\mathbf{W}^{-1}\mathbf{J}$ 又是 GN 近似海森。两重近似 $+$ ML 偏差（\cref{eq:nlopt-mlbias}）使强非线性下它\emph{过乐观}。后果：下游（数据关联门限、主动 SLAM、一致性）误判。对策：用 NEES/NIS 检验一致性；强非线性处考虑去偏或采样法（\cref{ch:eskf}）。
\end{pitfall}

\begin{practice}
\begin{enumerate}
  \item 由 \cref{eq:gn-normal} 说明为什么收敛后左端 $\boldsymbol{\Lambda}$ 就是逆协方差，并写出取某变量边缘协方差的做法（不整体求逆）。
  \item （概念）解释 NEES $>N$（过自信）在工程上意味着什么、可能由什么引起。
  \item （推导，选做）对一维 $g(x)=fb/x$（立体深度）写出 \cref{eq:nlopt-mlbias} 的标量版，定性说明深度越大偏差越大。
\end{enumerate}
\end{practice}

% ════════════════════════════════════════════════════════════════
\section{高斯变分推断：MAP 与拉普拉斯的推广（进阶/选读）}
\label{sec:nlopt-gvi}

\begin{note}
\textbf{本节定位：进阶 / 选读 / 前沿。} 前面所有方法（GN/LM、Laplace 协方差）求的都是后验\emph{众数}加一个“以众数为心”的高斯近似——这对 SLAM 主管线（BA、位姿图、VIO）已足够，初读可跳过本节。但若你在意\emph{后验的形状本身}（均值、不确定度是否被系统性低估），或想把“估状态”与“估系统参数/协方差/偏置”纳入\emph{同一}框架，本节给出的\emph{高斯变分推断}（Gaussian Variational Inference，GVI）正是 MAP$+$Laplace 的严格推广。它与本章主线无缝衔接：第 \ref{sec:nlopt-cov} 节把“GN 左端 $=$ 逆协方差”讲成 Laplace 近似，本节将证明\emph{这恰是 GVI 在“只用均值算期望”时的退化}；它也接回 \cref{sec:est-fg} 的因子图——GVI 同样\emph{精确稀疏}，可复用 BA/平滑的稀疏后端。本节主要吸收 Barfoot\cite{barfoot2024state} 第 5 章（其原始论文为 Barfoot--Forbes--Yoon 的 ESGVI\cite{barfoot2020esgvi}），并对照 Opper--Archambeau\cite{opper2009variational} 的机器学习路线与 Amari\cite{amari1998natural} 的自然梯度。
\end{note}

\paragraph{动机：点估计丢掉了后验的“形状”。}
到此为止，本章把状态估计当成\emph{优化}：找一个点 $\hat{\mathbf{x}}$ 使代价最小，再用逆海森补一个协方差。这是\emph{自底向上}的视角（从线性高斯出发、逐步推广到非线性）。但贝叶斯的本来目标（\cref{thm:nlopt-map} 的 Step 1）是\emph{整个后验} $p(\mathbf{x}\mid\mathbf{z})$，不只是它的峰值。当观测非线性、后验被掰歪时（\cref{sec:ekf} 立体相机例），“峰值 $+$ 逆海森”这套近似会:(i) 把估计钉在\emph{众数}上而非\emph{均值}（\cref{sec:nlopt-cov} 的 ML 偏差之源）；(ii) 协方差\emph{过自信}（Laplace 只看峰值附近的曲率，看不到分布的真实展宽）。本节换一个\emph{自顶向下}的视角：直接问“哪一个\emph{高斯} $q(\mathbf{x})$ 最像真后验 $p(\mathbf{x}\mid\mathbf{z})$”，并把\emph{均值与协方差一起优化}。

\paragraph{反面：为什么不干脆用更多 sigma 点 / 粒子？}
\cref{ch:eskf} 已给出“抓均值”的采样路线（UKF/ISPKF、粒子滤波）。但它们是\emph{滤波}式的：在\emph{先验}处布点、过非线性、重组矩——点是按\emph{先验}协方差撒的（\cref{eq:sigma-points}）。问题是后验往往比先验窄得多、且位置偏移，先验点未必落在后验的高概率区。GVI 的关键差别是\emph{在不断改进的\textbf{后验}估计处布点}（下文“先验 vs 后验统计线性化”），并以一个明确的\emph{泛函}（而非“匹配前两阶矩”的启发式）为优化目标。粒子滤波则受维度灾难所限（\cref{tab:filter-taxo}），批量大状态下不可行。GVI 在“有明确目标 $+$ 可扩展到大状态”两点上同时改进。

\paragraph{历史：变分推断与 GVI。}
\emph{变分推断}（variational inference / variational Bayes）源自机器学习\cite{barfoot2024state}：用一族简单分布 $q$ 逼近难算的后验 $p$，把“积分推断”变成“优化 $q$”。早期多用\emph{可因式分解}（mean-field）的 $q$；用\emph{多元高斯} $q$ 的做法（即 GVI）一度因参数量 $O(N^2)$（协方差有 $N^2$ 个元）而被冷落，直到 Opper--Archambeau\cite{opper2009variational} 指出“高斯先验 $+$ 可因式分解似然”下变分参数其实只有 $O(N)$，并给出对高斯参数求导的关键公式。Barfoot 等\cite{barfoot2020esgvi} 把它系统地接进机器人批量估计，利用似然因式分解得到\emph{精确稀疏}的逆协方差（ESGVI），并配上 cubature/sigma 点实现\emph{无导数}推断。本节即沿这条线展开。

% ----------------------------------------------------------------
\subsection{损失泛函：负 ELBO 与 KL 散度}
\label{sec:nlopt-gvi-elbo}

\paragraph{为什么用 $\mathrm{KL}(q\Vert p)$ 而非 $\mathrm{KL}(p\Vert q)$。}
我们要让高斯 $q(\mathbf{x})=\mathcal{N}(\boldsymbol{\mu},\boldsymbol{\Sigma})$ 尽量接近真后验 $p(\mathbf{x}\mid\mathbf{z})$，用 KL 散度（\cref{ch:slam_est} 的 \cref{eq:gauss-entropy} 一带）度量“接近”。KL 不对称，有两种取法\cite{barfoot2024state}：
\begin{equation}
  \mathrm{KL}(p\Vert q)=\mathbb{E}_p\!\big[\ln p(\mathbf{x}\mid\mathbf{z})-\ln q(\mathbf{x})\big],\qquad
  \mathrm{KL}(q\Vert p)=\mathbb{E}_q\!\big[\ln q(\mathbf{x})-\ln p(\mathbf{x}\mid\mathbf{z})\big].
  \label{eq:gvi-kl-two}
\end{equation}
二者的本质差别在\emph{期望对谁取}：$\mathrm{KL}(p\Vert q)$ 的期望对\emph{真后验} $p$（我们\emph{不知道}、正想求的东西，无法直接算）；$\mathrm{KL}(q\Vert p)$ 的期望对\emph{我们自己的高斯估计} $q$（已知、可采样）。\textbf{这一点是 GVI 选 $\mathrm{KL}(q\Vert p)$ 的全部理由}：它把推断变成“对 $q$ 取期望再优化 $q$”的可迭代问题。\footnote{两种取法的几何后果不同：$\mathrm{KL}(q\Vert p)$ 是\emph{零强迫}（zero-forcing/mode-seeking）——$q$ 倾向贴住后验的一个主峰、可能低估展宽（多峰时只抓一个峰）；$\mathrm{KL}(p\Vert q)$ 是\emph{矩匹配}（mass-covering）——倾向覆盖整个支撑、可能跨峰摊平。Bishop\cite{bishop2006prml} 对此有详细讨论。机器人后验通常单峰，故 $\mathrm{KL}(q\Vert p)$ 的代价（轻微低估方差）可接受，而它带来的可计算性是决定性的。$\mathrm{KL}(q\Vert p)$ 的选择还自然引出对系统\emph{参数}的估计（\cref{sec:nlopt-gvi-param}）。}

\paragraph{推导负 ELBO 损失 $V(q)$。}
代入贝叶斯式 $p(\mathbf{x}\mid\mathbf{z})=p(\mathbf{x},\mathbf{z})/p(\mathbf{z})$ 到 $\mathrm{KL}(q\Vert p)$：
\begin{equation}
  \mathrm{KL}(q\Vert p)=\mathbb{E}_q[-\ln p(\mathbf{x},\mathbf{z})]
  \underbrace{-\tfrac12\ln\!\big((2\pi e)^N\det\boldsymbol{\Sigma}\big)}_{-(\text{$q$ 的熵})}
  +\underbrace{\ln p(\mathbf{z})}_{\text{常数}},
  \label{eq:gvi-kl-expand}
\end{equation}
其中用了高斯熵 $-\mathbb{E}_q[\ln q]=\tfrac12\ln((2\pi e)^N\det\boldsymbol{\Sigma})$（\cref{eq:gauss-entropy}）。末项 $\ln p(\mathbf{z})$（对数证据/log-evidence）\emph{不含 $q$}，优化时可丢。定义待最小化的\emph{损失泛函}
\begin{equation}
  \boxed{\ V(q)=\mathbb{E}_q\big[\phi(\mathbf{x})\big]+\tfrac12\ln\!\big(\det\boldsymbol{\Sigma}^{-1}\big),\qquad \phi(\mathbf{x})=-\ln p(\mathbf{x},\mathbf{z}).\ }
  \label{eq:gvi-V}
\end{equation}
这里\emph{刻意}把熵项写成 $+\tfrac12\ln\det\boldsymbol{\Sigma}^{-1}$（差一常数 $-\tfrac{N}{2}\ln(2\pi e)$），因为后文要利用\emph{信息阵} $\boldsymbol{\Sigma}^{-1}$ 的稀疏（协方差 $\boldsymbol{\Sigma}$ 稠密、$\boldsymbol{\Sigma}^{-1}$ 稀疏，\cref{sec:est-fg}）——故全程用 $(\boldsymbol{\mu},\boldsymbol{\Sigma}^{-1})$ 作为 $q$ 的完整描述。$\phi(\mathbf{x})=-\ln p(\mathbf{x},\mathbf{z})$ 就是本章一路用的\emph{负对数联合}（即 \cref{thm:nlopt-map} 那个 $\tfrac12\sum\|\mathbf{e}\|_{\boldsymbol{\Sigma}}^2$ 加常数）。

\begin{insight}[$V(q)$ 的两项在“拔河”，以及与 ELBO 的关系]
$V(q)$ 第一项 $\mathbb{E}_q[\phi]$ 拉 $q$ 去\emph{贴合数据}（让高概率区落在低代价处）；第二项 $\tfrac12\ln\det\boldsymbol{\Sigma}^{-1}$ 惩罚 $q$ \emph{过于自信}（$\boldsymbol{\Sigma}^{-1}$ 越大、即越确定，惩罚越大）。两者平衡出“既贴数据又不过分确定”的高斯。$V(q)$ 恰是所谓\emph{证据下界}（Evidence Lower BOund，ELBO）的\emph{相反数}：由 \cref{eq:gvi-kl-expand} 且 $\mathrm{KL}\ge0$ 得 $\ln p(\mathbf{z})\ge-V(q)+\text{const}$，即 $-V(q)$ 是对数证据 $\ln p(\mathbf{z})$ 的一个下界。\emph{最小化 $V$} $=$ \emph{最大化 ELBO} $=$ \emph{令 $q$ 逼近 $p$}（KL 最小）——三句话同义。这把 \cref{thm:nlopt-map} 的“MAP $=$ 最小化负对数”升级成了“GVI $=$ 最小化负 ELBO”。
\end{insight}

% ----------------------------------------------------------------
\subsection{GVI 更新：对均值与信息阵的牛顿式迭代}
\label{sec:nlopt-gvi-update}

\paragraph{$V(q)$ 对高斯参数的导数。}
要极小化 $V$，需它对 $\boldsymbol{\mu}$ 与 $\boldsymbol{\Sigma}^{-1}$ 的导数。经一番矩阵微积分（分子布局，记号同本章；推导见 \cref{der:nlopt-gvi-deriv}），可得\cite{opper2009variational,barfoot2024state}
\begin{subequations}\label{eq:gvi-derivs}
\begin{align}
  \frac{\partial V}{\partial\boldsymbol{\mu}^\top}&=\boldsymbol{\Sigma}^{-1}\,\mathbb{E}_q\big[(\mathbf{x}-\boldsymbol{\mu})\,\phi(\mathbf{x})\big],\label{eq:gvi-dmu}\\
  \frac{\partial^2 V}{\partial\boldsymbol{\mu}^\top\partial\boldsymbol{\mu}}&=\boldsymbol{\Sigma}^{-1}\,\mathbb{E}_q\big[(\mathbf{x}-\boldsymbol{\mu})(\mathbf{x}-\boldsymbol{\mu})^\top\phi(\mathbf{x})\big]\boldsymbol{\Sigma}^{-1}-\boldsymbol{\Sigma}^{-1}\,\mathbb{E}_q[\phi(\mathbf{x})],\label{eq:gvi-d2mu}\\
  \frac{\partial V}{\partial\boldsymbol{\Sigma}^{-1}}&=-\tfrac12\mathbb{E}_q\big[(\mathbf{x}-\boldsymbol{\mu})(\mathbf{x}-\boldsymbol{\mu})^\top\phi(\mathbf{x})\big]+\tfrac12\boldsymbol{\Sigma}\,\mathbb{E}_q[\phi(\mathbf{x})]+\tfrac12\boldsymbol{\Sigma}.\label{eq:gvi-dSinv}
\end{align}
\end{subequations}
对比 \cref{eq:gvi-d2mu} 与 \cref{eq:gvi-dSinv}，消去公共的 $\mathbb{E}_q[(\mathbf{x}-\boldsymbol{\mu})(\mathbf{x}-\boldsymbol{\mu})^\top\phi]$ 项，得一条\textbf{关键的 Hessian--梯度关系}：
\begin{equation}
  \frac{\partial^2 V}{\partial\boldsymbol{\mu}^\top\partial\boldsymbol{\mu}}=\boldsymbol{\Sigma}^{-1}-2\,\boldsymbol{\Sigma}^{-1}\frac{\partial V}{\partial\boldsymbol{\Sigma}^{-1}}\boldsymbol{\Sigma}^{-1}.
  \label{eq:gvi-hessgrad}
\end{equation}
它把“对 $\boldsymbol{\Sigma}^{-1}$ 的一阶导”与“对 $\boldsymbol{\mu}$ 的二阶导”绑在一起——正是构造更新格式的支点。

\paragraph{迭代格式（类牛顿）。}
一阶条件 $\partial V/\partial\boldsymbol{\mu}=\mathbf{0}$、$\partial V/\partial\boldsymbol{\Sigma}^{-1}=\mathbf{0}$ 一般无法对 $(\boldsymbol{\mu},\boldsymbol{\Sigma}^{-1})$ 解析求解，故迭代。把 $V$ 在当前 $q^{(i)}$ 处展开（对 $\delta\boldsymbol{\mu}$ 取二阶、对 $\delta\boldsymbol{\Sigma}^{-1}$ 取一阶——后者二阶导难算且协方差本就是 $\mathbf{x}$ 的二次量），令其下降。\emph{协方差更新}：在 \cref{eq:gvi-hessgrad} 中令 $\partial V/\partial\boldsymbol{\Sigma}^{-1}=\mathbf{0}$（取极值），立得
\begin{equation}
  \big(\boldsymbol{\Sigma}^{-1}\big)^{(i+1)}=\frac{\partial^2 V}{\partial\boldsymbol{\mu}^\top\partial\boldsymbol{\mu}}\bigg|_{q^{(i)}}.
  \label{eq:gvi-cov-update}
\end{equation}
\emph{均值更新}：仿 MAP 的牛顿步（\cref{eq:nlopt-newton}），对 $\delta\boldsymbol{\mu}$ 求导置零，得\emph{线性系统}
\begin{equation}
  \underbrace{\Big(\tfrac{\partial^2 V}{\partial\boldsymbol{\mu}^\top\partial\boldsymbol{\mu}}\big|_{q^{(i)}}\Big)}_{(\boldsymbol{\Sigma}^{-1})^{(i+1)}}\,\delta\boldsymbol{\mu}=-\Big(\tfrac{\partial V}{\partial\boldsymbol{\mu}^\top}\big|_{q^{(i)}}\Big),\qquad \boldsymbol{\mu}^{(i+1)}=\boldsymbol{\mu}^{(i)}+\delta\boldsymbol{\mu}.
  \label{eq:gvi-mean-update}
\end{equation}
\textbf{注意更新后的信息阵 $(\boldsymbol{\Sigma}^{-1})^{(i+1)}$ 正好充当均值更新的系数矩阵}（左端），与 GN 正规方程 \cref{eq:gn-normal} 同构——这不是巧合，而是“信息阵 $=$ 代价的曲率”这一主线在 GVI 中的再现。

\begin{theorem}[GVI 的局部收敛保证]\label{thm:nlopt-gvi-conv}
按 \cref{eq:gvi-cov-update,eq:gvi-mean-update} 更新，只要 $\delta\boldsymbol{\mu},\delta\boldsymbol{\Sigma}^{-1}$ 不同时为零，损失严格下降：
\begin{equation}
  V\!\big(q^{(i+1)}\big)-V\!\big(q^{(i)}\big)\approx
  -\tfrac12\underbrace{\delta\boldsymbol{\mu}^\top\big(\boldsymbol{\Sigma}^{-1}\big)^{(i+1)}\delta\boldsymbol{\mu}}_{\ge0}
  -\tfrac12\underbrace{\operatorname{tr}\!\big(\boldsymbol{\Sigma}^{(i)}\delta\boldsymbol{\Sigma}^{-1}\boldsymbol{\Sigma}^{(i)}\delta\boldsymbol{\Sigma}^{-1}\big)}_{\ge0}\ \le 0.
  \label{eq:gvi-descent}
\end{equation}
两项分别在 $\delta\boldsymbol{\mu}=\mathbf{0}$、$\delta\boldsymbol{\Sigma}^{-1}=\mathbf{0}$ 时取等。因此迭代单调降低 $V$，停在 $V$ 的一个局部极小（即 $\partial V/\partial\boldsymbol{\mu},\partial V/\partial\boldsymbol{\Sigma}^{-1}$ 均零处）。这是\emph{局部}保证（基于上述泰勒展开），与牛顿/GN 同性质。
\end{theorem}
\begin{proof}
把 \cref{eq:gvi-cov-update,eq:gvi-mean-update} 代回 $V$ 的二阶（对 $\delta\boldsymbol{\mu}$）／一阶（对 $\delta\boldsymbol{\Sigma}^{-1}$）泰勒展开即得 \cref{eq:gvi-descent}（细节见 \cref{der:nlopt-gvi-deriv}）。第一项 $\ge0$ 因 $(\boldsymbol{\Sigma}^{-1})^{(i+1)}$ 正定（作为协方差的逆）。第二项 $\ge0$ 需证：对实正定 $\mathbf{A}=\boldsymbol{\Sigma}^{(i)}$、实对称 $\mathbf{B}=\delta\boldsymbol{\Sigma}^{-1}$，有 $\operatorname{tr}(\mathbf{A}\mathbf{B}\mathbf{A}\mathbf{B})\ge0$ 且等号 $\iff\mathbf{B}=\mathbf{0}$。用 $\operatorname{vec}$ 与 Kronecker 积（\cref{ch:matderiv} 的 $\operatorname{tr}(\mathbf{X}^\top\mathbf{Y})=\operatorname{vec}(\mathbf{X})^\top\operatorname{vec}(\mathbf{Y})$、$\operatorname{vec}(\mathbf{A}\mathbf{B}\mathbf{A})=(\mathbf{A}\otimes\mathbf{A})\operatorname{vec}(\mathbf{B})$）：
\[
  \operatorname{tr}(\mathbf{A}\mathbf{B}\mathbf{A}\mathbf{B})=\operatorname{vec}(\mathbf{A}\mathbf{B}\mathbf{A})^\top\operatorname{vec}(\mathbf{B})=\operatorname{vec}(\mathbf{B})^\top(\mathbf{A}\otimes\mathbf{A})\operatorname{vec}(\mathbf{B})\ge0,
\]
因 $\mathbf{A}$ 正定 $\Rightarrow\mathbf{A}\otimes\mathbf{A}$ 正定，二次型非负、且 $=0\iff\operatorname{vec}(\mathbf{B})=\mathbf{0}\iff\mathbf{B}=\mathbf{0}$。\hfill$\square$
\end{proof}

\begin{insight}[GVI $\Rightarrow$ MAP$+$Laplace 是“只用均值算期望”的退化]
\cref{eq:gvi-cov-update,eq:gvi-mean-update} 里的期望 $\mathbb{E}_q[\cdot]$ 若\emph{只用 $q$ 的均值点近似}（即把分布塌成一个点 $\delta(\mathbf{x}-\boldsymbol{\mu})$，等价于只取一个 sigma 点），则（用下文 \cref{eq:gvi-stein} 的 Stein 形式）更新退化为
\[
  (\boldsymbol{\Sigma}^{-1})^{(i+1)}=\tfrac{\partial^2\phi}{\partial\mathbf{x}^\top\partial\mathbf{x}}\Big|_{\boldsymbol{\mu}},\qquad
  (\boldsymbol{\Sigma}^{-1})^{(i+1)}\delta\boldsymbol{\mu}=-\tfrac{\partial\phi}{\partial\mathbf{x}^\top}\Big|_{\boldsymbol{\mu}},
\]
这\emph{正是} \cref{sec:nlopt-cov} 的 MAP 牛顿步 $+$ Laplace 协方差（$\boldsymbol{\Sigma}^{-1}=$ 代价海森）！\textbf{所以 MAP$+$Laplace 是 GVI 的一阶（单点）近似}：GVI 用\emph{整个高斯}上的期望取代“只在峰值求导”，从而看见后验的展宽与偏移。这条退化关系是本节与 \cref{sec:nlopt-cov} 的正式接缝。
\end{insight}

% ----------------------------------------------------------------
\subsection{自然梯度视角}
\label{sec:nlopt-gvi-natgrad}

\paragraph{GVI 更新 $=$ 自然梯度下降。}
上面的更新还有一个漂亮的几何解读：它就是在变分参数上做\emph{自然梯度下降}（natural gradient descent，NGD）\cite{amari1998natural}。把 $q$ 的参数堆成一列 $\boldsymbol{\alpha}=[\boldsymbol{\mu};\operatorname{vec}(\boldsymbol{\Sigma}^{-1})]$。\emph{普通}梯度下降 $\delta\boldsymbol{\alpha}=-\partial V/\partial\boldsymbol{\alpha}^\top$ 在欧氏几何里走；但概率分布族的“距离”不是欧氏的（参数差一点、分布可能差很多），其自然度规是\emph{费舍尔信息阵} $\boldsymbol{\mathcal{I}}_{\boldsymbol{\alpha}}$（\cref{eq:laplace} 的费舍尔信息在此推广到“对变分参数”）。NGD 用它预条件：
\begin{equation}
  \delta\boldsymbol{\alpha}=-\boldsymbol{\mathcal{I}}_{\boldsymbol{\alpha}}^{-1}\frac{\partial V}{\partial\boldsymbol{\alpha}^\top}.
  \label{eq:gvi-ngd}
\end{equation}
对高斯 $q$，其变分参数的费舍尔信息恰好块对角 $\boldsymbol{\mathcal{I}}_{\boldsymbol{\alpha}}=\operatorname{diag}\!\big(\boldsymbol{\Sigma}^{-1},\,\tfrac12\boldsymbol{\Sigma}\otimes\boldsymbol{\Sigma}\big)$。代入 \cref{eq:gvi-ngd} 并用 $\operatorname{vec}(\mathbf{A}\mathbf{B}\mathbf{C})=(\mathbf{C}^\top\otimes\mathbf{A})\operatorname{vec}(\mathbf{B})$ 拆开，得
\begin{equation}
  \boldsymbol{\Sigma}^{-1}\delta\boldsymbol{\mu}=-\frac{\partial V}{\partial\boldsymbol{\mu}^\top},\qquad
  \delta\boldsymbol{\Sigma}^{-1}=-2\,\boldsymbol{\Sigma}^{-1}\frac{\partial V}{\partial\boldsymbol{\Sigma}^{-1}}\boldsymbol{\Sigma}^{-1},
  \label{eq:gvi-ngd-components}
\end{equation}
\textbf{与 \cref{sec:nlopt-gvi-update} 的更新逐字相同}（第二式即把 \cref{eq:gvi-cov-update} 写成增量形）。于是“类牛顿的 GVI 更新”与“自然梯度下降”是同一件事——这解释了 GVI 为何收敛得比朴素（欧氏）梯度快：它沿信息几何的“最陡”方向走。

% ----------------------------------------------------------------
\subsection{Stein 引理：把期望化成无导数形式}
\label{sec:nlopt-gvi-stein}

\paragraph{用 Stein 引理改写更新。}
\cref{eq:gvi-derivs} 里的期望含 $\phi$ 与 $(\mathbf{x}-\boldsymbol{\mu})$ 的乘积，直接算需 $\phi$ 的高阶矩。\cref{ch:slam_est} 的\emph{Stein 引理}（\cref{eq:stein}：$\mathbb{E}_q[(\mathbf{x}-\boldsymbol{\mu})f(\mathbf{x})]=\boldsymbol{\Sigma}\,\mathbb{E}_q[\partial f/\partial\mathbf{x}^\top]$）能把“坐标 $\times$ 函数”的期望换成“函数导数”的期望。它还有一条\emph{二阶}版（对二次可微 $f$）：
\begin{equation}
  \mathbb{E}_q\big[(\mathbf{x}-\boldsymbol{\mu})(\mathbf{x}-\boldsymbol{\mu})^\top f(\mathbf{x})\big]
  =\boldsymbol{\Sigma}\,\mathbb{E}_q\!\Big[\tfrac{\partial^2 f}{\partial\mathbf{x}^\top\partial\mathbf{x}}\Big]\boldsymbol{\Sigma}+\boldsymbol{\Sigma}\,\mathbb{E}_q[f],
  \label{eq:gvi-stein2}
\end{equation}
（由 \cref{eq:stein} 对 $\boldsymbol{\mu}$ 再求一次导、或直接分部积分可得。）把 \cref{eq:gvi-stein2} 与 \cref{eq:stein} 用到 \cref{eq:gvi-cov-update,eq:gvi-mean-update}（取 $f=\phi$），更新\emph{大幅简化}为
\begin{equation}
  \boxed{\ \big(\boldsymbol{\Sigma}^{-1}\big)^{(i+1)}=\mathbb{E}_{q^{(i)}}\!\Big[\tfrac{\partial^2\phi}{\partial\mathbf{x}^\top\partial\mathbf{x}}\Big],\qquad
  \big(\boldsymbol{\Sigma}^{-1}\big)^{(i+1)}\delta\boldsymbol{\mu}=-\mathbb{E}_{q^{(i)}}\!\Big[\tfrac{\partial\phi}{\partial\mathbf{x}^\top}\Big].\ }
  \label{eq:gvi-stein}
\end{equation}
即\emph{“信息阵 $=$ 代价 Hessian 的期望、梯度 $=$ 代价梯度的期望”}——形式与 MAP 牛顿步\emph{一模一样}，唯一差别是“在均值处求值”换成“对 $q$ 取期望”。这立刻坐实了前述退化洞见：只取均值点 $\Rightarrow$ 期望塌成点值 $\Rightarrow$ MAP$+$Laplace。

\begin{insight}[Stein 引理的双向用法]
Stein 引理在 GVI 里\emph{两次}出场、方向相反：(i) \emph{正用}（\cref{eq:gvi-derivs}$\to$\cref{eq:gvi-stein}）把“坐标$\times\phi$”的期望换成“$\phi$ 的导数”的期望，导出免高阶矩的牛顿式更新；(ii) \emph{反用}（下文 \cref{sec:nlopt-gvi-sparse}）再把“$\phi$ 的导数”的期望换回“坐标$\times\phi$”的期望，从而\emph{绕过对 $\phi$ 求导}——为 cubature/sigma 点的\emph{无导数}实现铺路。一来一回，既得到了简洁更新、又卸掉了可微性要求。
\end{insight}

% ----------------------------------------------------------------
\subsection{精确稀疏 ESGVI：接回因子图后端}
\label{sec:nlopt-gvi-sparse}

\paragraph{因式分解 $\Rightarrow$ 期望只需边缘。}
朴素 GVI 每步 $O(N^3)$（要 $N\times N$ 的稠密 $\boldsymbol{\Sigma}$），大状态下不可行。救星与 BA 同源：\emph{似然因式分解}。SLAM/BA 的联合似然按因子拆开（\cref{thm:nlopt-map} 的独立性、\cref{def:factor-graph} 的因子图），故负对数联合可写成因子之和
\begin{equation}
  \phi(\mathbf{x})=\sum_{k=1}^{K}\phi_k(\mathbf{x}_k),\qquad \phi_k(\mathbf{x}_k)=-\ln p(\mathbf{x}_k,\mathbf{z}_k),\quad \mathbf{x}_k=\mathbf{P}_k\mathbf{x},
  \label{eq:gvi-factored}
\end{equation}
其中 $\mathbf{x}_k$ 是第 $k$ 个因子涉及的\emph{少数}变量（$\mathbf{P}_k$ 为抽取它们的投影/膨胀矩阵）。代入 \cref{eq:gvi-stein} 的期望（以 Hessian 项为例）：
\begin{equation}
  \mathbb{E}_q\!\Big[\tfrac{\partial^2\phi}{\partial\mathbf{x}^\top\partial\mathbf{x}}\Big]
  =\sum_{k=1}^{K}\mathbf{P}_k^\top\,\mathbb{E}_{q}\!\Big[\tfrac{\partial^2\phi_k}{\partial\mathbf{x}_k^\top\partial\mathbf{x}_k}\Big]\mathbf{P}_k
  =\sum_{k=1}^{K}\mathbf{P}_k^\top\,\mathbb{E}_{q_k}\!\Big[\tfrac{\partial^2\phi_k}{\partial\mathbf{x}_k^\top\partial\mathbf{x}_k}\Big]\mathbf{P}_k.
  \label{eq:gvi-esgvi-hess}
\end{equation}
\textbf{最后一步是 ESGVI 的命门}：因 $\phi_k$ 只依赖 $\mathbf{x}_k$，对其余变量求导为零，故对\emph{全}高斯 $q$ 的期望\emph{精确}退化为对\emph{边缘} $q_k(\mathbf{x}_k)=\mathcal{N}(\mathbf{P}_k\boldsymbol{\mu},\,\mathbf{P}_k\boldsymbol{\Sigma}\mathbf{P}_k^\top)$ 的期望（\emph{不是近似}）。标量项 $\mathbb{E}_q[\phi]=\sum_k\mathbb{E}_{q_k}[\phi_k]$、梯度项同理。

\begin{insight}[ESGVI $=$ GVI 接回 BA 稀疏后端]
\cref{eq:gvi-esgvi-hess} 有三层含义，每层都把 GVI 拉回本章主线：
\begin{enumerate}
  \item \textbf{逆协方差精确稀疏}：$\boldsymbol{\Sigma}^{-1}=\sum_k\mathbf{P}_k^\top(\cdots)\mathbf{P}_k$ 的非零块由因子结构决定，且\emph{迭代中稀疏模式不变}——与 \cref{sec:est-fg} 的“$\boldsymbol{\Lambda}$ 稀疏镜像因子图”一字不差（批量轨迹估计是块三对角、SLAM 是箭头阵，正是 \cref{eq:ba-block} 的结构）。
  \item \textbf{永不构造稠密 $\boldsymbol{\Sigma}$}：只需各因子边缘 $q_k$，即只需 $\boldsymbol{\Sigma}$ 中\emph{对应 $\boldsymbol{\Sigma}^{-1}$ 非零块}的那些子块——巨大节省。均值更新 \cref{eq:gvi-mean-update} 的左端 $\boldsymbol{\Sigma}^{-1}$ 稀疏，可\emph{复用 BA 的稀疏 Cholesky/Schur 求解器}（\cref{sec:nlopt-solve,sec:nlopt-sparse}）。
  \item \textbf{同阶复杂度}：ESGVI 与 MAP \emph{同} $O(N)$ 量级（瓶颈都是稀疏分解），只是系数更大（要算边缘期望，而非单点求导）。
\end{enumerate}
一句话：\textbf{GVI 不是另起炉灶，而是在 BA/平滑的同一张因子图、同一套稀疏分解之上，把“点估计”升级成“分布估计”}。
\end{insight}

\paragraph{Takahashi 稀疏协方差回收：只回收要用的协方差块。}
还差一步：边缘 $q_k$ 需要 $\boldsymbol{\Sigma}$ 的子块 $\boldsymbol{\Sigma}_{kk}=\mathbf{P}_k\boldsymbol{\Sigma}\mathbf{P}_k^\top$，而我们手里只有稀疏的 $\boldsymbol{\Sigma}^{-1}$、不愿求整个稠密 $\boldsymbol{\Sigma}=(\boldsymbol{\Sigma}^{-1})^{-1}$。\emph{Takahashi 等}（1973，源于电路理论；Erisman--Tinney 1975 证明其封闭性）给出只回收\emph{需要的}协方差块的递归。先对信息阵作稀疏 $\mathbf{L}\mathbf{D}\mathbf{L}^\top$ 分解（$\mathbf{D}$ 对角、$\mathbf{L}$ 单位下三角且稀疏，$\boldsymbol{\Sigma}^{-1}=\mathbf{L}\mathbf{D}\mathbf{L}^\top$）。由 $\mathbf{L}\mathbf{D}\mathbf{L}^\top\boldsymbol{\Sigma}=\mathbf{I}$ 左乘 $(\mathbf{L}\mathbf{D})^{-1}$ 得 $\mathbf{L}^\top\boldsymbol{\Sigma}=\mathbf{D}^{-1}\mathbf{L}^{-1}$，转置并两边加 $\boldsymbol{\Sigma}-\boldsymbol{\Sigma}\mathbf{L}$：
\begin{equation}
  \boldsymbol{\Sigma}=\mathbf{L}^{-\top}\mathbf{D}^{-1}+\boldsymbol{\Sigma}(\mathbf{I}-\mathbf{L}).
  \label{eq:gvi-takahashi}
\end{equation}
因 $\boldsymbol{\Sigma}$ 对称、且 $\mathbf{L}^{-\top}$ 只影响上半块（可丢），\cref{eq:gvi-takahashi} 给出\emph{后向}递归（$j\ge k$，$\delta$ 为 Kronecker delta）：
\begin{equation}
  \boldsymbol{\Sigma}_{j,k}=\delta(j,k)\,\mathbf{D}_{k,k}^{-1}-\sum_{\ell=k+1}^{K}\boldsymbol{\Sigma}_{j,\ell}\,\mathbf{L}_{\ell,k}.
  \label{eq:gvi-takahashi-rec}
\end{equation}
\textbf{关键封闭性}：计算 $\boldsymbol{\Sigma}^{-1}$ 非零块对应的 $\boldsymbol{\Sigma}$ 块时，\cref{eq:gvi-takahashi-rec} 右端只用到\emph{已在计算清单上}的块（外加“盒子四角”规则补的几块：若 $\mathbf{L}_{k,i},\mathbf{L}_{j,i}\ne\mathbf{0}$ 则 $\mathbf{L}_{j,k}\ne\mathbf{0}$）——故\emph{只回收稀疏 $\boldsymbol{\Sigma}^{-1}$ 所需的那些协方差块即可}，不必碰稠密的其余部分。对批量估计（块三对角 $\boldsymbol{\Sigma}^{-1}$），这套回收可\emph{搭车}在均值求解 \cref{eq:gvi-mean-update} 的同一次分解上，复杂度同阶。这与 \cref{eq:nlopt-elim-cond} 变量消元的 $\boldsymbol{\Sigma}_j=(\mathbf{R}_j^\top\mathbf{R}_j)^{-1}$ 协方差回收是同一思想的完整版。

\paragraph{cubature / sigma 点：无导数实现。}
最后用 Stein 引理\emph{反用}（\cref{eq:gvi-stein2,eq:stein}）把“$\phi_k$ 的导数期望”换回“坐标$\times\phi_k$ 的期望”，再用\emph{cubature}（高斯求积，即 sigma 点）数值算这些边缘期望。一个 $N_k$ 维因子，取 $L$ 个 sigma 点 $\mathbf{x}_{k,\ell}=\boldsymbol{\mu}_k+\sqrt{\boldsymbol{\Sigma}_{kk}}\,\boldsymbol{\xi}_{k,\ell}$（$\boldsymbol{\xi}$ 为单位 sigma 点），则
\begin{equation}
  \mathbb{E}_{q_k}[\phi_k]\approx\sum_{\ell}w_{k,\ell}\,\phi_k(\mathbf{x}_{k,\ell}),\quad
  \mathbb{E}_{q_k}[(\mathbf{x}_k-\boldsymbol{\mu}_k)\phi_k]\approx\sum_{\ell}w_{k,\ell}(\mathbf{x}_{k,\ell}-\boldsymbol{\mu}_k)\phi_k(\mathbf{x}_{k,\ell}),
  \label{eq:gvi-cubature}
\end{equation}
二阶矩同理。权 $w$ 与单位点 $\boldsymbol{\xi}$ 取法即各种 cubature 规则：\emph{无迹变换}（\cref{eq:sigma-points} 的 $2N_k+1$ 点，权 $w_0=\kappa/(N_k+\kappa)$ 等）、\emph{球面--径向}（$2N_k$ 点、等权 $1/(2N_k)$）、\emph{Gauss--Hermite}（$M$ 阶对 $\le2M-1$ 次多项式精确，但需 $M^{N_k}$ 点）。\textbf{因期望在\emph{因子}级（$N_k$ 小），点数可控}，故 Gauss--Hermite 在机器人问题里反而常可负担。

\begin{insight}[ESGVI 的无导数推断：连鲁棒核都“免实现 IRLS”]
反用 Stein 后，\emph{无需对 $\phi_k$ 求任何导数}——只要能\emph{求值} $\phi_k(\mathbf{x}_{k,\ell})$。这有两个深远后果：(i) 观测/运动模型可以是\emph{黑箱}（不可微也行），与 UKF 同样的“免雅可比”优点，但作用在\emph{整个因子}（含先验项）上而非仅观测模型；(ii) 若因子里含\emph{鲁棒核}（\cref{sec:nlopt-robust}），它被\emph{自动}吸收进 $\phi_k$ 的取值里，\textbf{无须再单独实现 IRLS 重加权}（\cref{eq:irls}）——cubature 把“非线性 $+$ 鲁棒”一并处理。这是 ESGVI 相对 MAP$+$IRLS 的一个实在便利。
\end{insight}

\paragraph{$\dagger$ Gauss--Newton 型变体 ESGVI-GN 与正定性。}
若负对数因子取最小二乘形 $\phi(\mathbf{x})=\tfrac12\mathbf{e}(\mathbf{x})^\top\mathbf{W}^{-1}\mathbf{e}(\mathbf{x})$，可对 $V$ 用 Jensen 不等式构造一个\emph{更省}的代用泛函 $V'(q)=\tfrac12\mathbb{E}_q[\mathbf{e}]^\top\mathbf{W}^{-1}\mathbb{E}_q[\mathbf{e}]+\tfrac12\ln\det\boldsymbol{\Sigma}^{-1}$（把期望挪进二次型外），它\emph{隐式}地以 $\bar{\mathbf{E}}^\top\mathbf{W}^{-1}\bar{\mathbf{E}}$ 近似 Hessian（$\bar{\mathbf{E}}=\mathbb{E}_q[\partial\mathbf{e}/\partial\mathbf{x}]$ 为\emph{统计雅可比}），从而像 GN 一样\emph{绕过牛顿的二阶项}：
\begin{equation}
  \big(\boldsymbol{\Sigma}^{-1}\big)^{(i+1)}=\bar{\mathbf{E}}^\top\mathbf{W}^{-1}\bar{\mathbf{E}},\qquad
  \bar{\mathbf{E}}^\top\mathbf{W}^{-1}\bar{\mathbf{E}}\,\delta\boldsymbol{\mu}=-\bar{\mathbf{E}}^\top\mathbf{W}^{-1}\bar{\mathbf{e}}.
  \label{eq:gvi-gn}
\end{equation}
这\emph{就是}把 GN（\cref{eq:gn-normal}）里的“单点雅可比/残差”换成“对 $q$ 的期望版” $\bar{\mathbf{E}},\bar{\mathbf{e}}$。两个好处：cubature 积分阶数减半（$\mathbf{e}$ 而非 $\phi\sim\mathbf{e}^2$，点数 $M^{N_k}$ 中 $M$ 减半，常是“可算/不可算”之差）；且由 Jensen，$V'$ 给的 $\boldsymbol{\Sigma}^{-1}$ \emph{更保守}（不过自信），缓解了 $\mathrm{KL}(q\Vert p)$ 固有的“低估方差”倾向。代价是多一层近似，是否优于 MAP 视问题而定，常作 full ESGVI 的\emph{廉价预处理器}。另注：GVI 是\emph{无约束}优化，未强制 $\boldsymbol{\Sigma}^{-1}\succ0$；初值差时可能跑出非正定——补救是先跑 MAP$+$Laplace 到收敛再用其结果\emph{热启动} GVI（实践极有效、也省算力，因只在后期才用全期望精修），或借自然梯度（\cref{sec:nlopt-gvi-natgrad}）做带约束更新\cite{barfoot2024state}。

% ----------------------------------------------------------------
\subsection{$\dagger$ GVI 看线性系统与参数估计（EM/系统辨识）}
\label{sec:nlopt-gvi-param}

\paragraph{线性高斯：GVI 恢复批量解。}
GVI 推广 MAP，自然应在\emph{线性高斯}时退回经典批量解。把线性提升模型 $\phi(\mathbf{x})=\tfrac12(\mathbf{B}\mathbf{u}-\mathbf{A}^{-1}\mathbf{x})^\top\mathbf{Q}^{-1}(\mathbf{B}\mathbf{u}-\mathbf{A}^{-1}\mathbf{x})+\tfrac12(\mathbf{y}-\mathbf{C}\mathbf{x})^\top\mathbf{R}^{-1}(\mathbf{y}-\mathbf{C}\mathbf{x})$（提升记号同 \cref{ch:slam_est}）代入 \cref{eq:gvi-stein}：因 $\phi$ 是 $\mathbf{x}$ 的\emph{二次}型，其 Hessian 为常数、期望与分布无关，立得
\begin{equation}
  \hat{\boldsymbol{\Sigma}}^{-1}=\underbrace{\mathbf{A}^{-\top}\mathbf{Q}^{-1}\mathbf{A}^{-1}+\mathbf{C}^\top\mathbf{R}^{-1}\mathbf{C}}_{\text{块三对角}},\qquad
  \hat{\boldsymbol{\Sigma}}^{-1}\hat{\mathbf{x}}=\mathbf{A}^{-\top}\mathbf{Q}^{-1}\mathbf{B}\mathbf{u}+\mathbf{C}^\top\mathbf{R}^{-1}\mathbf{y},
  \label{eq:gvi-linear}
\end{equation}
\emph{正是} \cref{ch:slam_est} 的线性高斯批量正规方程（块三对角信息阵，等价 RTS 平滑）。即\textbf{线性时 GVI 一步收敛、与批量 MAP 重合；非线性时才给出与 MAP 不同（更接近均值）的解}——这与 \cref{thm:nlopt-map} 那个“线性时一步出解”的例子首尾呼应。

\paragraph{EM 估参数：把 $\boldsymbol{\theta}$ 也学出来。}
$\mathrm{KL}(q\Vert p)$ 的选择带来一个 MAP 难有的红利：把系统\emph{参数} $\boldsymbol{\theta}$（如 $\mathbf{A},\mathbf{B},\mathbf{C}$ 与协方差 $\mathbf{Q},\mathbf{R}$）也一并估。引入 $\boldsymbol{\theta}$ 后，数据负对数似然分解为 $-\ln p(\mathbf{z}\mid\boldsymbol{\theta})=-\mathrm{KL}(q\Vert p)+V(q\mid\boldsymbol{\theta})$，其中 $V(q\mid\boldsymbol{\theta})=\mathbb{E}_q[\phi(\mathbf{x}\mid\boldsymbol{\theta})]+\tfrac12\ln\det\boldsymbol{\Sigma}^{-1}$ 仍是负 ELBO。\emph{期望最大化}（EM）交替：\textbf{E 步}固定 $\boldsymbol{\theta}$、跑 GVI 至收敛得 $q(\mathbf{x})$（这步 ESGVI 已现成）；\textbf{M 步}固定 $q$、令 $\partial V(q\mid\boldsymbol{\theta})/\partial\boldsymbol{\theta}=\mathbf{0}$ 解 $\boldsymbol{\theta}$。因似然因式分解，M 步导数 $\partial V/\partial\boldsymbol{\theta}=\sum_k\mathbb{E}_{q_k}[\partial\phi_k/\partial\boldsymbol{\theta}]$ 同样只需\emph{边缘}期望（仍用 E 步已回收的 $\boldsymbol{\Sigma}$ 块 $+$ cubature）。例：估观测协方差 $\mathbf{W}$，对 $\mathbf{W}^{-1}$ 求导置零得
\begin{equation}
  \mathbf{W}=\frac1K\sum_{k=1}^{K}\mathbb{E}_{q_k}\!\big[\mathbf{e}_k(\mathbf{x}_k)\mathbf{e}_k(\mathbf{x}_k)^\top\big],
  \label{eq:gvi-em-W}
\end{equation}
即“误差外积的（边缘期望）样本平均”——形似 \cref{ch:eskf} 的自适应协方差，但这里用\emph{平滑}估计且\emph{含状态不确定度的膨胀}修正（期望 $\mathbb{E}_{q_k}$ 自动把 $\boldsymbol{\Sigma}_{kk}$ 算进去）。这就是用结构化隐变量学系统矩阵的\emph{系统辨识}（精神上同 HMM 的 Baum--Welch），且\emph{始终不需稠密 $\boldsymbol{\Sigma}$}。\footnote{这把 \cref{sec:nlopt-robust} 的“鲁棒核 $=$ 协方差逆 Wishart 先验 MAP”推向更一般的“在变分框架内联合估计状态与噪声模型”：鲁棒化、协方差自适应、系统辨识，在 GVI 里是同一台机器的不同旋钮。}

% ----------------------------------------------------------------
\subsection{GVI 滤波校正步：与 ISPKF 的“先验 vs 后验线性化”对照}
\label{sec:nlopt-gvi-filter}

\paragraph{把 GVI 用在单步校正。}
GVI 不止用于批量，也能只做\emph{一步滤波校正}（这恰是立体相机例的设定）。预测步给出先验 $\check{\mathbf{x}},\check{\mathbf{P}}$，校正把先验项与观测项相加 $\phi(\mathbf{x})=\tfrac12(\mathbf{x}-\check{\mathbf{x}})^\top\check{\mathbf{P}}^{-1}(\mathbf{x}-\check{\mathbf{x}})+\tfrac12(\mathbf{y}-\mathbf{g}(\mathbf{x}))^\top\mathbf{R}^{-1}(\mathbf{y}-\mathbf{g}(\mathbf{x}))$。设当前 GVI 后验 $q=\mathcal{N}(\mathbf{x}_{\mathrm{op}},\mathbf{P}_{\mathrm{op}})$，用 \cref{eq:gvi-stein}（GN 近似 $\phi$ 的二阶项）得更新
\begin{equation}
  \mathbf{P}_{\mathrm{op}}\leftarrow\Big(\check{\mathbf{P}}^{-1}+\mathbb{E}_q\big[\mathbf{G}^\top\mathbf{R}^{-1}\mathbf{G}\big]\Big)^{-1},\quad
  \mathbf{x}_{\mathrm{op}}\leftarrow\mathbf{x}_{\mathrm{op}}+\mathbf{P}_{\mathrm{op}}\Big(\mathbb{E}_q\big[\mathbf{G}^\top\mathbf{R}^{-1}(\mathbf{y}-\mathbf{g})\big]+\check{\mathbf{P}}^{-1}(\check{\mathbf{x}}-\mathbf{x}_{\mathrm{op}})\Big),
  \label{eq:gvi-filter}
\end{equation}
$\mathbf{G}=\partial\mathbf{g}/\partial\mathbf{x}$，期望用 cubature 算（亦可反用 Stein 走无导数）。切换到 GN 型代用泛函（\cref{eq:gvi-gn}）并施 SMW，可整理成标准卡尔曼形 $\mathbf{K}=\check{\mathbf{P}}\bar{\mathbf{G}}^\top(\mathbf{R}+\bar{\mathbf{G}}\check{\mathbf{P}}\bar{\mathbf{G}}^\top)^{-1}$，其中 $\bar{\mathbf{G}}=\mathbb{E}_q[\partial\mathbf{g}/\partial\mathbf{x}]$ 是\emph{统计雅可比}——可用 Stein 反式 $\bar{\mathbf{G}}=\mathbb{E}_q[(\mathbf{g}-\bar{\mathbf{g}})(\mathbf{x}-\mathbf{x}_{\mathrm{op}})^\top]\mathbf{P}_{\mathrm{op}}^{-1}$ 由 sigma 点算出。

\begin{insight}[先验 vs 后验统计线性化：GVI 与 ISPKF 的分水岭]
\cref{ch:eskf} 的 ISPKF（\cref{ex:three-methods}、\cref{tab:filter-taxo}）也“用 sigma 点统计线性化、迭代趋均值”，与 \cref{eq:gvi-filter} 形似，但有一处\emph{本质}差别——\textbf{sigma 点按哪个分布撒}：
\begin{itemize}
  \item \textbf{ISPKF：先验线性化}。sigma 点按\emph{先验}协方差 $\check{\mathbf{P}}$ 撒（迭代只动均值/工作点，撒点的协方差仍是先验）。故称\emph{先验统计线性化}。
  \item \textbf{GVI：后验线性化}。sigma 点按\emph{不断改进的后验}协方差 $\mathbf{P}_{\mathrm{op}}$ 撒（\cref{eq:gvi-filter} 里期望 $\mathbb{E}_q$ 的 $q$ 就是当前后验）。故称\emph{后验统计线性化}。
\end{itemize}
后验通常比先验\emph{窄且偏移}，按后验撒点更贴近真正起作用的高概率区，故 GVI 的线性化更准、更逼近真后验。这是 GVI 在下例几乎无偏的根本原因，也是它\emph{不}等同于 ISPKF 的关键\cite{barfoot2024state}。
\end{insight}

\paragraph{（选读）立体相机：GVI vs MAP/IEKF vs ISPKF 三方对照（闭环既有贯穿例）。}
回到 \cref{sec:ekf} 那个贯穿全书的一维立体相机例（深度 $x$、视差 $y=fb/x+n$，参数同 \cref{eq:stereo-bayes} 一带）。\cref{ch:eskf} 已算出：在“众数 vs 均值”上，IEKF/MAP 落在被非线性拉偏的\emph{众数}、ISPKF 接近\emph{均值}（\cref{ex:three-methods} 的数值 $\hat{x}_{\text{iekf}}=\hat{x}_{\text{map}}=24.5694$、$\bar{x}=24.7770$、$\hat{x}_{\text{ispkf}}=24.7414$）。现在加入 GVI：用 $M=3$ 个 sigma 点、按\emph{后验}协方差迭代估均值与方差，得 $\hat{x}_{\text{gvi}}=24.7792$——在所有方法里\emph{最贴近真后验均值} $\bar{x}=24.7770$。
\begin{table}[H]\centering\small
\caption{立体相机校正步：四法对照（参数同 \cref{sec:ekf}；$\bar{x}=24.7770$ 为真后验均值；偏差为 $10^5$ 次试验的 $\hat{e}_{\text{mean}}=\mathbb{E}[\hat{x}]-\check{x}$）。GVI 几乎无偏，因它\emph{显式}地最小化“与真后验的 KL”、且按后验线性化。}
\label{tab:nlopt-gvi-stereo}
\begin{tabular}{lllll}
\toprule
方法 & 线性化在何处 & 抓后验的 & 单步估 $\hat{x}$ [m] & 系统偏差 $\hat{e}_{\text{mean}}$ \\
\midrule
MAP / IEKF & 工作点（$\to$众数） & 众数 & $24.5694$ & $\approx-33.0$\,cm \\
ISPKF & \emph{先验} $\check{\mathbf{P}}$ & 均值（近似） & $24.7414$ & $\approx-3.84$\,cm \\
GVI（本节） & \emph{后验} $\mathbf{P}_{\mathrm{op}}$ & 均值$+$形状（KL 最优） & $\mathbf{24.7792}$ & $\boldsymbol{\approx0.28}$\,\textbf{cm} \\
真后验（基准） & — & — & $\bar{x}=24.7770$ & — \\
\bottomrule
\end{tabular}
\end{table}
\noindent 在“系统偏差”这一指标上，GVI（$\approx0.28$\,cm）比 MAP（$\approx-33.0$\,cm）小\emph{两个数量级}、也优于 ISPKF（$\approx-3.84$\,cm）；均方误差三者相当（$\hat{e}_{\text{sq}}\approx4.28$\,m$^2$，与 \cref{sec:ekf} 的 $4.41$\,m$^2$ 同量级）。\textbf{为什么 GVI 几乎无偏}：它\emph{直接}以“与真后验的 KL 最小”为目标（而非顺手取个峰值），且按后验撒点使统计线性化落在对的位置——于是它估到的均值几乎就是真后验均值，而“偏差 $=$ 估计期望 $-$ 真值期望”自然趋零。这正是本章开头 \cref{sec:nlopt-map} 那个“MAP 给众数、众数 $\ne$ 均值故有偏”pitfall 的\emph{彻底解药}：要无偏，就别只取众数，去估\emph{整个}（最优）高斯。

% ----------------------------------------------------------------
\subsection{$\dagger$ 补证：IEKF 收敛到后验众数即 MAP}
\label{sec:nlopt-gvi-iekf}

承 \cref{sec:nlopt-gn} 的伏笔，这里把“IEKF $=$ 单步 MAP 的 GN”补成定理与完整证明（不外引）。它也正是上节“GVI 在只用均值算期望时退回 MAP/IEKF”的具体例证。

\begin{theorem}[IEKF 收敛到后验众数（$=$ 单步 MAP）]\label{thm:nlopt-iekf-map}
单步滤波校正问题（先验 $\check{\mathbf{x}},\check{\mathbf{P}}$、观测 $\mathbf{y}=\mathbf{g}(\mathbf{x})+\mathbf{n}$，$\mathbf{n}\sim\mathcal{N}(\mathbf{0},\mathbf{R})$）的 MAP 估计即
\begin{equation}
  \hat{\mathbf{x}}=\arg\min_{\mathbf{x}}\ J(\mathbf{x}),\qquad
  J(\mathbf{x})=\tfrac12\|\mathbf{x}-\check{\mathbf{x}}\|_{\check{\mathbf{P}}}^2+\tfrac12\|\mathbf{y}-\mathbf{g}(\mathbf{x})\|_{\mathbf{R}}^2.
  \label{eq:iekf-map-cost}
\end{equation}
迭代 EKF（\cref{eq:iekf}）的不动点 $\mathbf{x}_{\mathrm{op}}^\star$ 恰满足 $J$ 的一阶最优条件 $\partial J/\partial\mathbf{x}|_{\mathbf{x}_{\mathrm{op}}^\star}=\mathbf{0}$；且 IEKF 的每步更新\emph{就是}对 \cref{eq:iekf-map-cost} 的一次高斯牛顿迭代。故 IEKF（若收敛）收敛到 \cref{eq:iekf-map-cost} 的（局部）极小，即后验 $p(\mathbf{x}\mid\mathbf{y})\propto\exp(-J)$ 的\emph{众数}。
\end{theorem}
\begin{proof}
\emph{Step 1（GN 步）}。\cref{eq:iekf-map-cost} 是标准非线性最小二乘：把先验当一个“残差” $\mathbf{e}_p=\mathbf{x}-\check{\mathbf{x}}$（权 $\check{\mathbf{P}}^{-1}$、雅可比 $\mathbf{I}$）、观测当残差 $\mathbf{e}_m=\mathbf{y}-\mathbf{g}(\mathbf{x})$（权 $\mathbf{R}^{-1}$、雅可比 $-\mathbf{G}$，$\mathbf{G}=\partial\mathbf{g}/\partial\mathbf{x}|_{\mathbf{x}_{\mathrm{op}}}$）。堆叠雅可比 $\mathbf{J}=[\mathbf{I};-\mathbf{G}]$、权 $\mathbf{W}^{-1}=\operatorname{diag}(\check{\mathbf{P}}^{-1},\mathbf{R}^{-1})$，GN 正规方程 \cref{eq:gn-normal} 在工作点 $\mathbf{x}_{\mathrm{op}}$ 处为
\[
  \big(\check{\mathbf{P}}^{-1}+\mathbf{G}^\top\mathbf{R}^{-1}\mathbf{G}\big)\delta\mathbf{x}
  =-\check{\mathbf{P}}^{-1}(\mathbf{x}_{\mathrm{op}}-\check{\mathbf{x}})+\mathbf{G}^\top\mathbf{R}^{-1}\big(\mathbf{y}-\mathbf{g}(\mathbf{x}_{\mathrm{op}})\big).
\]
\emph{Step 2（化成 IEKF 形）}。记 $\boldsymbol{\Sigma}^{-1}=\check{\mathbf{P}}^{-1}+\mathbf{G}^\top\mathbf{R}^{-1}\mathbf{G}$。对它用 SMW（\cref{ch:slam_est}）得 $\boldsymbol{\Sigma}=\check{\mathbf{P}}-\mathbf{K}\mathbf{G}\check{\mathbf{P}}$，其中卡尔曼增益 $\mathbf{K}=\check{\mathbf{P}}\mathbf{G}^\top(\mathbf{G}\check{\mathbf{P}}\mathbf{G}^\top+\mathbf{R})^{-1}$。解出 $\delta\mathbf{x}$、令 $\mathbf{x}_{\mathrm{op}}^{+}=\mathbf{x}_{\mathrm{op}}+\delta\mathbf{x}$，整理（用 $\boldsymbol{\Sigma}\mathbf{G}^\top\mathbf{R}^{-1}=\mathbf{K}$、$\boldsymbol{\Sigma}\check{\mathbf{P}}^{-1}=\mathbf{I}-\mathbf{K}\mathbf{G}$）得
\[
  \mathbf{x}_{\mathrm{op}}^{+}=\check{\mathbf{x}}+\mathbf{K}\big(\mathbf{y}-\mathbf{g}(\mathbf{x}_{\mathrm{op}})-\mathbf{G}(\check{\mathbf{x}}-\mathbf{x}_{\mathrm{op}})\big),
\]
\emph{这与 IEKF 更新 \cref{eq:iekf} 逐字相同}（再令 $\mathbf{x}_{\mathrm{op}}\leftarrow\mathbf{x}_{\mathrm{op}}^{+}$ 重线性化）。故 IEKF 的每次迭代就是对 $J$ 的一次 GN 步。
\emph{Step 3（不动点 $=$ 驻点）}。设迭代收敛到 $\mathbf{x}_{\mathrm{op}}^\star$，则 $\delta\mathbf{x}=\mathbf{0}$，由 Step 1 的正规方程右端为零：
\[
  \check{\mathbf{P}}^{-1}(\mathbf{x}_{\mathrm{op}}^\star-\check{\mathbf{x}})-\mathbf{G}^{\star\top}\mathbf{R}^{-1}\big(\mathbf{y}-\mathbf{g}(\mathbf{x}_{\mathrm{op}}^\star)\big)=\mathbf{0}.
\]
而这正是 \cref{eq:iekf-map-cost} 的梯度 $\partial J/\partial\mathbf{x}=\check{\mathbf{P}}^{-1}(\mathbf{x}-\check{\mathbf{x}})-\mathbf{G}^\top\mathbf{R}^{-1}(\mathbf{y}-\mathbf{g}(\mathbf{x}))$ 在 $\mathbf{x}_{\mathrm{op}}^\star$ 处为零的条件。故不动点是 $J$ 的驻点；GN 在正定 $\boldsymbol{\Sigma}^{-1}$ 下每步为下降方向，收敛点为（局部）极小，即后验 $\propto\exp(-J)$ 的众数。\hfill$\square$
\end{proof}

\begin{insight}
\cref{thm:nlopt-iekf-map} 把三条线缝在一起：\textbf{IEKF（滤波）$=$ 单步 MAP 的 GN（优化）$=$ GVI 只用均值算期望的退化（推断）}。沿这条缝看：EKF 是其“只迭代一次”的草率版（\cref{prop:ekf}）；ISPKF 把求众数换成求均值（\cref{tab:filter-taxo}）；GVI 则把“点 $+$ 逆海森”升级成“整个最优高斯”。滤波、优化、变分推断，本是同一座山的三条上山路。
\end{insight}

% ----------------------------------------------------------------
\begin{pitfall}[概念误区：GVI 只是“迭代更多次的 UKF / 带 sigma 点的 MAP”]
三者貌合神离。\textbf{UKF/ISPKF} 按\emph{先验}撒点、以“匹配前两阶矩”为启发式、无统一目标泛函；\textbf{MAP$+$Laplace} 只在\emph{众数}处求导、协方差是事后补的逆海森（GVI 的单点退化）；\textbf{GVI} 有明确目标——最小化与真后验的 $\mathrm{KL}(q\Vert p)$（即负 ELBO \cref{eq:gvi-V}），按\emph{后验}撒点，\emph{联合}优化均值与协方差。后果上的差别：GVI 的均值几乎无偏（\cref{tab:nlopt-gvi-stereo}）、协方差不像 Laplace 那样系统性过自信。但 GVI 也更贵（每步要边缘期望 $+$ 协方差回收），且 $\mathrm{KL}(q\Vert p)$ 仍可能轻微低估方差（用 \cref{eq:gvi-gn} 的保守变体缓解）、对多峰后验只抓一个峰。\emph{选型}：SLAM 主管线点估计够用就 MAP/GN；要校准的不确定度或联合估参数/协方差，才上 GVI。
\end{pitfall}

\begin{pitfall}[实现陷阱：在朴素 GVI 里显式构造稠密 $\boldsymbol{\Sigma}$]
GVI 的 $\boldsymbol{\Sigma}$ 一般\emph{稠密}（$N\times N$），大状态下既存不下也求不动逆。\textbf{切勿}照搬 \cref{eq:gvi-cov-update} 字面去求整个 $\boldsymbol{\Sigma}=(\boldsymbol{\Sigma}^{-1})^{-1}$。正解是 ESGVI（\cref{sec:nlopt-gvi-sparse}）：只存稀疏 $\boldsymbol{\Sigma}^{-1}$、用 Takahashi 递归（\cref{eq:gvi-takahashi-rec}）只回收因子边缘要用的那几个 $\boldsymbol{\Sigma}$ 块、复用 BA 的稀疏分解。把它当“点估计 BA 多带一组协方差块回收”来实现，复杂度同阶。
\end{pitfall}

\begin{practice}
\begin{enumerate}
  \item （推导）由 $\mathrm{KL}(q\Vert p)$ 与高斯熵 \cref{eq:gauss-entropy} 推出负 ELBO 损失 \cref{eq:gvi-V}，说明为何末项 $\ln p(\mathbf{z})$ 可丢、为何用 $\boldsymbol{\Sigma}^{-1}$ 而非 $\boldsymbol{\Sigma}$ 记熵项。
  \item （推导，草稿纸）用 Hessian--梯度关系 \cref{eq:gvi-hessgrad} 导出协方差更新 \cref{eq:gvi-cov-update}；再把更新代回泰勒展开验证下降式 \cref{eq:gvi-descent} 第一项。
  \item 证明 \cref{thm:nlopt-gvi-conv} 用到的 $\operatorname{tr}(\mathbf{A}\mathbf{B}\mathbf{A}\mathbf{B})\ge0$（$\mathbf{A}\succ0$、$\mathbf{B}$ 对称），并说明等号条件。
  \item （概念）解释“只用 $q$ 的均值点近似期望”为何让 GVI 退回 MAP$+$Laplace（\cref{eq:gvi-stein}），并据此说明 MAP 为何有偏。
  \item （核心）说清 ESGVI 的两步如何接回因子图：(i) 似然因式分解 $\Rightarrow$ 期望只需边缘 $q_k$、$\boldsymbol{\Sigma}^{-1}$ 精确稀疏（\cref{eq:gvi-esgvi-hess}）；(ii) Takahashi 递归只回收稀疏 $\boldsymbol{\Sigma}^{-1}$ 对应的 $\boldsymbol{\Sigma}$ 块（\cref{eq:gvi-takahashi-rec}）。与 \cref{sec:nlopt-sparse} 的 Schur/变量消元对照。
  \item （对照）用“先验 vs 后验统计线性化”一句话说清 GVI 与 ISPKF（\cref{ex:three-methods}）的本质差别，并解释 \cref{tab:nlopt-gvi-stereo} 中 GVI 偏差为何远小于 MAP。
  \item （编程，选做）对一维立体相机例（\cref{eq:stereo-map}）实现 \cref{eq:gvi-filter} 的 GVI 校正（$M=3$ sigma 点、按 $\mathbf{P}_{\mathrm{op}}$ 撒点、迭代），复现 $\hat{x}_{\text{gvi}}\approx24.78$ 与偏差 $\approx0.28$\,cm，并与 MAP/IEKF、ISPKF 对比。
  \item （证明）完整重做 \cref{thm:nlopt-iekf-map}：从 \cref{eq:iekf-map-cost} 的 GN 步出发，经 SMW 化成 IEKF 更新 \cref{eq:iekf}，再说明不动点 $=$ 后验众数。
\end{enumerate}
\end{practice}

% ════════════════════════════════════════════════════════════════
\section{可微优化：把后端嵌进端到端学习（进阶/选读）}
\label{sec:nlopt-diffopt}

\begin{note}
\textbf{本节定位：进阶 / 选读 / 前沿。}\;到此为止，本章都把优化当作\emph{给定}的问题来\emph{解}（GN/LM 求最优解、Laplace/GVI 求最优解附近的分布）。本节反过来问一个新问题：\emph{如何对优化的解 $\mathbf{x}^*$ 关于问题的参数 $\mathbf{y}$ 求导}——从而让一个\emph{学习前端}（神经网络）能从\emph{几何后端}（BA/位姿图）的最优解反向受益。这与 \cref{sec:nlopt-gvi} 的 GVI 同属“优化的高阶视角”：GVI 升级\emph{输出}（点$\to$分布），本节升级\emph{接口}（让梯度穿过优化）。它接回 \cref{sec:nlopt-cov} 的“最优解二阶结构”（隐式微分要用海森）与 \cref{thm:nlopt-iekf-map} 的 IEKF（同样关心最优解对扰动如何变）。SLAM 主管线点估计够用就不必本节；要把“学习的前端”与“几何的后端”缝成\emph{端到端可训练}的系统，才需要它。本节方法用本书口吻陈述，库放\cref{sec:nlopt-practice} 与延伸。
\end{note}

\paragraph{动机：分开训练前后端会累积误差。}\cite{carlone2026handbook}
现代 SLAM 常\emph{前端学习 $+$ 后端几何}：前端用神经网络做特征检测/匹配或运动初估（\cref{ch:vo}），后端用几何优化（BA/位姿图）求全局一致解。若只把学习前端的输出\emph{单向}喂给后端，两模块各自为政、误差累积。一个新趋势是让信息\emph{双向}流动：后端把它的最优解\emph{反馈}给前端，使前端直接从后端估计里学习。这天然是一个\emph{双层优化}（bilevel optimization，BLO）：上层训练前端（参数 $\mathbf{y}$，无物理意义，如网络权重），\emph{约束于}下层的几何优化（参数 $\mathbf{x}$，有物理意义，如相机位姿/路标）。

\paragraph{双层优化的结构。}\cite{carlone2026handbook}
记上层损失 $\mathcal{U}(\mathbf{y},\mathbf{x}^*)$、下层非线性最小二乘 $\mathcal{L}(\mathbf{y},\mathbf{x})$：
\begin{equation}
  \mathbf{y}^*=\arg\min_{\mathbf{y}}\ \mathcal{U}\big(\mathbf{y},\mathbf{x}^*(\mathbf{y})\big)
  \quad\text{s.t.}\quad
  \mathbf{x}^*(\mathbf{y})=\arg\min_{\mathbf{x}}\ \mathcal{L}(\mathbf{y},\mathbf{x}).
  \label{eq:nlopt-blo}
\end{equation}
\emph{两例（SLAM 语境）}：(a) \textbf{学习特征 $+$ BA}——前端网络（参数 $\mathbf{y}$）做特征提取/匹配，后端 BA（变量 $\mathbf{x}=$ 位姿$+$路标）解重投影最小二乘；上层 $\mathcal{U}$ 取特征匹配误差、下层 $\mathcal{L}$ 取重投影误差，最优位姿 $\mathbf{x}^*$ 充当训练网络的监督信号，故优化 \cref{eq:nlopt-blo} 可经“反传 BA 重投影误差”进一步压低匹配误差。(b) \textbf{学习前端 $+$ 位姿图}——前端网络估相对位姿、后端位姿图（\cref{sec:nlopt-posegraph}）消漂移；上下层都用位姿图误差，差别只在上层优化 $\mathbf{y}$、下层优化 $\mathbf{x}$，于是前端能经“反传后端位姿残差”学到\emph{全局}几何一致性。求解 BLO 靠对 $\mathbf{y}$ 做梯度下降，故需上层梯度
\begin{equation}
  \nabla_{\mathbf{y}}\mathcal{U}
  =\frac{\partial\mathcal{U}(\mathbf{y},\mathbf{x}^*)}{\partial\mathbf{y}}
  +\frac{\partial\mathcal{U}(\mathbf{y},\mathbf{x}^*)}{\partial\mathbf{x}^*}\,
  \frac{\partial\mathbf{x}^*(\mathbf{y})}{\partial\mathbf{y}}.
  \label{eq:nlopt-blo-grad}
\end{equation}
两个\emph{直接}梯度项易得；难点是\emph{间接}项 $\partial\mathbf{x}^*/\partial\mathbf{y}$——“最优解如何随问题参数变”。有两条求法：\emph{展开微分}（unrolled）与\emph{隐式微分}（implicit）。

\paragraph{途径一：展开微分。}\cite{carlone2026handbook}
把下层迭代器（如 GN/梯度下降的每步 $\mathbf{x}_t=\Phi_t(\mathbf{x}_{t-1};\mathbf{y})$）\emph{展开}成一张可自动微分（AutoDiff）的计算图，用 $\mathbf{x}_T\approx\mathbf{x}^*$ 代入、对整条链 $\mathbf{x}_T=(\Phi_T\circ\cdots\circ\Phi_0)(\mathbf{y})$ 反传即得 $\partial\mathbf{x}_T/\partial\mathbf{y}$。可用反向模式（reverse-mode，类似 BP）或前向模式。缺点：迭代步数 $T$ 大、$\mathbf{x}$ 与 $\mathbf{y}$ 高维时\emph{费时费内存}（要存全部中间量）。补救：\emph{截断展开}（只存最后 $M$ 步，理论上梯度质量可与完整展开相当）；更极端地只做\emph{一步}迭代，并用\emph{有限差分}绕过显式海森——对 $\mathbf{x}_0^\pm=\mathbf{x}_0\pm\epsilon\,\partial\mathcal{U}/\partial\mathbf{x}_1$ 算 $\big(\partial\mathcal{L}(\mathbf{x}_0^+,\mathbf{y})/\partial\mathbf{y}-\partial\mathcal{L}(\mathbf{x}_0^-,\mathbf{y})/\partial\mathbf{y}\big)/(2\epsilon)$ 近似间接项整体。\emph{注意}：若变量在流形（李群），扰动须用\emph{右扰动} $\boxplus$（\cref{sec:nlopt-manifold}）而非普通加法，雅可比/海森-向量积也按右扰动算（\cref{ch:lie}）。

\paragraph{途径二：隐式微分（自包含推导）。}\cite{carlone2026handbook}
展开微分要把整个迭代过程记下来；隐式微分则\emph{绕开迭代}、直接对下层的\emph{最优性条件}求导，是本节的核心。其思想即微积分里的隐函数求导：由方程 $R(\mathbf{x},\mathbf{y})=\mathbf{0}$ 定义的隐函数 $\mathbf{x}(\mathbf{y})$，不必显式解出 $\mathbf{x}$，只需对 $R=\mathbf{0}$ 两边\emph{全微分}再解出 $\mathrm{d}\mathbf{x}/\mathrm{d}\mathbf{y}$。在这里，下层最优解 $\mathbf{x}^*(\mathbf{y})$ 是驻点，满足一阶最优性条件（设 $\mathcal{L}$ 对 $\mathbf{x},\mathbf{y}$ 二次可微）
\begin{equation}
  \frac{\partial\mathcal{L}\big(\mathbf{x}^*(\mathbf{y}),\mathbf{y}\big)}{\partial\mathbf{x}^*}=\mathbf{0}.
  \label{eq:nlopt-implicit-opt}
\end{equation}
把 \cref{eq:nlopt-implicit-opt} 看成关于 $\mathbf{y}$ 的恒等式，两边对 $\mathbf{y}$ 求全导数（用链式法则，注意 $\mathbf{x}^*$ 经 $\mathbf{y}$ 隐含依赖）：
\begin{equation}
  \underbrace{\frac{\partial^2\mathcal{L}}{\partial\mathbf{x}^*\partial\mathbf{y}}}_{\textstyle\mathbf{H}_{xy}}
  +\underbrace{\frac{\partial^2\mathcal{L}}{\partial\mathbf{x}^*\partial\mathbf{x}^*}}_{\textstyle\mathbf{H}_{xx}}\,
  \frac{\partial\mathbf{x}^*(\mathbf{y})}{\partial\mathbf{y}}=\mathbf{0}.
  \label{eq:nlopt-implicit-totaldiff}
\end{equation}
$\mathbf{H}_{xx}=\partial^2\mathcal{L}/\partial\mathbf{x}^*\partial\mathbf{x}^*$ 是下层海森（在 NLS 里即 \cref{thm:gn} 的 $\boldsymbol{\Lambda}=\mathbf{J}^\top\mathbf{W}^{-1}\mathbf{J}$ 之类的近似海森），$\mathbf{H}_{xy}=\partial^2\mathcal{L}/\partial\mathbf{x}^*\partial\mathbf{y}$ 是交叉二阶导。$\mathbf{x}^*$ 是极小、$\mathbf{H}_{xx}\succ0$ 可逆，解出\emph{间接梯度}
\begin{equation}
  \boxed{\ \frac{\partial\mathbf{x}^*(\mathbf{y})}{\partial\mathbf{y}}=-\,\mathbf{H}_{xx}^{-1}\,\mathbf{H}_{xy}\ }
  \label{eq:nlopt-implicit-grad}
\end{equation}
这就把“变量 $\mathbf{x},\mathbf{y}$ 间的间接梯度”化成了“$\mathcal{L}$ 的直接二阶导”，代价是一次海森求逆。代入 \cref{eq:nlopt-blo-grad} 得
\begin{equation}
  \nabla_{\mathbf{y}}\mathcal{U}
  =\frac{\partial\mathcal{U}}{\partial\mathbf{y}}
  -\underbrace{\frac{\partial\mathcal{U}}{\partial\mathbf{x}^*}}_{\textstyle\mathbf{v}^\top}\,
  \mathbf{H}_{xx}^{-1}\,\mathbf{H}_{xy}
  =\frac{\partial\mathcal{U}}{\partial\mathbf{y}}-\mathbf{q}^\top\mathbf{H}_{xy},
  \qquad \mathbf{q}:=\mathbf{H}_{xx}^{-1}\mathbf{v}.
  \label{eq:nlopt-implicit-final}
\end{equation}

\paragraph{为何不直接求海森逆：海森-向量积。}\cite{carlone2026handbook}
\cref{eq:nlopt-implicit-grad} 的 $\mathbf{H}_{xx}^{-1}$ 直接算\emph{不现实}：上下层若各含一个仅 $10^6$ 参数的网络，参数本身存 $4$\,MB，其海森却要 $(10^6)^2\times4$\,Byte $=4$\,TB——存都存不下。但 \cref{eq:nlopt-implicit-final} 只需 $\mathbf{q}=\mathbf{H}_{xx}^{-1}\mathbf{v}$ 这\emph{一个}向量，等价于解线性系统 $\mathbf{H}_{xx}\mathbf{q}=\mathbf{v}$（与本章一贯的“分解/迭代解线性系统、绝不显式求逆”\cref{sec:nlopt-solve} 同调）。用共轭梯度或梯度下降解它，每步只需\emph{海森-向量积}（Hessian-vector product，HVP）$\mathbf{H}_{xx}\mathbf{q}$，而\emph{HVP $=$“梯度-向量积”的梯度}：
\begin{equation}
  \mathbf{H}_{xx}\,\mathbf{q}=\frac{\partial^2\mathcal{L}}{\partial\mathbf{x}\partial\mathbf{x}}\,\mathbf{q}
  =\frac{\partial}{\partial\mathbf{x}}\Big(\underbrace{\tfrac{\partial\mathcal{L}}{\partial\mathbf{x}}\cdot\mathbf{q}}_{\text{标量}}\Big),
  \label{eq:nlopt-hvp}
\end{equation}
即先算梯度 $\partial\mathcal{L}/\partial\mathbf{x}$ 与 $\mathbf{q}$ 的内积（一个标量），再对它做一次 AutoDiff——\textbf{全程不显式构造或存储 $\mathbf{H}_{xx}$}。同理 \cref{eq:nlopt-implicit-final} 末项的 $\mathbf{H}_{xy}\mathbf{q}$ 也用 HVP 算。这把“反传穿过 BA”落成了可执行的算法：解一个以近似海森为系数的小线性系统、配 HVP。

\begin{insight}[展开 vs 隐式：存计算图还是解线性系统]
两条途径各有取舍。\emph{展开微分}把下层迭代当普通计算图、复用深度学习的 AutoDiff，实现直接，但要存全部中间迭代（$T$ 步、高维时昂贵），且梯度质量依赖展开够不够长。\emph{隐式微分}只用最优性条件 \cref{eq:nlopt-implicit-opt}，与到达 $\mathbf{x}^*$ 的\emph{路径无关}、内存省（只解一个线性系统 \cref{eq:nlopt-hvp}），但要求 $\mathbf{x}^*$ \emph{确实收敛到驻点}、且实现更繁。一个常用近似：\emph{忽略}间接项、只取直接梯度 $\nabla_{\mathbf{y}}\mathcal{U}\approx\partial\mathcal{U}/\partial\mathbf{y}|_{\mathbf{x}_T}$（等价于把下层解 $\mathbf{x}_T$ 当上层的\emph{常量}），更省但引入误差 $\sim|\partial\mathcal{U}/\partial\mathbf{x}^*\cdot\partial\mathbf{x}^*/\partial\mathbf{y}|$——当间接梯度含\emph{小}二阶导之积时可接受。
\end{insight}

\begin{pitfall}[把可微 BA 当“BA 只是网络的一层”]
有一类做法只把位姿误差“穿过”BA 反传、监督仍来自\emph{真值位姿}（BA 退化成网络的一个特殊层）——这\emph{没}真正双向连接前后端，前端学不到后端的几何先验，常\emph{次优}。真正的 BLO（\cref{eq:nlopt-blo}）让前端从\emph{几何优化本身}（而非外部真值）学习，可\emph{无需外部监督}就改进。实现上别忘了：下层变量含位姿时，\cref{eq:nlopt-implicit-grad,eq:nlopt-hvp} 的导数都按\emph{右扰动}（\cref{sec:nlopt-manifold}）算，否则梯度方向与流形几何不符（现代库已支持李群的 Jacobian-/Hessian-向量积）。
\end{pitfall}

\paragraph{生态与趋势。}\cite{carlone2026handbook,agarwal2012ceres}
把可微优化嵌进学习的库有 Theseus、PyPose、CvxpyLayer 等（在 PyTorch 等框架上提供可微 NLS 层、并支持李群 AutoDiff 与二阶优化器）；它们与本章 \cref{sec:nlopt-practice} 的 Ceres/g2o 互补——后者解几何优化、前者让该优化\emph{可微}。\cref{sec:nlopt-practice} 提过 FastTriggs（\cref{sec:nlopt-robust} 的一种只用一阶导的 IRLS 校正，比 Triggs 更稳）即出自此类库。代表系统：把 BA 层挂进网络的 BA-Net、DROID-SLAM，优化“位姿-码图”的 DeepFactors，以及\emph{双向}连接前后端、无外部监督即提升的 iSLAM。注意这些库/系统本节只作指引，方法本身（unrolled/implicit/HVP）已自包含于上文。

% ════════════════════════════════════════════════════════════════
\section{可证最优：从局部到全局（进阶/选读）}
\label{sec:nlopt-certifiable}

\begin{note}
\textbf{本节定位：进阶 / 选读 / 前沿。}\;\cref{sec:nlopt-why} 说过 GN/LM 只保证\emph{局部}最优、初值差易陷局部极小（一个经典反例是同一位姿图实例从不同随机初值收敛到多个不同的局部极小），并埋下“可证最优求解器（SE-Sync）留作前沿”的钩子；\cref{sec:nlopt-robust} 末也把它列为应对鲁棒 SLAM 的前沿出路。本节兑现该钩子，回答一个 \cref{sec:nlopt-why} 没答的问题：\emph{能否不靠初值、可证全局最优地解位姿图}？答案是\emph{半定松弛}（semidefinite relaxation）——把非凸的二次约束二次规划（QCQP）松弛成凸的半定规划（SDP），在松弛\emph{精确}时还附带一张事后\emph{最优性证书}。这与本章主线（GN/LM 求局部解）形成“局部 vs 全局”的对照，与 \cref{sec:nlopt-diffopt}/\cref{sec:nlopt-gvi} 同属“非线性优化的高阶视角”。方法用本书口吻陈述，定理“陈述 $+$ 引用证明”但陈述自包含。（注：本节只讲\emph{求解器}；“最优解离真值多近”的信息论极限——费舍尔信息与图拉普拉斯的联系——属估计精度的范畴，见 \cref{ch:slam_est}。）
\end{note}

\paragraph{背景：SLAM 非凸，却常意外收敛——但并非总是。}\cite{carlone2026handbook,rosen2019sesync}
位姿图、旋转平均等 SLAM 子问题已知\emph{NP 难}，一般无高效全局算法（除非 P$=$NP）。早期 SLAM 用里程计初值（漂移可压到行程的 $1\%$）常意外收敛到接近真值的解，一度被归功于好初值；但后续发现即便初值差，随机梯度下降/LM/预条件共轭梯度也常恢复近最优解——说明 SLAM \emph{有内在结构}使它格外好解。然而反例也存在：现成的局部优化（梯度下降、拟牛顿）在\emph{良定}的 SLAM 实例上仍可能给出离谱错误（如停车场多层位姿图从随机初值陷入劣质局部极小\cite{rosen2019sesync}），且局部法\emph{无法判别}自己是否陷在局部极小，带来可靠性隐患。这催生了基于\emph{凸松弛}（而非局部优化）的\emph{可证最优}（certifiably correct/optimal）方法：算出解后能量化其离全局最优有多远、甚至证明其为全局最优。这不与 NP 难矛盾——最坏情形下它仍可能给不出证书、只给次优界；但实践中它对绝大多数 SLAM 问题都能出证书，且\emph{给不出证书本身就是预警}（提示下游别信任此估计、或令机器人启动失效保护）。

% ----------------------------------------------------------------
\subsection{Shor 松弛：QCQP 到 SDP}
\label{sec:nlopt-certifiable-shor}

\paragraph{母机：把秩-1 提升后丢掉秩约束。}\cite{rosen2019sesync,carlone2026handbook}
许多 SLAM 问题可写成\emph{二次约束二次规划}（quadratically constrained quadratic program，QCQP）——目标与约束都是二次型：
\begin{equation}
  p^*=\min_{\mathbf{x}\in\mathbb{R}^n}\ \mathbf{x}^\top\mathbf{C}\mathbf{x}
  \quad\text{s.t.}\quad \mathbf{x}^\top\mathbf{A}_i\mathbf{x}=b_i,\ i=1,\dots,m,
  \label{eq:nlopt-qcqp}
\end{equation}
$\mathbf{C},\mathbf{A}_i$ 对称。\emph{Shor 松弛}用几步代数把它松成凸问题。先用迹的循环性把二次型写成线性形：对任意对称 $\mathbf{M}$，$\mathbf{x}^\top\mathbf{M}\mathbf{x}=\operatorname{tr}(\mathbf{M}\mathbf{x}\mathbf{x}^\top)$。于是 \cref{eq:nlopt-qcqp} 里 $\mathbf{x}$ \emph{只经外积} $\mathbf{X}\triangleq\mathbf{x}\mathbf{x}^\top$ 出现。注意：$\mathbf{X}=\mathbf{x}\mathbf{x}^\top$ \emph{当且仅当} $\mathbf{X}$ 对称、半正定（$\mathbf{X}\succeq0$）且\emph{秩 1}。代换变量得等价问题
\begin{equation}
  p^*=\min_{\mathbf{X}\in\mathbb{S}^n}\ \operatorname{tr}(\mathbf{C}\mathbf{X})
  \quad\text{s.t.}\quad \operatorname{tr}(\mathbf{A}_i\mathbf{X})=b_i,\ \mathbf{X}\succeq0,\ \operatorname{rank}(\mathbf{X})=1.
  \label{eq:nlopt-qcqp-lifted}
\end{equation}
目标与等式约束现在对矩阵变量 $\mathbf{X}$ \emph{线性}、$\mathbf{X}\succeq0$ \emph{凸}——唯一的非凸来自\emph{秩-1}约束。Shor 松弛就一招：\textbf{丢掉秩约束}，得凸的\emph{半定规划}
\begin{equation}
  d^*=\min_{\mathbf{X}\in\mathbb{S}^n}\ \operatorname{tr}(\mathbf{C}\mathbf{X})
  \quad\text{s.t.}\quad \operatorname{tr}(\mathbf{A}_i\mathbf{X})=b_i,\ \mathbf{X}\succeq0.
  \label{eq:nlopt-sdp}
\end{equation}
SDP 是凸的、可多项式时间解。

\paragraph{次优性界与“松弛精确即全局最优”。}\cite{rosen2019sesync}
因 \cref{eq:nlopt-sdp} 把可行集\emph{放大}了（丢了秩约束），$d^*\le p^*$。这已经有用：对任一局部解 $\hat{\mathbf{x}}$（如 GN/LM 得到的），其次优性被 $f(\hat{\mathbf{x}})-p^*\le f(\hat{\mathbf{x}})-d^*$ 界住（$f(\mathbf{x})=\mathbf{x}^\top\mathbf{C}\mathbf{x}$），右端可算（$d^*$ 由 SDP 给出）——无需知道真正的 $p^*$。更妙：若解出的 SDP 最优 $\mathbf{X}^*$ \emph{恰好秩 1}，$\mathbf{X}^*=\mathbf{x}^*\mathbf{x}^{*\top}$，则 $\mathbf{x}^*$ 在原 \cref{eq:nlopt-qcqp} 可行且达到 $d^*$，故\emph{就是原非凸问题的全局最优}。关键事实是：对许多 SLAM 任务，这一“秩 1（精确）”情形\emph{真的经常发生}，于是能从凸 SDP 反解出非凸问题的全局解。（“提升”的代价是维度从 $n$ 涨到 $\tfrac{n(n+1)}{2}$——这既解释了松弛为何能绕开非凸，也提示大问题需专用求解器。）

% ----------------------------------------------------------------
\subsection{SE-Sync：可证最优的位姿图优化}
\label{sec:nlopt-certifiable-sesync}

把 Shor 松弛用到\emph{位姿图优化}（PGO，\cref{sec:nlopt-posegraph}）即得 SE-Sync\cite{rosen2019sesync}——一个可证最优的 PGO 求解器。PGO 是最简单常用的 SLAM 形式，也是第一个被证明可凸松弛的 SLAM 问题，思路后被推广到一大类问题。分三步：化 PGO 为只含旋转的 QCQP $\to$ Shor 松弛并证其精确 $\to$ 用 Riemannian Staircase 高效求解。

\begin{insight}[前置：代数图论（关联阵、拉普拉斯、Fiedler 值）]
SE-Sync 的化简离不开把图写成矩阵（与 \cref{sec:nlopt-sparse} 的“稀疏镜像因子图”一脉相承）。对有向图 $\overrightarrow{\mathcal{G}}=(\mathcal{V},\mathcal{E})$（$n$ 节点、$m$ 边）：\emph{关联矩阵}（incidence）$\mathbf{A}\in\mathbb{R}^{m\times n}$ 每行对应一条边 $e=(i,j)$，仅第 $i$ 列为 $-1$、第 $j$ 列为 $+1$；\emph{拉普拉斯矩阵} $\mathbf{L}\triangleq\mathbf{A}^\top\mathbf{A}\in\mathbb{R}^{n\times n}$ 的对角元 $\mathbf{L}_{ii}$ 是节点 $i$ 的度、非对角元 $\mathbf{L}_{ij}=-1$（$i,j$ 间有边、不论方向）。$\mathbf{L}$ 的最小特征值恒为 $0$（对应全 $1$ 向量，因每行和为零）；零特征值的个数 $=$ 连通分量数，故连通图恰有一个零特征值，其\emph{第二小}特征值衡量图“连得多好”，称\emph{代数连通度}或 \emph{Fiedler 值}。后文 SE-Sync 的数据矩阵正由\emph{加权}拉普拉斯（边权为测量精度）搭成；图连得越好，松弛越易精确（也越利于估计精度，\cref{ch:slam_est}）。
\end{insight}

\paragraph{第一步：PGO 的极大似然与只含旋转的 QCQP。}\cite{rosen2019sesync}
PGO 由相对位姿测量 $\tilde{\mathbf{T}}_{ij}\approx\mathbf{T}_i^{-1}\mathbf{T}_j$ 估一组位姿 $\mathbf{T}_i=(\mathbf{t}_i,\mathbf{R}_i)\in\SE(d)$。取一个让 MLE 代数形式简单的噪声模型：平移 $\tilde{\mathbf{t}}_{ij}$ 受各向同性高斯噪声（精度 $\tau_{ij}$）、旋转 $\tilde{\mathbf{R}}_{ij}$ 受\emph{各向同性 Langevin 噪声}（精度 $\kappa_{ij}$；Langevin 分布是旋转流形 $\SO(d)$ 上类比高斯的指数族分布，密度 $\propto\exp(\kappa\operatorname{tr}(\mathbf{M}^\top\mathbf{R}))$）。这给出极大似然
\begin{equation}
  p^*_{\mathrm{MLE}}=\min_{\substack{\mathbf{t}_i\in\mathbb{R}^d,\ \mathbf{R}_i\in\SO(d)}}
  \sum_{(i,j)\in\overrightarrow{\mathcal{E}}}\kappa_{ij}\big\|\mathbf{R}_j-\mathbf{R}_i\tilde{\mathbf{R}}_{ij}\big\|_F^2
  +\tau_{ij}\big\|\mathbf{t}_j-\mathbf{t}_i-\mathbf{R}_i\tilde{\mathbf{t}}_{ij}\big\|_2^2.
  \label{eq:nlopt-pgo-mle}
\end{equation}
目标是\emph{二次}最小二乘，约束 $\mathbf{R}_i\in\SO(d)$ 在 $d{=}2,3$ 可写成二次约束——故这是 QCQP。\emph{关键化简}：固定旋转 $\mathbf{R}$，\cref{eq:nlopt-pgo-mle} 对平移 $\mathbf{t}$ 是线性最小二乘，可\emph{解析消去}平移 $\mathbf{t}^*(\mathbf{R})$（解含加权拉普拉斯 $\mathbf{L}(\mathcal{W}^\tau)$ 的伪逆，正是上面 insight 的对象），代回得\emph{只含旋转}的等价问题
\begin{equation}
  p^*_{\mathrm{MLE}}=\min_{\mathbf{R}\in\SO(d)^n}\ \operatorname{tr}\!\big(\tilde{\mathbf{Q}}\,\mathbf{R}^\top\mathbf{R}\big),
  \label{eq:nlopt-pgo-rot}
\end{equation}
$\mathbf{R}=(\mathbf{R}_1,\dots,\mathbf{R}_n)$，数据矩阵 $\tilde{\mathbf{Q}}$ 由旋转的\emph{连接拉普拉斯}加上一项由平移测量经关联阵投影构成（此 $\tilde{\mathbf{Q}}$ 非过程噪声 $\boldsymbol{\Sigma}_w$，只是沿用 Handbook 记号的数据矩阵；其可由稀疏分解高效构造）。\cref{eq:nlopt-pgo-rot} 形如标准\emph{多旋转平均}（multiple rotation averaging），只是数据矩阵更复杂。

\paragraph{第二步：Shor 松弛及其精确性。}\cite{rosen2019sesync}
先把 $\mathbf{R}\in\SO(d)^n$ 松到 $\mathbf{R}\in\mathrm{O}(d)^n$（正交阵由二次正交性约束定义，使之成齐次 QCQP），再仿 \cref{eq:nlopt-sdp}：把 $\mathbf{R}^\top\mathbf{R}$ 换成一个秩-$d$ 的 $\mathbf{Z}\succeq0$、丢秩约束，得 PGO 的半定松弛
\begin{equation}
  p^*_{\mathrm{SDP}}=\min_{\mathbf{Z}\in\mathbb{S}^{dn}_+}\ \operatorname{tr}(\tilde{\mathbf{Q}}\mathbf{Z})
  \quad\text{s.t.}\quad \operatorname{BlockDiag}_{d\times d}(\mathbf{Z})=(\mathbf{I}_d,\dots,\mathbf{I}_d),
  \label{eq:nlopt-pgo-sdp}
\end{equation}
（块对角约束即每个 $d\times d$ 对角块为 $\mathbf{I}_d$，对应正交性）。由 Shor 松弛 $p^*_{\mathrm{SDP}}\le p^*_{\mathrm{MLE}}$；若解出的 $\mathbf{Z}^*$ 能分解为 $\mathbf{Z}^*=\mathbf{R}^{*\top}\mathbf{R}^*$ 且 $\mathbf{R}^*\in\SO(d)^n$，则 $\mathbf{R}^*$ 即 \cref{eq:nlopt-pgo-rot} 的全局最优。能否如此，由下面定理保证（证明见\cite{rosen2019sesync}）。

\begin{theorem}[PGO 半定松弛的精确性]\label{thm:nlopt-sesync-exact}
设 $\bar{\mathbf{Q}}$ 是用\emph{真值}相对位姿构造的数据矩阵。存在常数 $\beta(\bar{\mathbf{Q}})>0$，使得当测量噪声足够小（$\|\tilde{\mathbf{Q}}-\bar{\mathbf{Q}}\|_2<\beta$）时：(1) 半定松弛 \cref{eq:nlopt-pgo-sdp} 有\emph{唯一}解 $\mathbf{Z}^*$；(2) $\mathbf{Z}^*=\mathbf{R}^{*\top}\mathbf{R}^*$，其中 $\mathbf{R}^*\in\SO(d)^n$ 是极大似然 \cref{eq:nlopt-pgo-rot} 的全局最小。
\end{theorem}

\noindent 即：噪声不太大时，解凸的 \cref{eq:nlopt-pgo-sdp} 就能恢复 PGO 的\emph{全局}最优（再经 $\mathbf{t}^*(\mathbf{R}^*)$ 回平移）。这是 SE-Sync 之所以可证最优的根。

\paragraph{第三步：Riemannian Staircase 高效解 SDP。}\cite{rosen2019sesync,carlone2026handbook}
直接解 \cref{eq:nlopt-pgo-sdp}（稠密 $\mathbf{Z}$ 维度 $dn$）对通用内点法太大。诀窍是\emph{利用低秩解}：\cref{thm:nlopt-sesync-exact} 表明欲求的 $\mathbf{Z}^*$ 秩为 $d$、即便不精确也秩不大，故可用 \emph{Burer--Monteiro 分解} $\mathbf{Z}=\mathbf{Y}^\top\mathbf{Y}$（$\mathbf{Y}\in\mathbb{R}^{r\times dn}$，$r$ 略大于 $d$）把 \cref{eq:nlopt-pgo-sdp} 重写成
\begin{equation}
  \min_{\mathbf{Y}\in\mathrm{St}(d,r)^n}\ \operatorname{tr}\!\big(\tilde{\mathbf{Q}}\,\mathbf{Y}^\top\mathbf{Y}\big),
  \qquad \mathrm{St}(d,r)=\{\mathbf{Y}\in\mathbb{R}^{r\times d}:\mathbf{Y}^\top\mathbf{Y}=\mathbf{I}_d\},
  \label{eq:nlopt-pgo-bm}
\end{equation}
块对角约束化为“每块 $\mathbf{Y}_i$ 的列是标准正交标架”，即 $\mathbf{Y}_i$ 在\emph{Stiefel 流形} $\mathrm{St}(d,r)$ 上——于是 $\mathbf{Y}^\top\mathbf{Y}\succeq0$ 自动满足、无需显式半定约束，且变量从 $\tfrac{dn(dn+1)}{2}$ 降到 $rnd$（$r\ll dn$）。\cref{eq:nlopt-pgo-bm} 的可行集是 Stiefel 流形之积，可\emph{复用本章的流形优化机器}（\cref{sec:nlopt-manifold} 的 retraction $+$ 二阶黎曼优化）求局部解，远快于通用约束规划。重引入正交约束又带来非凸，但下面定理给出“何时局部解即全局解”：

\begin{theorem}[全局最优的充分条件]\label{thm:nlopt-staircase}
若 $\mathbf{Y}\in\mathrm{St}(d,r)^n$ 是 \cref{eq:nlopt-pgo-bm} 的一个\emph{行秩亏}的二阶临界点，则 $\mathbf{Y}$ 是 \cref{eq:nlopt-pgo-bm} 的全局最小，且 $\mathbf{Z}^*=\mathbf{Y}^\top\mathbf{Y}$ 是半定松弛 \cref{eq:nlopt-pgo-sdp} 的解。
\end{theorem}

\noindent 这立刻给出 \emph{Riemannian Staircase}（黎曼阶梯）：从小秩 $r\ge d{+}1$ 起，用二阶黎曼局部优化解 \cref{eq:nlopt-pgo-bm} 得二阶临界点 $\mathbf{Y}^*$；若 $\mathbf{Y}^*$ 行秩亏，\cref{thm:nlopt-staircase} 证它已是全局最优；否则\emph{升一级} $r\leftarrow r{+}1$ 再解。因 $r\ge dn{+}1$ 时任意 $\mathbf{Y}$ 必行秩亏，阶梯有限步必终止——而实践常一两级即足。它像一个\emph{轻量元算法}“裹住”SLAM 常用的快速二阶局部求解器：既保留其速度，又\emph{额外}保证全局最优。

\paragraph{舍入与最优性证书。}\cite{rosen2019sesync}
得低秩因子 $\mathbf{Y}^*$ 后\emph{舍入}回 $\SO(d)^n$：松弛精确时 $\mathbf{Y}^*$ 实秩为 $d$，对其作\emph{薄 SVD} 即得 $\mathbf{R}^*$；不精确时取 $\mathbf{Y}^*$ 的秩-$d$ 最优近似、再把各 $d\times d$ 块用 SVD 投影回 $\SO(d)$ 得可行近似解 $\hat{\mathbf{R}}$。最后给\emph{最优性证书}：SE-Sync 返回可行解 $\hat{\mathbf{T}}$ 与下界 $p^*_{\mathrm{SDP}}\le p^*_{\mathrm{MLE}}$，从而任一可行 $\mathbf{T}$ 的次优性被 $F(\tilde{\mathbf{Q}}\mathbf{R}^\top\mathbf{R})-p^*_{\mathrm{SDP}}$ 界住；\emph{当松弛精确时}舍入解\emph{达到}该下界
\begin{equation}
  F\!\big(\tilde{\mathbf{Q}}\hat{\mathbf{R}}^\top\hat{\mathbf{R}}\big)=p^*_{\mathrm{SDP}},
  \label{eq:nlopt-sesync-cert}
\end{equation}
事后验证 \cref{eq:nlopt-sesync-cert} 成立即\emph{证明} $\hat{\mathbf{T}}$ 为全局最优——这就是“可证最优”一词的实义：不只给解，还给一张可独立核验的最优性证书。SE-Sync 在标准 PGO 基准上恢复全局最优，且速度不逊于（甚至快于）传统局部求解器。

\begin{insight}[可证最优 $=$ 凸松弛 $+$ 事后证书；与本章局部法互补]
SE-Sync 把 \cref{sec:nlopt-why} 那个“局部法可能陷局部极小、还\emph{不自知}”的隐患\emph{对症解决}：凸松弛绕过初值依赖、证书让你\emph{知道}解是否全局最优。但它\emph{不取代} GN/LM——Riemannian Staircase 内层用的正是本章的二阶流形局部优化；可证最优是“在局部求解器外裹一层全局保证”。两者关系恰如本章一贯的分工：局部法快、可证法稳而有保证。给不出证书时（极端噪声）退化为“给一个次优界 $+$ 预警”，仍比盲目相信局部解强。
\end{insight}

% ----------------------------------------------------------------
\subsection{扩展：地标、测距、各向异性与外点}
\label{sec:nlopt-certifiable-ext}

SE-Sync 的机器可推广到更多 SLAM 问题，但不同问题对“松弛是否精确”给出不同答案\cite{rosen2019sesync,carlone2026handbook}。

\paragraph{地标 SLAM：路标当纯平移变量。}\cite{rosen2019sesync}
带方位$+$距离（即相对位置）测量的地标 SLAM 可\emph{原样}套用 SE-Sync——把路标当“无旋转分量的位姿”（纯平移变量）并入即可。当路标数 $\gg$ 位姿数时，瓶颈在构造数据矩阵 $\tilde{\mathbf{Q}}$；此处用\emph{Schur 补技巧}（正是 \cref{eq:schur} 边缘化路标那一招）可使其复杂度\emph{随路标数线性}。有趣的是，\emph{增加}路标数与位姿-路标连通度反而\emph{改善}松弛的精确性（容许更大噪声仍精确）——呼应上文 insight“图连得越好越易精确”。

\paragraph{测距辅助 SLAM：引入单位向量凑成 QCQP。}\cite{carlone2026handbook}
纯距离测量 $\tilde{r}_{ij}$（如 UWB）给出 $(\|\mathbf{t}_j-\mathbf{t}_i\|-\tilde{r}_{ij})^2$ 项，含\emph{未平方}的范数、\emph{非}二次，Shor 松弛不直接适用。一个巧法：引入辅助\emph{单位向量} $\mathbf{b}_{ij}$（$\|\mathbf{b}_{ij}\|=1$），把该项重写成 $\gamma_{ij}\|\mathbf{t}_j-\mathbf{t}_i-\tilde{r}_{ij}\mathbf{b}_{ij}\|^2$——既推断出 $i,j$ 间方位 $\mathbf{b}_{ij}$、又使该项\emph{二次}（单位范数约束也二次），于是整体复成 QCQP，再 Shor 松弛 $+$ Riemannian Staircase 求解（即 CORA 算法）。但测距 SLAM 的松弛\emph{不再}对有界噪声一概精确：多机仅靠机间\emph{测距}（无相对位姿）通常不精确；图连通度再次是松弛精确与否的关键。

\paragraph{各向异性噪声与外点：超出 QCQP，需矩松弛。}\cite{carlone2026handbook,yang2021teaser}
真实传感器噪声常\emph{各向异性}（如双目沿视线方向更不确定）；而 \cref{sec:nlopt-robust} 的外点要用鲁棒核。这两类问题严格说\emph{不再是} QCQP——各向异性使地标项变成\emph{四次}多项式；外点经 B-R 对偶（\cref{thm:nlopt-br}）引入权 $w_{ij}$ 后目标含\emph{三次}项——但都属\emph{多项式优化问题}（POP）。Shor 松弛有个推广叫\emph{矩松弛}（moment / Lasserre 松弛）专治 POP：取单项式向量 $[\mathbf{x}]_r$、作\emph{矩矩阵} $\mathbf{X}_{2r}=[\mathbf{x}]_r[\mathbf{x}]_r^\top$，则任意次数 $\le2r$ 的多项式都是其元素的线性组合，仿 Shor 丢秩约束即得 SDP；它还系统地加\emph{冗余约束}（如矩矩阵重复元相等、$h_i(\mathbf{x})=0\Rightarrow x_k h_i(\mathbf{x})=0$）以提升松弛质量，并给出一个\emph{阶 $r$ 递增}的松弛层级，温和条件下有限阶必精确（SLAM 的各向异性/外点问题常在\emph{低阶}即精确，且用\emph{稀疏}单项式基可缩小 SDP）。\textbf{代价}：矩松弛的大量冗余约束使 SDP \emph{退化}（degenerate），令 Riemannian Staircase \emph{失效}（约束规范条件不成立、对偶解不唯一），需专用且较慢的求解器——这是可证最优在鲁棒/各向异性问题上仍待突破处。点云配准的 TEASER\cite{yang2021teaser}（截断最小二乘 $+$ 半定松弛 $+$ 事后证书）是这条线在前端配准的代表（\cref{ch:vo}）。

\begin{practice}
\begin{enumerate}
  \item （推导）由 \cref{eq:nlopt-qcqp} 经 $\mathbf{x}^\top\mathbf{M}\mathbf{x}=\operatorname{tr}(\mathbf{M}\mathbf{x}\mathbf{x}^\top)$ 推出提升形 \cref{eq:nlopt-qcqp-lifted}，说明“非凸只剩秩-1 约束”。
  \item （概念）解释 $d^*\le p^*$ 为何成立，以及“SDP 解秩 1 $\Rightarrow$ 原 QCQP 全局最优”。由此说明最优性证书 \cref{eq:nlopt-sesync-cert} 的含义。
  \item （概念）SE-Sync 三步里，“解析消去平移得只含旋转的 QCQP \cref{eq:nlopt-pgo-rot}”用到了哪个图矩阵？它与 \cref{sec:nlopt-sparse} 的 Schur 消元有何相似？
  \item （概念）Riemannian Staircase 为何能复用本章 \cref{sec:nlopt-manifold} 的流形优化？\cref{thm:nlopt-staircase} 的“行秩亏二阶临界点”起什么作用？
  \item （概念）为什么测距项 $(\|\mathbf{t}_j-\mathbf{t}_i\|-\tilde r_{ij})^2$ 不是二次、而引入单位向量 $\mathbf{b}_{ij}$ 后就是 QCQP？各向异性/外点为何要用矩松弛而非 Shor？
\end{enumerate}
\end{practice}

% ════════════════════════════════════════════════════════════════
\section{实践：手写 GN / Ceres / g2o}
\label{sec:nlopt-practice}

现代做法是\emph{声明残差、交给库}：库负责线性化、稀疏分解、信赖域调度。本节给同一曲线拟合的三套完整实现并对比（与 \cref{sec:nlopt-gn} 手写 GN 互为印证）。

\paragraph{Ceres（自动求导 $+$ 信赖域）。}\cite{gaoxiang2019slam14}
Ceres 把问题写成 $\min\tfrac12\sum_i\rho_i(\|f_i(\cdot)\|^2)$ s.t. $l_j\le x_j\le u_j$（核 $\rho$ 取恒等即无鲁棒）。用户：定义参数块、残差块（带模板 \texttt{()} 的仿函数 $\to$ 自动求导）、加入 \texttt{Problem}、\texttt{Solve}。\footnote{Ceres 的“自动求导”靠的就是\emph{仿函数（functor）模板化}：残差类的 \texttt{operator()} 写成 \texttt{template<typename T>}，用户用 \texttt{double} 写一遍前向计算、却把雅可比作为\emph{模板类型参数}隐式传入——Ceres 把模板实例化成一个叫 \emph{Jet}（对偶数 dual number，即把每个实数 $a$ 扩成 $a+\sum_i\dot a_i\varepsilon_i$、其中 $\varepsilon_i\varepsilon_j=0$）的类型，像以 \texttt{Jet} 为模板实参调用 \texttt{operator()} 那样跑同一份代码，靠运算符重载在编译期把一阶导一并算出（前向模式：$f(a+\varepsilon)=f(a)+f'(a)\varepsilon$，因 $\varepsilon^2=0$ 高阶项自动归零）\cite{gaoxiang2019slam14,agarwal2012ceres}。本质上是“编译期模板元 $+$ 对偶数”，比有限差分的数值导数精确，又免去手推解析雅可比。}

\begin{codebox}[C++]{Ceres 曲线拟合（slambook2/ch6/ceresCurveFitting.cpp，关键体）}
// 残差仿函数：模板 () 运算符使 Ceres 能自动求导
struct CURVE_FITTING_COST {
  CURVE_FITTING_COST(double x, double y) : _x(x), _y(y) {}
  template<typename T>
  bool operator()(const T* const abc, T* residual) const {  // abc 有 3 维
    residual[0] = T(_y) - ceres::exp(abc[0]*T(_x)*T(_x) + abc[1]*T(_x) + abc[2]);
    return true;                                             // y - exp(a x^2 + b x + c)
  }
  const double _x, _y;
};
// main 中：
double abc[3] = {ae, be, ce};                               // 参数块
ceres::Problem problem;
for (int i=0; i<N; i++)
  problem.AddResidualBlock(
    new ceres::AutoDiffCostFunction<CURVE_FITTING_COST, 1, 3>(  // 误差维 1, 参数维 3
      new CURVE_FITTING_COST(x_data[i], y_data[i])),
    nullptr,    // 核函数: nullptr 即不用鲁棒核(rho=恒等); 换 HuberLoss 即启用 IRLS
    abc);
ceres::Solver::Options options;
options.linear_solver_type = ceres::DENSE_NORMAL_CHOLESKY;  // 稠密正规方程 + Cholesky
// BA 中改为 ceres::SPARSE_SCHUR 即自动对路标做 Schur 边缘化(见 sec sparse)
ceres::Solver::Summary summary;
ceres::Solve(options, &problem, &summary);
// 结果: a,b,c = 0.890908 2.1719 0.943628; 8 次迭代; 约 1.3 ms
\end{codebox}

\noindent Ceres 默认用 LM/信赖域：其日志里 \texttt{tr\_ratio} 即增益比 $\rho$（始终近 1 说明二次近似良好）、\texttt{tr\_radius} 从 $10^4$ 单调升到 $9.4\times10^6$（对应 $\rho>\tfrac34$ 则放大信赖域，\cref{sec:nlopt-lm}）。注意 Ceres 报的 cost 是 $\tfrac12\sum e^2\approx50.97$，手写 GN 报 $\sum e^2\approx101.94$，差因子 2（系数 $\tfrac12$）。

\paragraph{g2o（图优化）。}\cite{gaoxiang2019slam14}
g2o（general graphic optimization）把变量建成\emph{顶点}、残差建成\emph{边}：顶点用\emph{超图}（边可连一/二/多个顶点）。曲线拟合只有一个顶点 $(a,b,c)$、100 条\emph{一元边}（自连）。顶点须实现\emph{更新} \texttt{oplusImpl}（位姿则用右扰动）、边须实现 \texttt{computeError} 与 \texttt{linearizeOplus}。

\begin{codebox}[C++]{g2o 曲线拟合（slambook2/ch6/g2oCurveFitting.cpp，关键体）}
// 顶点: 优化变量维 3, 类型 Vector3d
class CurveFittingVertex : public g2o::BaseVertex<3, Eigen::Vector3d> {
public:
  EIGEN_MAKE_ALIGNED_OPERATOR_NEW
  virtual void setToOriginImpl() override { _estimate << 0,0,0; }
  virtual void oplusImpl(const double* update) override {     // 向量空间: 直接加
    _estimate += Eigen::Vector3d(update);                     // 位姿则改为右乘 Exp
  }
  virtual bool read(istream&) { return true; }
  virtual bool write(ostream&) const { return true; }
};
// 边: 观测维 1, 类型 double, 连接 CurveFittingVertex
class CurveFittingEdge : public g2o::BaseUnaryEdge<1, double, CurveFittingVertex> {
public:
  EIGEN_MAKE_ALIGNED_OPERATOR_NEW
  CurveFittingEdge(double x) : BaseUnaryEdge(), _x(x) {}
  virtual void computeError() override {
    const auto* v = static_cast<const CurveFittingVertex*>(_vertices[0]);
    const Eigen::Vector3d abc = v->estimate();
    _error(0,0) = _measurement - std::exp(abc[0]*_x*_x + abc[1]*_x + abc[2]);
  }
  virtual void linearizeOplus() override {                    // 解析雅可比
    const auto* v = static_cast<const CurveFittingVertex*>(_vertices[0]);
    const Eigen::Vector3d abc = v->estimate();
    double y = exp(abc[0]*_x*_x + abc[1]*_x + abc[2]);
    _jacobianOplusXi[0] = -_x*_x*y;   // de/da
    _jacobianOplusXi[1] = -_x*y;      // de/db
    _jacobianOplusXi[2] = -y;         // de/dc
  }
  virtual bool read(istream&) { return true; }
  virtual bool write(ostream&) const { return true; }
  double _x;
};
// main: 选 GN/LM/DogLeg 之一; 加顶点; 加边(setInformation = Sigma^-1); optimize(10)
// 边可挂鲁棒核: edge->setRobustKernel(new g2o::RobustKernelHuber());  // IRLS + Huber
// 结果: estimated model = 0.890912 2.1719 0.943629
\end{codebox}

\noindent g2o 顶点的 \texttt{oplusImpl} 处理 $\mathbf{x}_{k+1}=\mathbf{x}_k\boxplus\Delta\mathbf{x}$：曲线参数在向量空间故直接加；\emph{位姿不在向量空间}，须用右乘更新 $\mathbf{T}\,\mathrm{Exp}(\delta\boldsymbol{\xi})$（\cref{sec:nlopt-manifold}，BA 中的关键）。日志 \texttt{chi2} 即 $\sum e^\top\boldsymbol{\Sigma}^{-1}e\approx101.94$；\texttt{schur=0} 表示本例无需 Schur（只一个顶点）。顶点与边里的 \texttt{read}/\texttt{write} 之所以可\emph{留空}（直接 \texttt{return true}）：它们只服务于把图\emph{存盘/读盘}（如 \texttt{.g2o} 文本格式）的磁盘序列化，与\emph{内存中}的优化迭代无关——只要不需要把问题落地成文件，把它们留空不影响求解\cite{gaoxiang2019slam14}。

\paragraph{三者对比（同一曲线拟合）。}\cite{gaoxiang2019slam14}

\begin{table}[H]\centering\small
\caption{手写 GN / Ceres / g2o 对比（曲线拟合，真值 $a,b,c=1,2,1$）}
\label{tab:nlopt-three}
\begin{tabular}{lllll}
\toprule
方法 & 求解器 / 求导 & 迭代 & 用时 & 终估 $(a,b,c)$ \\
\midrule
手写 GN & LDLT(Cholesky) / 解析雅可比 & 9 & $\approx0.2$ ms & $0.8909,2.1719,0.9436$ \\
Ceres & DENSE\_NORMAL\_CHOLESKY / 自动求导 $+$ LM & 8 & $\approx1.3$ ms & $0.8909,2.1719,0.9436$ \\
g2o & LinearSolverDense / 解析雅可比 $+$ GN & 10 & $\approx0.9$ ms & $0.8909,2.1719,0.9436$ \\
\bottomrule
\end{tabular}
\end{table}

\noindent 结果一致，速度\emph{手写 $>$ g2o $>$ Ceres}（“通用性与效率互相矛盾”；Ceres 此处用了自动求导且配置非纯 GN，故看起来慢）。

\paragraph{解析雅可比 vs 自动求导。}\cite{gaoxiang2019slam14}
能推\emph{解析}雅可比（用 \cref{ch:lie} 右扰动）就告诉库，避免数值/自动求导的精度与效率损失——这与本书“逐公式推导”的基调一致。Ceres 提供模板元自动求导（编译期、基于对偶数 Jet 的精确前向自动微分，非数值差分）与运行时数值求导；g2o 提供运行时数值求导。多数问题若能推出解析形式，可避免数值求导的种种问题。BA 中位姿用流形参数化、线性求解器选 \texttt{SPARSE\_SCHUR} 即自动对路标 Schur 边缘化（\cref{sec:nlopt-sparse}）。

\begin{practice}
\begin{enumerate}
  \item 用 Sophus $+$ Eigen 手写一个 SE(3) 位姿图的 GN/LM 迭代，\texttt{oplusImpl} 用右乘，对比有无 Huber 核时对外点的鲁棒性。
  \item 在 Ceres 里把线性求解器从 \texttt{DENSE} 换成 \texttt{SPARSE\_SCHUR}，在一个 BA 例上测时间，验证 Schur 的加速。
  \item （概念）解释 g2o 的 \texttt{setInformation(Sigma.inverse())} 与本章白化/信息矩阵加权的对应。
\end{enumerate}
\end{practice}

% ════════════════════════════════════════════════════════════════
\section*{本章常见误解汇总}
\addcontentsline{toc}{section}{本章常见误解汇总}
\begin{table}[H]\centering\small
\begin{tabular}{p{0.46\textwidth} p{0.46\textwidth}}
\toprule
\textbf{常见误解} & \textbf{正确理解} \\
\midrule
非线性最小二乘能像线性那样一步解 & 一般无闭式解，须迭代线性化（\cref{sec:nlopt-why}）\\
MAP 给的是后验“均值” & 给的是“众数”，非线性时 $\ne$ 均值，故 ML 有偏（\cref{sec:nlopt-cov}）\\
信息矩阵加权可有可无 & 不加则量纲混乱、丢各向异性；须白化（\cref{sec:nlopt-solve}）\\
高斯牛顿就是牛顿法 & GN 线性化\emph{残差}、用 $\mathbf{J}^\top\mathbf{W}^{-1}\mathbf{J}$ 近似海森，丢残差$\times$二阶导项（\cref{thm:gn}）\\
LM 的 $\lambda$ 越大越好 & 大趋最速下降（稳但慢）、小趋 GN（快但可能不稳）；按 $\rho$ 自适应（\cref{sec:nlopt-lm}）\\
增益比 $\rho$ 与鲁棒核 $\rho$ 是一回事 & 前者是标量（调信赖域），后者是函数（降权外点），仅撞名 \\
直接对信息矩阵求逆解方程 & 用 Cholesky/QR 分解 $+$ 前后代，绝不显式求逆（\cref{sec:nlopt-solve}）\\
Cholesky 和 QR 给不同的解 & 给相同 $\mathbf{R}$、相同解；QR 更稳、Cholesky 更快（\cref{sec:nlopt-solve}）\\
变量排序会改变解 & 只改 fill-in（速度），解完全相同（COLAMD，\cref{sec:nlopt-sparse}）\\
BA 上万维解不动 & 利用因子图稀疏 $+$ Schur 消元，降到 $O(\text{相机}^3)$（\cref{eq:schur}）\\
g2o BA 不必声明边缘化顶点 & 用稀疏 Schur 时\emph{必须}对路标 \texttt{setMarginalized}，否则报错（\cref{sec:nlopt-ba}）\\
边缘化路标和边缘化关键帧后果一样 & 删路标 fill-in 落相机块（无害）；删关键帧落路标块（破坏对角，\cref{sec:nlopt-ba}）\\
位姿图不固定基准也能解 & 只含相对约束 $\Rightarrow$ gauge 自由度 $\Rightarrow$ $\mathbf{H}$ 奇异；须固定一帧（\cref{sec:nlopt-posegraph}）\\
位姿当普通向量做加法更新 & 用右扰动 $\boxplus$（$\mathbf{T}\,\mathrm{Exp}(\delta\boldsymbol{\xi})$），无约束（\cref{sec:nlopt-manifold}）\\
鲁棒核是工程黑魔法 & 它 $=$ IRLS 重加权 $=$ 逆 Wishart 先验下的 MAP $=$ B-R 对偶一实例（\cref{eq:nlopt-robust-cauchy},\cref{thm:nlopt-br}）\\
前端剔除干净就不必后端鲁棒核 & 几何/一致性约束是内点\emph{必要非充分}条件，残余外点仍毒化二次代价（\cref{sec:nlopt-robust-frontend}）\\
RANSAC 万能 & 外点多/最小集大时期望迭代 $1/w^n$ 爆炸；位姿图宜用 PCM 最大团（\cref{eq:nlopt-ransac},\cref{eq:nlopt-pcm}）\\
GN 逆海森就是真协方差 & 是 Laplace$+$GN 双重近似，强非线性下过乐观（\cref{eq:laplace}）\\
GVI 只是迭代更多的 UKF & GVI 有明确目标（负 ELBO）、按\emph{后验}撒点、联合优均值与协方差，几乎无偏（\cref{sec:nlopt-gvi}）\\
MAP$+$Laplace 与 GVI 无关 & MAP$+$Laplace $=$ GVI 只用均值点算期望的退化（\cref{eq:gvi-stein}）\\
可微 BA 就是把 BA 当网络一层 & 那只单向反传、监督来自真值；真双层优化让前端从几何后端学（\cref{eq:nlopt-blo}）\\
GN/LM 给的就是全局最优 & 只局部最优、且不自知；半定松弛/SE-Sync 可证全局最优$+$证书（\cref{sec:nlopt-certifiable}）\\
\bottomrule
\end{tabular}
\end{table}

% ════════════════════════════════════════════════════════════════
\section*{本章小结}
\addcontentsline{toc}{section}{本章小结}

\begin{table}[H]\centering\small
\caption{符号表}
\begin{tabular}{lll}
\toprule
符号 & 含义 & 首次出现 \\
\midrule
$J(\mathbf{x})$ & 目标函数 $\tfrac12\sum_i\|\mathbf{e}_i\|_{\boldsymbol{\Sigma}_i}^2$ & \cref{eq:nls} \\
$\mathbf{e}(\mathbf{x}),\mathbf{J}$ & 残差及其雅可比 $\mathbf{J}=\partial\mathbf{e}/\partial\mathbf{x}$ & \cref{eq:gn-linearize} \\
$\nabla J,\mathbf{H}_J$ & 目标的梯度 / 海森（区别于残差雅可比 $\mathbf{J}$） & \cref{eq:taylor-J} \\
$\boldsymbol{\Lambda},\mathbf{W}$ & 信息矩阵 $\mathbf{J}^\top\mathbf{W}^{-1}\mathbf{J}$ / 堆叠权 $\operatorname{diag}(\boldsymbol{\Sigma}_i)$ & \cref{eq:gn-normal} \\
$\lambda,\rho$ & LM 阻尼 / 增益比 & \cref{eq:lm},\cref{eq:lm-rho} \\
$\mathbf{A}_i,\mathbf{b}_i$ & 白化雅可比 / 右端 $\boldsymbol{\Sigma}_i^{-1/2}\mathbf{J}_i,\ \boldsymbol{\Sigma}_i^{-1/2}(-\mathbf{e}_i)$ & \cref{eq:nlopt-Ai} \\
$\mathbf{R}$（此章）& Cholesky 上三角因子 $\boldsymbol{\Lambda}=\mathbf{R}^\top\mathbf{R}$（非旋转） & \cref{eq:cholesky} \\
$\mathbf{B},\mathbf{C},\mathbf{E},\mathbf{S}$ & BA 相机/路标/耦合/Schur 补 $\mathbf{S}=\mathbf{B}-\mathbf{E}\mathbf{C}^{-1}\mathbf{E}^\top$ & \cref{eq:ba-block} \\
$\mathbf{s}_j,\tau$ & 变量消元的分隔集 / 新因子 & \cref{eq:nlopt-elim-cond} \\
$\boxplus,\boxminus$ & 流形增量 / 差 $\mathbf{T}\boxplus\delta\boldsymbol{\xi}=\mathbf{T}\,\mathrm{Exp}(\delta\boldsymbol{\xi})$ & \cref{eq:retraction} \\
$\rho(u),u_i$ & 鲁棒核 / 标量马氏范数 & \cref{eq:nlopt-mestim} \\
$\mathbf{Y}_i$ & IRLS 重加权后的等效协方差 & \cref{eq:irls} \\
$\mathbf{M}_i,\boldsymbol{\Psi}_i,\nu_i$ & 待估协方差 / 逆 Wishart scale / 自由度 & \cref{eq:nlopt-invwishart} \\
$\check{\mathbf{x}},\hat{\mathbf{x}},\mathbf{x}_{\mathrm{op}}$ & 先验 / 后验 / 线性化工作点 & \cref{sec:nlopt-cov} \\
$\epsilon_{\text{nees}},\epsilon_{\text{nis}}$ & 归一化估计误差/创新平方 & \cref{eq:nlopt-nees},\cref{eq:nlopt-nis} \\
$q(\mathbf{x}),V(q),\phi$ & 高斯近似后验 / 负 ELBO 损失 / 负对数联合 $-\ln p(\mathbf{x},\mathbf{z})$ & \cref{eq:gvi-V} \\
$q_k,\boldsymbol{\Sigma}_{kk},\mathbf{P}_k$ & 第 $k$ 因子的边缘 / 其协方差块 / 抽取投影 & \cref{eq:gvi-esgvi-hess} \\
$C,F,\mathsf{S},\gamma$ & 几何约束 / 一致性函数 / 共识(内点)集 / 松弛阈值 & \cref{eq:nlopt-geomconstraint},\cref{eq:nlopt-pcm} \\
$w_i,\Phi_\rho(w_i),\beta$ & B-R 对偶权 / 外点过程 / 截断阈值 & \cref{eq:nlopt-br},\cref{eq:nlopt-truncquad} \\
$\rho_\mu,\mu$ & GNC 光滑核 / 控制凸性的光滑因子 & \cref{eq:nlopt-gnc-tq} \\
$\mathbf{y},\mathcal{U},\mathcal{L}$ & 上层参数(网络) / 上层损失 / 下层 NLS 损失 & \cref{eq:nlopt-blo} \\
$\mathbf{H}_{xx},\mathbf{H}_{xy},\mathbf{q}$ & 下层海森 / 交叉二阶导 / HVP 求解的辅助向量 & \cref{eq:nlopt-implicit-grad},\cref{eq:nlopt-hvp} \\
$\mathbf{X},\mathbf{Z}$（此节）& Shor/SE-Sync 提升的半正定矩阵变量（非状态） & \cref{eq:nlopt-sdp},\cref{eq:nlopt-pgo-sdp} \\
$\tilde{\mathbf{Q}},\mathrm{St}(d,r)$ & PGO 数据矩阵(非过程噪声) / Stiefel 流形 & \cref{eq:nlopt-pgo-rot},\cref{eq:nlopt-pgo-bm} \\
$p^*,d^*,p^*_{\mathrm{SDP}}$ & QCQP 最优 / SDP 最优(下界) / 松弛最优值 & \cref{eq:nlopt-sdp},\cref{eq:nlopt-pgo-sdp} \\
\bottomrule
\end{tabular}
\end{table}

\begin{table}[H]\centering\small
\caption{定理速查表}
\begin{tabular}{lll}
\toprule
名称 & 一句话说明（你能做什么） & 出处 \\
\midrule
MAP$\to$最小二乘 & 把“找最可能状态”化成可解的加权二次型 & \cref{thm:nlopt-map} \\
迭代下降 & 非线性问题靠“线性化$\to$解增量$\to$更新”反复逼近 & \cref{sec:nlopt-why} \\
牛顿 / GN & GN 用 $\mathbf{J}^\top\mathbf{W}^{-1}\mathbf{J}$ 替海森，免算二阶导 & \cref{thm:gn} \\
LM & 用 $\lambda,\rho$ 在 GN 与最速下降间安全插值 & \cref{prop:lm} \\
白化 & 左乘 $\boldsymbol{\Sigma}^{-1/2}$ 使各观测同量纲可加 & \cref{eq:nlopt-whiten} \\
Cholesky/QR/$\sqrt{}$SAM & 分解 $+$ 前后代解线性系统，绝不求逆 & \cref{eq:cholesky} \\
Schur 消元 & 边缘化路标，复杂度降到 $O(\text{相机}^3)$ & \cref{eq:schur} \\
变量消元 $=$ 分解 & 消元一步$=$Schur 补，回代得 MAP 与协方差 & \cref{eq:nlopt-elim-cond} \\
BA 代价与稀疏 & 重投影误差 $+$ 箭头 $\mathbf{H}=[\mathbf{B},\mathbf{E};\mathbf{E}^\top,\mathbf{C}]$ & \cref{eq:nlopt-ba-cost} \\
位姿图 & 去路标，只优化位姿、边为相对运动 & \cref{eq:nlopt-pg-cost} \\
流形 GN & 右扰动雅可比 $+$ $\boxplus$ 更新，把约束优化变无约束 & \cref{eq:retraction} \\
前端剔除 RANSAC/PCM & 几何/一致性约束选最大内点集；PCM$=$最大团 & \cref{eq:nlopt-consensusmax},\cref{eq:nlopt-pcm} \\
IRLS & 把鲁棒核变回“逐步重加权”的普通最小二乘 & \cref{eq:irls} \\
鲁棒 $=$ MAP & 鲁棒核 $\Leftrightarrow$ 逆 Wishart 协方差先验下的 MAP & \cref{eq:nlopt-robust-cauchy} \\
（选读）Black--Rangarajan 对偶 & 鲁棒核 $\Leftrightarrow$ 权 $+$ 外点过程；IRLS$=$交替最小化 & \cref{thm:nlopt-br} \\
（选读）GNC & 从凸光滑核渐进收紧到原核，抗坏初值 & \cref{alg:nlopt-gnc} \\
拉普拉斯协方差 & 收敛后逆海森即后验协方差（$=$ CRLB） & \cref{eq:laplace} \\
ML 偏差 & 非线性使估计有偏，可按曲率闭式去偏 & \cref{eq:nlopt-mlbias} \\
NEES/NIS & 用卡方检验判断协方差是否过/欠自信 & \cref{eq:nlopt-nees} \\
（选读）GVI（负 ELBO） & 估\emph{整个}最优高斯而非众数；MAP$+$Laplace 是其单点退化 & \cref{eq:gvi-V} \\
（选读）ESGVI 稀疏 & 似然因式分解 $\Rightarrow\boldsymbol{\Sigma}^{-1}$ 精确稀疏，复用 BA 后端 & \cref{eq:gvi-esgvi-hess} \\
（选读）IEKF $=$ 单步 MAP & IEKF 不动点 $=$ 后验众数（GN 证） & \cref{thm:nlopt-iekf-map} \\
（选读）双层优化/隐式微分 & 对优化解 $\mathbf{x}^*$ 求导 $\partial\mathbf{x}^*/\partial\mathbf{y}=-\mathbf{H}_{xx}^{-1}\mathbf{H}_{xy}$，配 HVP & \cref{eq:nlopt-implicit-grad} \\
（选读）Shor 松弛 & QCQP 提升丢秩约束$\to$凸 SDP；秩 1 即全局最优 & \cref{eq:nlopt-sdp} \\
（选读）SE-Sync & PGO$\to$旋转 QCQP$\to$SDP，松弛精确则全局最优$+$证书 & \cref{thm:nlopt-sesync-exact} \\
（选读）Riemannian Staircase & Burer--Monteiro 低秩$+$Stiefel 优化；秩亏临界点$=$全局最优 & \cref{thm:nlopt-staircase} \\
\bottomrule
\end{tabular}
\end{table}

\begin{table}[H]\centering\small
\caption{知识点总表}
\begin{tabular}{cllll}
\toprule
\# & 知识点 & 核心要点 & 难度 & 脉络层 \\
\midrule
1 & MAP$\to$最小二乘 & 高斯$+$独立$+$信息加权 & \(\bigstar\bigstar\bigstar\) & 骨干 \\
2 & 迭代下降 & 线性化$\to$解增量$\to$更新 & \(\bigstar\bigstar\) & 骨干 \\
3 & 梯度/牛顿/GN & $\mathbf{J}^\top\mathbf{W}^{-1}\mathbf{J}$ 近似海森 & \(\bigstar\bigstar\bigstar\) & 骨干 \\
4 & LM/信赖域 & $\lambda,\rho$ 插值与接受/拒绝 & \(\bigstar\bigstar\bigstar\) & 骨干 \\
5 & 白化·Cholesky/QR & $\sqrt{}$SAM；不求逆 & \(\bigstar\bigstar\bigstar\) & 骨干 \\
6 & 稀疏·Schur·消元 & fill-in$=$共视；COLAMD & \(\bigstar\bigstar\bigstar\bigstar\) & 骨干 \\
6b & BA·重投影·箭头 $\mathbf{H}$ & $\mathbf{e}=\mathbf{z}-h(\mathbf{T},\mathbf{p})$；Ceres/g2o BA & \(\bigstar\bigstar\bigstar\bigstar\) & 骨干 \\
6c & 位姿图·滑窗边缘化 & 去路标；删帧致 fill-in & \(\bigstar\bigstar\bigstar\) & 次优 \\
7 & 流形右扰动 & $\boxplus$；对接 Ceres/GTSAM & \(\bigstar\bigstar\bigstar\bigstar\) & 次优 \\
8 & 前端剔除 RANSAC/PCM & 共识最大化；一致性图$=$最大团；必要非充分 & \(\bigstar\bigstar\bigstar\bigstar\) & 前沿 \\
8b & 后端鲁棒核/IRLS/逆 Wishart & 鲁棒$=$协方差先验 MAP；B-R 对偶；GNC 抗坏初值 & \(\bigstar\bigstar\bigstar\bigstar\) & 前沿 \\
9 & 协方差/偏差/一致性 & Laplace$=$CRLB；ML 偏差；NEES & \(\bigstar\bigstar\bigstar\bigstar\) & 次优 \\
9b & GVI/ESGVI（选读） & 负 ELBO；后验线性化；接因子图；立体相机偏差 0.28cm & \(\bigstar\bigstar\bigstar\bigstar\bigstar\) & 前沿 \\
9c & 可微优化（选读） & 双层优化；隐式微分 $\partial\mathbf{x}^*/\partial\mathbf{y}$；海森-向量积 & \(\bigstar\bigstar\bigstar\bigstar\bigstar\) & 前沿 \\
9d & 可证最优（选读） & Shor 松弛；SE-Sync；Riemannian Staircase；最优性证书 & \(\bigstar\bigstar\bigstar\bigstar\bigstar\) & 前沿 \\
10 & 三库实践 & 手写$>$g2o$>$Ceres & \(\bigstar\bigstar\) & 次优 \\
\bottomrule
\end{tabular}
\end{table}

\paragraph{延伸阅读。}
\begin{itemize}
  \item 视觉 SLAM 十四讲\cite{gaoxiang2019slam14} 第 6、9、10 章：MAP$\to$最小二乘、GN/LM 直觉与代码、BA 重投影代价与稀疏 Schur（\cref{sec:nlopt-ba}）、Ceres/g2o 曲线拟合与 BAL BA、位姿图与“球”实践（\cref{sec:nlopt-posegraph}）、滑动窗口边缘化。
  \item Barfoot\cite{barfoot2024state} 第 4、5 章：批量 GN 的严谨推导（$\mathbf{J}^\top\mathbf{W}^{-1}\mathbf{J}$ 近似海森、IEKF$=$GN、Laplace 协方差、ML 偏差、IRLS、鲁棒$=$逆 Wishart MAP、SWF 边缘化、NEES/NIS）；第 5 章变分推断（GVI/ESGVI、负 ELBO、Takahashi 稀疏协方差回收、cubature 无导数推断、参数估计/系统辨识、立体相机 GVI vs MAP vs ISPKF，\cref{sec:nlopt-gvi}）。原始论文见 ESGVI\cite{barfoot2020esgvi}、变分高斯近似\cite{opper2009variational}、自然梯度\cite{amari1998natural}。
  \item SLAM Handbook\cite{carlone2026handbook} \& Dellaert/Kaess\cite{dellaert2017factor} 第 1 章：因子图、白化、Cholesky/QR/$\sqrt{}$SAM、稀疏$=$因子图、COLAMD/fill-in、变量消元$=$分解、Bayes 树/iSAM2。
  \item SLAM Handbook\cite{carlone2026handbook} 第 3 章（鲁棒）：外点成因（数据关联/感知混叠/模型违背）、前端 RANSAC 与成对一致性最大化（PCM，\cref{sec:nlopt-robust-frontend}）、M-估计核族、Black--Rangarajan 对偶（\cref{thm:nlopt-br}）、GNC（\cref{alg:nlopt-gnc}）、M3500/SubT/Victoria Park 工程直觉。原始文献：RANSAC\cite{fischler1981ransac}、Nistér 五点法\cite{nister2004fivepoint}、PCM\cite{mangelson2018pcm}、B-R 对偶\cite{black1996rangarajan}、GNC\cite{yang2020gnc}、自适应核族\cite{barron2019general}、开关约束\cite{sunderhauf2012switchable}。
  \item SLAM Handbook\cite{carlone2026handbook} 第 4 章（可微优化，\cref{sec:nlopt-diffopt}）：双层优化、展开/隐式微分、海森-向量积、FastTriggs、Theseus/PyPose 等库。
  \item SLAM Handbook\cite{carlone2026handbook} 第 6 章（可证最优，\cref{sec:nlopt-certifiable}）：Shor 半定松弛、SE-Sync\cite{rosen2019sesync}（旋转 QCQP$\to$SDP$\to$精确性$\to$Riemannian Staircase$\to$证书）、地标/测距/各向异性/外点扩展与矩松弛；点云配准的 TEASER\cite{yang2021teaser}。
  \item Hertzberg 等\cite{hertzberg2013integrating}：$\boxplus/\boxminus$ 流形封装——把通用优化干净地搬到流形。
\end{itemize}

\section*{本章与后续章节的关系}
\begin{table}[H]\centering\small
\begin{tabular}{lll}
\toprule
后续章节 & 与本章的关系 & 本章哪个知识点为其铺垫 \\
\midrule
\cref{ch:vo} 视觉里程计 & 重投影 BA 后端解最小二乘 & BA 重投影/箭头稀疏（\cref{sec:nlopt-ba}）、流形 GN、Schur、鲁棒核 \\
\cref{ch:lidar_slam} 激光 SLAM & 点云配准/位姿图优化 & 位姿图（\cref{sec:nlopt-posegraph}）、GN/LM、白化、Cholesky \\
\cref{ch:vio} VIO & IMU 预积分因子 $+$ 滑窗边缘化 & SWF 边缘化（\cref{sec:nlopt-ba} 删帧 fill-in $+$ 附录）、流形右扰动 \\
\cref{ch:calib} 相机标定 & 最大似然 LM 精化内外参 & GN/LM、解析雅可比、协方差 \\
\bottomrule
\end{tabular}
\end{table}

\section*{故障排查手册}
\addcontentsline{toc}{section}{故障排查手册}
\begin{table}[H]\centering\small
\begin{tabular}{p{0.22\textwidth} p{0.30\textwidth} p{0.40\textwidth}}
\toprule
症状 & 可能原因 & 排查步骤 / 对策 \\
\midrule
不收敛 / 发散 & 初值差、强非线性、步太大 & 给好初值（ICP/PnP）；改 LM 增大 $\lambda$；线搜索取 $\alpha<1$（\cref{sec:nlopt-lm}）\\
$\boldsymbol{\Lambda}$ 奇异 / 解爆掉 \texttt{NaN} & 不可观测、gauge 自由度、$\theta\to0$ 雅可比奇异 & 加先验/固定基准帧（\cref{ch:slam_est}）；小角度用泰勒分支 \\
BA 慢得无法实时 & 未用稀疏 Schur / 排序差 & 选 \texttt{SPARSE\_SCHUR}；好排序（COLAMD，\cref{sec:nlopt-sparse}）\\
回环处轨迹被拉飞 & 外点/误匹配主导二次代价 & 前端 RANSAC/PCM 预剔（\cref{sec:nlopt-robust-frontend}）$+$ 后端鲁棒核（Huber/Cauchy）；非凸核用 GNC（\cref{alg:nlopt-gnc}）\\
GNC 把内点也判成外点而崩 & 初值差（如里程计也含外点）& 用不依赖初值的前端 PCM 先剔、再 GNC（PCM$+$GNC 最稳，\cref{sec:nlopt-robust-gnc}）\\
解收敛却疑似局部极小、无从判断 & 非凸、初值差、局部法不自知 & 用 SE-Sync 等可证最优求解器（半定松弛$+$最优性证书，\cref{sec:nlopt-certifiable}）\\
位姿更新后不再是 SE(3) & 用了普通加法更新 & 右扰动 $\boxplus$（$\mathbf{T}\,\mathrm{Exp}(\delta\boldsymbol{\xi})$，\cref{sec:nlopt-manifold}）\\
雅可比侧与更新侧不符、震荡 & 左扰动雅可比配右乘更新 & 全程同侧；引用左扰动公式经 $\mathrm{Ad}$ 转右（\cref{ch:lie}）\\
协方差过乐观（NEES$>N$）& GN/Laplace 近似 $+$ ML 偏差 & NEES/NIS 检验；强非线性处去偏或采样（\cref{sec:nlopt-cov}）\\
\bottomrule
\end{tabular}
\end{table}

\section*{研究实践建议}
\addcontentsline{toc}{section}{研究实践建议}
给新手：先把手写 GN 曲线拟合（\cref{sec:nlopt-gn}）跑通，再换 Ceres/g2o 验证；调 BA 时先确认\emph{初值}与\emph{稀疏求解器}选对，再谈鲁棒核。给有经验者：能推解析雅可比就别用自动求导（\cref{ch:lie} 右扰动）；用 NEES/NIS 把协方差校准好（下游一致性依赖它）；非凸鲁棒核配 GNC 以避局部极小。

% ════════════════════════════════════════════════════════════════
\section*{附录：本章推导补遗}
\label{sec:nlopt-appendix}
\addcontentsline{toc}{section}{附录：本章推导补遗}

\begin{derivation}[高斯牛顿如何近似海森（\cref{thm:gn}）]
$J=\tfrac12\mathbf{e}^\top\mathbf{W}^{-1}\mathbf{e}$，梯度 $\nabla J=\mathbf{J}^\top\mathbf{W}^{-1}\mathbf{e}$。再求一次导：
\[
  \mathbf{H}_J=\frac{\partial}{\partial\mathbf{x}}\big(\mathbf{J}^\top\mathbf{W}^{-1}\mathbf{e}\big)
  =\mathbf{J}^\top\mathbf{W}^{-1}\mathbf{J}+\sum_i\big[\mathbf{W}^{-1}\mathbf{e}\big]_i\,\frac{\partial^2 e_i}{\partial\mathbf{x}\partial\mathbf{x}^\top}.
\]
第一项只含一阶导 $\mathbf{J}$；第二项含残差 $\mathbf{e}$ 与二阶导之积。接近最优时 $\mathbf{e}\to\mathbf{0}$，第二项被压制，故 $\mathbf{H}_J\approx\mathbf{J}^\top\mathbf{W}^{-1}\mathbf{J}$。代入牛顿步 $\mathbf{H}_J\Delta\mathbf{x}=-\nabla J$ 即 \cref{eq:gn-normal}。当残差大且模型强非线性时第二项不可略，GN 退化、需牛顿或 LM。\hfill$\square$
\end{derivation}

\begin{derivation}[Schur 消元（\cref{eq:schur}）]
对 \cref{eq:ba-block} 左乘消元矩阵 $\big[\begin{smallmatrix}\mathbf{I}&-\mathbf{E}\mathbf{C}^{-1}\\\mathbf{0}&\mathbf{I}\end{smallmatrix}\big]$：
\[
  \begin{bmatrix}\mathbf{B}-\mathbf{E}\mathbf{C}^{-1}\mathbf{E}^\top&\mathbf{0}\\\mathbf{E}^\top&\mathbf{C}\end{bmatrix}
  \begin{bmatrix}\Delta\mathbf{x}_c\\\Delta\mathbf{x}_p\end{bmatrix}
  =\begin{bmatrix}\mathbf{v}-\mathbf{E}\mathbf{C}^{-1}\mathbf{w}\\\mathbf{w}\end{bmatrix}.
\]
第一行只含 $\Delta\mathbf{x}_c$，即 $\mathbf{S}\Delta\mathbf{x}_c=\mathbf{v}-\mathbf{E}\mathbf{C}^{-1}\mathbf{w}$（$\mathbf{S}$ 为 $\mathbf{C}$ 的 Schur 补）；解出后回代第二行 $\mathbf{C}\Delta\mathbf{x}_p=\mathbf{w}-\mathbf{E}^\top\Delta\mathbf{x}_c$ 得 $\Delta\mathbf{x}_p$。因 $\mathbf{C}$ 块对角，$\mathbf{C}^{-1}$ 廉价；主成本在分解 $\mathbf{S}$（仅相机维）。$\mathbf{S}$ 的非对角填充即共视诱导的 fill-in。\hfill$\square$
\end{derivation}

\begin{derivation}[IRLS 的合法性（\cref{eq:irls}）]
M-估计目标 $J'=\sum_i\alpha_i\rho(u_i)$，$u_i=\|\mathbf{e}_i\|_{\boldsymbol{\Sigma}_i}$。链式求导
$\partial J'/\partial\mathbf{x}=\sum_i\alpha_i\dfrac{\rho'(u_i)}{u_i}\mathbf{e}_i^\top\boldsymbol{\Sigma}_i^{-1}\dfrac{\partial\mathbf{e}_i}{\partial\mathbf{x}}$。
令重加权目标 $J''=\tfrac12\sum_i\mathbf{e}_i^\top\mathbf{Y}_i^{-1}\mathbf{e}_i$，$\mathbf{Y}_i^{-1}=\dfrac{\alpha_i\rho'(u_i)}{u_i}\boldsymbol{\Sigma}_i^{-1}$，在冻结点 $\mathbf{x}_{\mathrm{op}}$ 处其梯度与 $J'$ 相同。故 $\partial J'/\partial\mathbf{x}=\partial J''/\partial\mathbf{x}=\mathbf{0}$ 同解：每步解普通加权最小二乘 $J''$、更新权 $\mathbf{Y}_i$、迭代，即 IRLS，收敛到 $J'$ 的极小（到达路径不同）。\hfill$\square$
\end{derivation}

\begin{derivation}[鲁棒 $=$ 逆 Wishart 先验 MAP（\cref{eq:nlopt-robust-cauchy}）]
把最优协方差 \cref{eq:nlopt-invwishart-M} $\mathbf{M}_i(\mathbf{x})=\tfrac1{\alpha_i}(\boldsymbol{\Psi}_i+\mathbf{e}_i\mathbf{e}_i^\top)$ 代回 \cref{eq:nlopt-invwishart-J}。逐项：
$\mathbf{e}_i^\top\mathbf{M}_i^{-1}\mathbf{e}_i$ 用 Sherman--Morrison 得 $\dfrac{\alpha_i\,\mathbf{e}_i^\top\boldsymbol{\Psi}_i^{-1}\mathbf{e}_i}{1+\mathbf{e}_i^\top\boldsymbol{\Psi}_i^{-1}\mathbf{e}_i}$；$-\alpha_i\ln\det(\mathbf{M}_i^{-1})=\alpha_i\ln\det(\mathbf{M}_i)$ 含 $\alpha_i\ln(1+\mathbf{e}_i^\top\boldsymbol{\Psi}_i^{-1}\mathbf{e}_i)$（用 $\det(\boldsymbol{\Psi}+\mathbf{e}\mathbf{e}^\top)=\det\boldsymbol{\Psi}(1+\mathbf{e}^\top\boldsymbol{\Psi}^{-1}\mathbf{e})$）。合并、丢与 $\mathbf{x}$ 无关常数，得 $J'(\mathbf{x})=\tfrac12\sum_i\alpha_i\ln(1+\mathbf{e}_i^\top\boldsymbol{\Psi}_i^{-1}\mathbf{e}_i)$，即加权 Cauchy 核（$\boldsymbol{\Psi}_i=\boldsymbol{\Sigma}_i$）。\hfill$\square$
\end{derivation}

\begin{derivation}[（选读）ML 偏差闭式（\cref{eq:nlopt-mlbias}，Box 1971）]
最优条件 $\sum_k\mathbf{G}_k(\hat{\mathbf{x}})^\top\boldsymbol{\Sigma}_{v,k}^{-1}(\mathbf{y}_k-\mathbf{g}_k(\hat{\mathbf{x}}))=\mathbf{0}$。在真值 $\mathbf{x}$ 处对 $\mathbf{g}_k,\mathbf{G}_k$ 二阶/一阶展开（$\delta\mathbf{x}=\hat{\mathbf{x}}-\mathbf{x}$，$\mathbf{n}_k=\mathbf{y}_k-\mathbf{g}_k(\mathbf{x})$，$\boldsymbol{\mathcal{G}}_{jk}=\partial^2 g_{jk}/\partial\mathbf{x}\partial\mathbf{x}^\top$）：
\[
  \mathbf{g}_k(\hat{\mathbf{x}})\approx\mathbf{g}_k+\mathbf{G}_k\delta\mathbf{x}+\tfrac12\sum_j\mathbf{1}_j\delta\mathbf{x}^\top\boldsymbol{\mathcal{G}}_{jk}\delta\mathbf{x},\quad
  \mathbf{G}_k(\hat{\mathbf{x}})\approx\mathbf{G}_k+\sum_j\mathbf{1}_j\delta\mathbf{x}^\top\boldsymbol{\mathcal{G}}_{jk}.
\]
设 $\delta\mathbf{x}=\mathbf{A}\mathbf{n}+\mathbf{b}(\mathbf{n})$（$\mathbf{A}$ 线性系数、$\mathbf{b}$ 二次）。代入展开、按 $\mathbf{n}$ 阶分组，得线性项 $\mathbf{L}\mathbf{n}$、两个二次项 $\mathbf{q}_1,\mathbf{q}_2$。要对任意 $\mathbf{n}$ 恒为零：考虑 $-\mathbf{n}$ 情形相减（二次项偶对称、相消），得 $2\mathbf{L}\mathbf{n}=\mathbf{0}$ $\Rightarrow$ $\mathbf{L}=\mathbf{0}$，定出 $\mathbf{A}=\mathbf{W}^{-1}\sum_k\mathbf{G}_k^\top\boldsymbol{\Sigma}_{v,k}^{-1}\mathbf{P}_k$（$\mathbf{P}_k$ 取第 $k$ 块噪声的投影，$\mathbf{W}=\sum_k\mathbf{G}_k^\top\boldsymbol{\Sigma}_{v,k}^{-1}\mathbf{G}_k$）。对二次项取期望：用恒等式 $\mathbf{A}\boldsymbol{\Sigma}\mathbf{A}^\top\equiv\mathbf{W}^{-1}$、$\mathbf{A}\boldsymbol{\Sigma}\mathbf{P}_k^\top\equiv\mathbf{W}^{-1}\mathbf{G}_k^\top$ 证 $\mathbb{E}[\mathbf{q}_2]=\mathbf{0}$，再由 $\mathbb{E}[\mathbf{n}^\top\mathbf{A}^\top\boldsymbol{\mathcal{G}}_{jk}\mathbf{A}\mathbf{n}]=\operatorname{tr}(\boldsymbol{\mathcal{G}}_{jk}\mathbf{A}\boldsymbol{\Sigma}\mathbf{A}^\top)=\operatorname{tr}(\boldsymbol{\mathcal{G}}_{jk}\mathbf{W}^{-1})$ 得
$\mathbb{E}[\mathbf{b}(\mathbf{n})]=-\tfrac12\mathbf{W}^{-1}\sum_k\mathbf{G}_k^\top\boldsymbol{\Sigma}_{v,k}^{-1}\sum_j\mathbf{1}_j\operatorname{tr}(\boldsymbol{\mathcal{G}}_{jk}\mathbf{W}^{-1})$。因 $\mathbb{E}[\mathbf{n}]=\mathbf{0}$，$\mathbb{E}[\delta\mathbf{x}]=\mathbb{E}[\mathbf{b}(\mathbf{n})]$ 即 \cref{eq:nlopt-mlbias}。\hfill$\square$
\end{derivation}

\begin{derivation}[（选读）GVI 损失导数与下降式（\cref{eq:gvi-derivs,eq:gvi-descent}）]
\label{der:nlopt-gvi-deriv}
\emph{对 $\boldsymbol{\mu}$ 的导数（\cref{eq:gvi-dmu}）}。损失 $V(q)=\mathbb{E}_q[\phi]+\tfrac12\ln\det\boldsymbol{\Sigma}^{-1}$ 的熵项不含 $\boldsymbol{\mu}$，故 $\partial V/\partial\boldsymbol{\mu}^\top=\partial\mathbb{E}_q[\phi]/\partial\boldsymbol{\mu}^\top$。把期望写成积分 $\mathbb{E}_q[\phi]=\int\phi(\mathbf{x})\,\mathcal{N}(\mathbf{x};\boldsymbol{\mu},\boldsymbol{\Sigma})\,\mathrm{d}\mathbf{x}$，对 $\boldsymbol{\mu}$ 求导只作用在高斯上。用 $\partial\mathcal{N}/\partial\boldsymbol{\mu}=\boldsymbol{\Sigma}^{-1}(\mathbf{x}-\boldsymbol{\mu})\mathcal{N}$（对高斯指数求导）得 $\partial\mathbb{E}_q[\phi]/\partial\boldsymbol{\mu}^\top=\boldsymbol{\Sigma}^{-1}\mathbb{E}_q[(\mathbf{x}-\boldsymbol{\mu})\phi]$，即 \cref{eq:gvi-dmu}。再求一次（用 $\partial[(\mathbf{x}-\boldsymbol{\mu})\mathcal{N}]/\partial\boldsymbol{\mu}=[(\mathbf{x}-\boldsymbol{\mu})(\mathbf{x}-\boldsymbol{\mu})^\top\boldsymbol{\Sigma}^{-1}-\mathbf{I}]\mathcal{N}$）得 \cref{eq:gvi-d2mu}。
\emph{对 $\boldsymbol{\Sigma}^{-1}$ 的导数（\cref{eq:gvi-dSinv}）}。熵项 $\tfrac12\ln\det\boldsymbol{\Sigma}^{-1}$ 对 $\boldsymbol{\Sigma}^{-1}$ 求导给 $\tfrac12\boldsymbol{\Sigma}$（用 $\partial\ln\det\mathbf{X}/\partial\mathbf{X}=\mathbf{X}^{-\top}$，\cref{ch:matderiv}）。期望项用 $\partial\mathcal{N}/\partial\boldsymbol{\Sigma}^{-1}=\tfrac12[\boldsymbol{\Sigma}-(\mathbf{x}-\boldsymbol{\mu})(\mathbf{x}-\boldsymbol{\mu})^\top]\mathcal{N}$ 得 $-\tfrac12\mathbb{E}_q[(\mathbf{x}-\boldsymbol{\mu})(\mathbf{x}-\boldsymbol{\mu})^\top\phi]+\tfrac12\boldsymbol{\Sigma}\mathbb{E}_q[\phi]$，合并即 \cref{eq:gvi-dSinv}。（此处按忽略 $\boldsymbol{\Sigma}^{-1}$ 对称性的“朴素”导数算；对称矩阵求导的修正 $\mathbf{A}\mapsto\mathbf{A}+\mathbf{A}^\top-\mathbf{A}\circ\mathbf{I}$ 不改“令导数为零”的零点，故求极值时可直接用朴素式，见 \cref{ch:matderiv}。）
\emph{Hessian--梯度关系（\cref{eq:gvi-hessgrad}）}。比较 \cref{eq:gvi-d2mu} 与 \cref{eq:gvi-dSinv}：后者含的 $\mathbb{E}_q[(\mathbf{x}-\boldsymbol{\mu})(\mathbf{x}-\boldsymbol{\mu})^\top\phi]$ 可由前者解出，代入即得 $\partial^2 V/\partial\boldsymbol{\mu}^\top\partial\boldsymbol{\mu}=\boldsymbol{\Sigma}^{-1}-2\boldsymbol{\Sigma}^{-1}(\partial V/\partial\boldsymbol{\Sigma}^{-1})\boldsymbol{\Sigma}^{-1}$。
\emph{下降式（\cref{eq:gvi-descent}）}。把 $V$ 对 $(\delta\boldsymbol{\mu},\delta\boldsymbol{\Sigma}^{-1})$ 展成“$\delta\boldsymbol{\mu}$ 二阶 $+$ $\delta\boldsymbol{\Sigma}^{-1}$ 一阶”：
\[
  V(q^{(i+1)})\approx V(q^{(i)})+\Big(\tfrac{\partial V}{\partial\boldsymbol{\mu}^\top}\Big)^\top\delta\boldsymbol{\mu}+\tfrac12\delta\boldsymbol{\mu}^\top\Big(\tfrac{\partial^2 V}{\partial\boldsymbol{\mu}^\top\partial\boldsymbol{\mu}}\Big)\delta\boldsymbol{\mu}+\operatorname{tr}\!\Big(\tfrac{\partial V}{\partial\boldsymbol{\Sigma}^{-1}}\delta\boldsymbol{\Sigma}^{-1}\Big).
\]
均值更新 \cref{eq:gvi-mean-update} 给 $(\partial V/\partial\boldsymbol{\mu}^\top)=-(\boldsymbol{\Sigma}^{-1})^{(i+1)}\delta\boldsymbol{\mu}$，代入前三项得 $-\tfrac12\delta\boldsymbol{\mu}^\top(\boldsymbol{\Sigma}^{-1})^{(i+1)}\delta\boldsymbol{\mu}$。协方差更新取 $\delta\boldsymbol{\Sigma}^{-1}=-2\boldsymbol{\Sigma}^{(i)}{}^{-1}(\partial V/\partial\boldsymbol{\Sigma}^{-1})\boldsymbol{\Sigma}^{(i)}{}^{-1}$（\cref{eq:gvi-ngd-components}），代入末项得 $-\tfrac12\operatorname{tr}(\boldsymbol{\Sigma}^{(i)}\delta\boldsymbol{\Sigma}^{-1}\boldsymbol{\Sigma}^{(i)}\delta\boldsymbol{\Sigma}^{-1})$，合并即 \cref{eq:gvi-descent}。\hfill$\square$
\end{derivation}

\paragraph{（选读）滑动窗口的工程动机与“删关键帧”的两难（十四讲 §10.1）。}\cite{gaoxiang2019slam14}
实际 SLAM 不能每时每刻跑整个 BA（关键帧只增不减，规模终将失控）。\emph{滑动窗口}只保留最近 $N$ 个关键帧（或按共视图取“时间近、空间散开”的帧，避免相机静止时窗口缩成一团退化），窗外的或\emph{固定}、或\emph{丢弃}。先看两端操作的难易对比\cite{gaoxiang2019slam14}：\textbf{往窗内\emph{新增}一个关键帧是平凡的}——它只是按正常 BA 把新位姿与它看到的路标连进来、照旧边缘化路标即可，稀疏的箭头结构（\cref{eq:ba-block}）丝毫不破；\textbf{唯独从窗内\emph{删}一个旧关键帧才会破坏稀疏}（详见下）。固定易懂；\emph{丢弃}却有理论两难：窗外变量并非与窗内无关（它们共视同一路标），简单删除会\emph{丢信息}，正确做法是\emph{边缘化}——把窗外信息\emph{压}进窗内、保留其影响。但边缘化\emph{旧关键帧}会带来一个新麻烦：被删关键帧 $\mathbf{x}_1$ 看到 $\mathbf{y}_1,\dots,\mathbf{y}_4$，对它做 Schur 消元会在这些\emph{路标之间}填入非零块（fill-in），\textbf{破坏右下角路标块 $\mathbf{C}$ 的对角结构}（\cref{eq:ba-block}），使 BA 不再能按箭头稀疏求解。对比 \cref{sec:nlopt-sparse}：边缘化\emph{路标}时 fill-in 落在相机块（$\mathbf{B}$ 本就不要求对角，无害）；边缘化\emph{关键帧}时 fill-in 落在路标块（致命）。工程补救（OKVIS\cite{leutenegger2015okvis} 等）：边缘化旧关键帧时\emph{连同}其独占路标一起边缘化（把路标信息转成剩余关键帧间的共视约束、维持 $\mathbf{C}$ 对角），其仍被新帧看到的路标则\emph{丢弃旧帧对它的那条观测}以保稀疏。概率上，边缘化即条件化 $p(\cdots)=p(\text{其余}\mid\mathbf{x}_1)\,p(\mathbf{x}_1)$ 后舍去 $p(\mathbf{x}_1)$；被删变量从此不再被使用，故滑动窗口适合 VO、不适合大尺度建图。下面给其矩阵实质。

\begin{derivation}[（选读）滑动窗口滤波：边缘化 $=$ Schur 补]
\label{der:nlopt-swf}
SWF（Sibley 2006\cite{sibley2010sliding}）在固定大小窗内迭代 GN 并滑动以求在线常数时间——介于离线批量与在线 EKF 之间。设窗内批量方程 $(\mathbf{H}^\top\mathbf{W}^{-1}\mathbf{H})\delta\mathbf{x}^\star=\mathbf{H}^\top\mathbf{W}^{-1}\mathbf{e}$ 为块三对角（运动链）。缩窗欲移除最旧态 $\mathbf{x}_0$：朴素删除会丢信息（初态知识不再参与）；正确做法是\emph{边缘化}——把窗内联合高斯（信息形式：$\mathbf{H}^\top\mathbf{W}^{-1}\mathbf{H}$ 是信息阵、$\mathbf{H}^\top\mathbf{W}^{-1}\mathbf{e}$ 是信息向量）对 $\mathbf{x}_0$ 边缘化。结果是对 $\mathbf{x}_1$ 块做 Schur 补：
\[
  \mathbf{A}_{11}\to\mathbf{A}_{11}-\mathbf{A}_{10}\mathbf{A}_{00}^{-1}\mathbf{A}_{10}^\top,\qquad
  \mathbf{b}_1\to\mathbf{b}_1-\mathbf{A}_{10}\mathbf{A}_{00}^{-1}\mathbf{b}_0,
\]
即把被移除态的信息\emph{压}进其邻接态（保留信息、不丢精度），$\mathbf{x}_0$ 出窗后在该窗迭代中保持常量。这与 \cref{eq:schur} 的 BA Schur 补、与变量消元（\cref{sec:nlopt-sparse}）是同一矩阵操作的不同应用。VIO 的滑窗边缘化（\cref{ch:vio}）即此。一般扩缩窗时，边缘化产生的“先验”块 $\bar{\mathbf{A}}_{kk}=\check{\boldsymbol{\Sigma}}_k^{-1}+\mathbf{F}_k^\top\boldsymbol{\Sigma}_{w,k+1}^{-1}\mathbf{F}_k+\mathbf{G}_k^\top\boldsymbol{\Sigma}_{v,k}^{-1}\mathbf{G}_k$ 与信息向量 $\bar{\mathbf{b}}_k$ 递归更新，其中等效先验信息 $\check{\boldsymbol{\Sigma}}_k^{-1}=\boldsymbol{\Sigma}_{w,k}^{-1}-\boldsymbol{\Sigma}_{w,k}^{-1}\mathbf{F}_{k-1}\bar{\mathbf{A}}_{k-1,k-1}^{-1}\mathbf{F}_{k-1}^\top\boldsymbol{\Sigma}_{w,k}^{-1}$（即上一步 Schur 补的对角块），它在窗内迭代时保持常量、无需重算——这正是 SWF 在线常数时间的关键\cite{barfoot2024state}。\hfill$\square$
\end{derivation}

\begin{derivation}[（选读）SWF 窗扩/窗缩的逐式递归块更新]
\label{der:nlopt-swf-recursion}
上一推导给出了“边缘化 $=$ Schur 补”的原理与递归先验块 $\bar{\mathbf{A}}_{kk},\check{\boldsymbol{\Sigma}}_k^{-1}$ 的最终形；这里把\emph{每一次滑动}拆成“\emph{窗扩}（右端加入新时刻）$\to$\emph{窗缩}（左端 Schur 边缘化最老态）”两步，并逐式写出三对角块的递归更新，使其可直接照搬成代码。沿用 \cref{sec:nlopt-appendix} 离散批量记号（$\mathbf{F}_k,\mathbf{G}_k$ 为运动/观测雅可比、$\mathbf{Q}_k',\mathbf{R}_k'$ 为相应噪声协方差、$\mathbf{e}_{v,k},\mathbf{e}_{y,k}$ 为运动/观测误差），并以窗长 $4$ 为例。

\emph{Step 1（初始窗）}。在初始窗 $\{\mathbf{x}_0,\dots,\mathbf{x}_3\}$ 上照常立批量 GN 方程并迭代收敛。因运动链只耦合相邻时刻，信息阵 $\mathbf{H}^\top\mathbf{W}^{-1}\mathbf{H}$ \emph{块三对角}：
\begin{equation}
\underbrace{\begin{bmatrix}
\bar{\mathbf{A}}_{00} & \mathbf{A}_{10}^\top & & \\
\mathbf{A}_{10} & \mathbf{A}_{11} & \mathbf{A}_{21}^\top & \\
 & \mathbf{A}_{21} & \mathbf{A}_{22} & \mathbf{A}_{32}^\top \\
 & & \mathbf{A}_{32} & \mathbf{A}_{33}
\end{bmatrix}}_{\mathbf{H}^\top\mathbf{W}^{-1}\mathbf{H}}
\begin{bmatrix}\delta\mathbf{x}_0^\star\\ \delta\mathbf{x}_1^\star\\ \delta\mathbf{x}_2^\star\\ \delta\mathbf{x}_3^\star\end{bmatrix}
=\begin{bmatrix}\bar{\mathbf{b}}_0\\ \mathbf{b}_1\\ \mathbf{b}_2\\ \mathbf{b}_3\end{bmatrix},
\label{eq:nlopt-swf-init}
\end{equation}
其对角块 $\bar{\mathbf{A}}_{00}=\check{\boldsymbol{\Sigma}}_0^{-1}+\mathbf{F}_0^\top\mathbf{Q}_1^{\prime-1}\mathbf{F}_0+\mathbf{G}_0^\top\mathbf{R}_0^{\prime-1}\mathbf{G}_0$ 含\emph{初始先验信息} $\check{\boldsymbol{\Sigma}}_0^{-1}$（给定）、$\mathbf{A}_{10}=-\mathbf{Q}_1^{\prime-1}\mathbf{F}_0$。

\emph{Step 2（窗扩：右端加入 $\mathbf{x}_4$）}。引入下一时刻 $\mathbf{x}_4$、并入最新输入与观测，只在右下角\emph{追加一行一列}（新块 $\mathbf{A}_{44},\mathbf{A}_{43}$ 与 $\mathbf{b}_4$ 全新，$\mathbf{A}_{33},\mathbf{b}_3$ 因新增运动项需重算），三对角结构\emph{不破}：
\begin{equation}
\begin{bmatrix}
\bar{\mathbf{A}}_{00} & \mathbf{A}_{10}^\top & & & \\
\mathbf{A}_{10} & \mathbf{A}_{11} & \mathbf{A}_{21}^\top & & \\
 & \mathbf{A}_{21} & \mathbf{A}_{22} & \mathbf{A}_{32}^\top & \\
 & & \mathbf{A}_{32} & \mathbf{A}_{33} & \mathbf{A}_{43}^\top \\
 & & & \mathbf{A}_{43} & \mathbf{A}_{44}
\end{bmatrix}
\begin{bmatrix}\delta\mathbf{x}_0^\star\\ \delta\mathbf{x}_1^\star\\ \delta\mathbf{x}_2^\star\\ \delta\mathbf{x}_3^\star\\ \delta\mathbf{x}_4^\star\end{bmatrix}
=\begin{bmatrix}\bar{\mathbf{b}}_0\\ \mathbf{b}_1\\ \mathbf{b}_2\\ \mathbf{b}_3\\ \mathbf{b}_4\end{bmatrix}.
\label{eq:nlopt-swf-expand}
\end{equation}
\textbf{这正是“新增帧平凡”的代数证据}（呼应 \cref{sec:nlopt-appendix} 正文）：加帧只是把块三对角阵向右下延拓一格，无 fill-in。此时窗长临时为 $5$。

\emph{Step 3（窗缩：左端边缘化 $\mathbf{x}_0$）}。为保持窗长，须移除最老态 $\mathbf{x}_0$。\emph{朴素删除}（直接划掉 $\mathbf{x}_0$ 的行列）会丢初态信息；\emph{正确做法}是把 \cref{eq:nlopt-swf-expand} 的信息形高斯对 $\mathbf{x}_0$ 边缘化，即对 $\mathbf{x}_1$ 块做 Schur 补（$\mathbf{x}_0$ 只与 $\mathbf{x}_1$ 相邻，故只污染 $(1,1)$ 块与 $\mathbf{b}_1$）：
\begin{equation}
\begin{bmatrix}
\underbrace{\mathbf{A}_{11}-\mathbf{A}_{10}\bar{\mathbf{A}}_{00}^{-1}\mathbf{A}_{10}^\top}_{\textstyle\bar{\mathbf{A}}_{11}} & \mathbf{A}_{21}^\top & & \\
\mathbf{A}_{21} & \mathbf{A}_{22} & \mathbf{A}_{32}^\top & \\
 & \mathbf{A}_{32} & \mathbf{A}_{33} & \mathbf{A}_{43}^\top \\
 & & \mathbf{A}_{43} & \mathbf{A}_{44}
\end{bmatrix}
\begin{bmatrix}\delta\mathbf{x}_1^\star\\ \delta\mathbf{x}_2^\star\\ \delta\mathbf{x}_3^\star\\ \delta\mathbf{x}_4^\star\end{bmatrix}
=\begin{bmatrix}\underbrace{\mathbf{b}_1-\mathbf{A}_{10}\bar{\mathbf{A}}_{00}^{-1}\bar{\mathbf{b}}_0}_{\textstyle\bar{\mathbf{b}}_1}\\ \mathbf{b}_2\\ \mathbf{b}_3\\ \mathbf{b}_4\end{bmatrix}.
\label{eq:nlopt-swf-contract}
\end{equation}
对比朴素删除版（缺 $-\mathbf{A}_{10}\bar{\mathbf{A}}_{00}^{-1}\mathbf{A}_{10}^\top$ 与 $-\mathbf{A}_{10}\bar{\mathbf{A}}_{00}^{-1}\bar{\mathbf{b}}_0$ 两项），\cref{eq:nlopt-swf-contract} 让 $\bar{\mathbf{A}}_{00},\mathbf{A}_{10},\bar{\mathbf{b}}_0$ \emph{仍参与解}——信息没丢。关键：\textbf{边缘化只在被删态\emph{唯一}的邻接块（这里 $(1,1)$）产生填充，三对角性保持}，故窗内仍可享稀疏求解（\cref{sec:nlopt-solve}）。这与 \cref{sec:nlopt-appendix} 正文“删\emph{关键帧}致路标块 fill-in（致命）”形成对照：在\emph{纯运动链}（无共享路标）里，被删态只接邻一个态，Schur 补不破三对角；BA 里被删关键帧接邻\emph{多个共视路标}，才会在路标块间撒下 fill-in。

\emph{Step 4（一般递归式）}。窗反复滑动后，第 $k$ 个时刻成为窗的左边界，其对角块 $\bar{\mathbf{A}}_{kk}$ 与信息向量 $\bar{\mathbf{b}}_k$ 由上一次边缘化\emph{递归}给出（把 \cref{eq:nlopt-swf-contract} 的 $0\to k{-}1$、$1\to k$ 一般化）：
\begin{subequations}\label{eq:nlopt-swf-rec}
\begin{align}
\bar{\mathbf{A}}_{kk} &= \check{\boldsymbol{\Sigma}}_k^{-1}+\mathbf{F}_k^\top\mathbf{Q}_{k+1}^{\prime-1}\mathbf{F}_k+\mathbf{G}_k^\top\mathbf{R}_k^{\prime-1}\mathbf{G}_k, \label{eq:nlopt-swf-Akk}\\
\bar{\mathbf{b}}_k &= \mathbf{c}_k+\check{\boldsymbol{\Sigma}}_k^{-1}\mathbf{e}_{v,k}-\mathbf{F}_k^\top\mathbf{Q}_{k+1}^{\prime-1}\mathbf{e}_{v,k+1}+\mathbf{G}_k^\top\mathbf{R}_k^{\prime-1}\mathbf{e}_{y,k}, \label{eq:nlopt-swf-bk}\\
\check{\boldsymbol{\Sigma}}_k^{-1} &= \mathbf{Q}_k^{\prime-1}-\mathbf{Q}_k^{\prime-1}\mathbf{F}_{k-1}\,\bar{\mathbf{A}}_{k-1,k-1}^{-1}\,\mathbf{F}_{k-1}^\top\mathbf{Q}_k^{\prime-1}, \label{eq:nlopt-swf-Pk}\\
\mathbf{c}_k &= \mathbf{Q}_k^{\prime-1}\mathbf{F}_{k-1}\,\bar{\mathbf{A}}_{k-1,k-1}^{-1}\Big(\mathbf{c}_{k-1}+\check{\boldsymbol{\Sigma}}_{k-1}^{-1}\mathbf{e}_{v,k-1}+\mathbf{G}_{k-1}^\top\mathbf{R}_{k-1}^{\prime-1}\mathbf{e}_{y,k-1}\Big), \label{eq:nlopt-swf-ck}
\end{align}
\end{subequations}
以 $\mathbf{c}_0=\mathbf{0}$、$\check{\boldsymbol{\Sigma}}_0^{-1}=$ 给定初始信息阵起步。其中 \cref{eq:nlopt-swf-Pk} 即上一步 Schur 补\emph{对角块}（边缘化在 $(k,k)$ 块留下的等效\emph{先验信息}），\cref{eq:nlopt-swf-ck} 是其对应的\emph{先验信息向量}修正——二者正是被边缘化的左侧历史\emph{压缩成}的一个“等效先验因子”，把无穷长历史浓缩为窗左界上的\emph{一个}高斯先验（与 EKF 把历史压进单一均值/协方差同构，但 SWF 仍\emph{迭代}窗内多个时刻、只锁定\emph{出窗}态的线性化点）。

\emph{Step 5（关键性质：常量先验）}。最要紧的一点：$\bar{\mathbf{A}}_{kk}$ 与 $\bar{\mathbf{b}}_k$ 中来自边缘化的 $\check{\boldsymbol{\Sigma}}_k^{-1}$ 与 $\mathbf{c}_k$ \emph{只依赖已出窗、已被锁定线性化点的态}，故对当前窗内变量而言是\emph{常量}：窗内 GN \emph{迭代重线性化}时它们\textbf{无需重算}，每滑一步只新做一次 Schur 补。正是这一点把每步代价钉成\emph{与轨迹长度无关的常数时间}，使 SWF 成为批量平滑（\cref{thm:nlopt-map}）的在线、定窗工程化形态——介于离线全批量与单时刻 IEKF（$=$ 窗长 $1$、且只迭代校正步）之间，亦即 STEAM/固定滞后平滑（fixed-lag smoothing，\cref{sec:est-steam}）的核心机理。\hfill$\square$
\end{derivation}

\begin{algo}[滑动窗滤波（SWF）：窗扩--窗缩滑动循环]\label{alg:nlopt-swf}
给窗长 $W$、初始先验 $\{\check{\mathbf{x}}_0,\check{\boldsymbol{\Sigma}}_0^{-1}\}$。
(1) \textbf{初始窗}：在 $\{\mathbf{x}_0,\dots,\mathbf{x}_{W-1}\}$ 上立块三对角批量 GN \cref{eq:nlopt-swf-init}，迭代收敛；
(2) \textbf{窗扩}：纳入新时刻 $\mathbf{x}_{k+W}$，按 \cref{eq:nlopt-swf-expand} 在右下追加 $\mathbf{A}_{k+W,k+W},\mathbf{A}_{k+W,k+W-1},\mathbf{b}_{k+W}$（临时窗长 $W{+}1$）；
(3) \textbf{窗缩}：对最老态 $\mathbf{x}_k$ 做 Schur 边缘化 \cref{eq:nlopt-swf-contract}，把 $\bar{\mathbf{A}}_{kk}^{-1}$ 压进邻接块、得新左界的常量先验 $\check{\boldsymbol{\Sigma}}_{k+1}^{-1},\mathbf{c}_{k+1}$（\cref{eq:nlopt-swf-Pk,eq:nlopt-swf-ck}）；锁定出窗态 $\mathbf{x}_k$ 的线性化点并输出其估计；
(4) \textbf{迭代新窗}：在 $\{\mathbf{x}_{k+1},\dots,\mathbf{x}_{k+W}\}$ 上迭代 GN 到收敛（先验块 \cref{eq:nlopt-swf-rec} 视为常量、不重算）；
(5) $k\leftarrow k+1$，回到 (2)，直至数据耗尽。
\textbf{要点}：输出取窗\emph{最左}（出窗）态可纳入未来信息但引入延迟；要低延迟可取窗内较新态。窗长 $1$ 且只迭代校正步即退化为 IEKF（\cref{thm:nlopt-iekf-map}）。
\end{algo}

\begin{derivation}[（选读）连续时间批量估计（高斯过程回归）]
把先验从离散运动链换成连续时间随机微分方程 $\dot{\mathbf{x}}(t)=f(\mathbf{x},\mathbf{v},\mathbf{w},t)$（白噪声功率谱密度 $\mathbf{w}\sim\mathcal{GP}(\mathbf{0},\mathbf{Q}_C\,\delta(t-t'))$），围绕工作轨迹 $\mathbf{x}_{\mathrm{op}}(t)$ 线性化得 LTV 形式 $\dot{\mathbf{x}}\approx\mathbf{A}(t)\mathbf{x}+\boldsymbol{\nu}(t)+\mathbf{L}(t)\mathbf{w}$（连续系统矩阵记 $\mathbf{A}(t)$、白噪声 PSD 记 $\mathbf{Q}_C$，字母与 \cref{sec:est-steam} 的 STEAM、\cref{thm:matderiv-vanloan} 的 Van Loan 离散化一致；离散累积噪声则记 $\mathbf{Q}'$ 以示区别），则 $\mathbf{x}(t)$ 是高斯过程，均值/协方差由转移函数 $\boldsymbol{\Phi}(t,s)$ 积分给出。在量测时刻堆叠成 lifted 先验 $\mathbf{x}\sim\mathcal{N}(\mathbf{F}\boldsymbol{\nu},\mathbf{F}\mathbf{Q}'\mathbf{F}^\top)$（此处 $\mathbf{F}$ 为\emph{堆叠}转移阵、$\mathbf{Q}'$ 为离散累积过程噪声，对应 \cref{sec:est-steam} 的 $\mathbf{F}^{-1}$ 与 $\mathbf{Q}(\Delta t)$），与离散批量\emph{同形}地得 GN 更新 $(\mathbf{F}^{-\top}\mathbf{Q}^{\prime-1}\mathbf{F}^{-1}+\mathbf{G}^\top\mathbf{R}^{\prime-1}\mathbf{G})\delta\mathbf{x}^*=\cdots$，仍块三对角。优点：可在\emph{任意}时刻 $O(1)$ 高斯插值查询轨迹 $\mathbf{x}_{\mathrm{op}}(s)=\check{\mathbf{x}}(s)+\check{\boldsymbol{\Sigma}}(s)\check{\boldsymbol{\Sigma}}^{-1}(\mathbf{x}_{\mathrm{op}}-\check{\mathbf{x}})$（$\check{\boldsymbol{\Sigma}}$ 为先验/查询协方差），适合异步/高频传感器（对接 STEAM、\cref{ch:vio}）。转移函数 $\boldsymbol{\Phi}(t,s)=\boldsymbol{\Upsilon}(t)\boldsymbol{\Upsilon}(s)^{-1}$ 可由归一化基本阵 $\dot{\boldsymbol{\Upsilon}}=\mathbf{A}(t)\boldsymbol{\Upsilon},\ \boldsymbol{\Upsilon}(0)=\mathbf{I}$ 数值积分得到；每次迭代代价仍随轨迹长度线性\cite{barfoot2024state}。\hfill$\square$
\end{derivation}
